% MERGE generado automáticamente a partir de libroVersion1CorrecionAndres.tex y libroVersion2CorrecionJA.tex
% Base estructural: versión JA. Contenido matemático/métrico ampliado: versión Andrés.
% ============================================================================
% Libro de Tesis / Trabajo Final de Grado - FPUNA (ACTUALIZADO A IMPLEMENTACIÓN)
% Autores: Andrés Román, Juan González
% Tutor: Dr. Cristian Cappo
% ============================================================================
\documentclass[12pt,oneside]{report}

% --- Idioma y codificación ---
\usepackage[spanish]{babel}
\usepackage[utf8]{inputenc}
\usepackage[T1]{fontenc}
\usepackage{lmodern}

% --- Página ---
\usepackage{geometry}
\geometry{a4paper, margin=2.5cm}

% --- Utilidades ---
\usepackage{graphicx}
\usepackage{booktabs}
\usepackage{longtable}
\usepackage{tabularx}
\usepackage{array}
\usepackage{ragged2e}
\usepackage{hyperref}
\usepackage{amsmath}
\usepackage{amssymb}
\usepackage{url}
\usepackage{xurl}   % mejor corte de URLs largas
\usepackage{float}  % para [H]
\usepackage{placeins} % \FloatBarrier
\usepackage{enumitem}
\usepackage{textcomp}
\usepackage{makecell}
\usepackage{adjustbox}
\usepackage{multirow}
\usepackage{microtype}
\usepackage[justification=raggedright,singlelinecheck=false]{caption}

% --- Figuras (hechas en LaTeX) ---
\usepackage{tikz}
\usepackage{pgfplots}
\pgfplotsset{compat=1.18}
\usetikzlibrary{arrows.meta,positioning,calc,shapes.geometric}

% --- Hipervínculos ---
\hypersetup{
  colorlinks=true,
  linkcolor=black,
  urlcolor=blue,
  citecolor=black
}

% --- Ruta opcional de imágenes (si ponés el logo en ./images/) ---
% \graphicspath{{images/}}

% --- Columnas útiles para tabularx ---
\newcolumntype{L}[1]{>{\RaggedRight\arraybackslash}p{#1}}
\newcolumntype{Y}{>{\RaggedRight\arraybackslash}X}
\newcolumntype{C}[1]{>{\Centering\arraybackslash}p{#1}}

% --- Macros útiles ---
\newcommand{\R}{\mathbb{R}}
\newcommand{\sigmoid}{\sigma}
\newcommand{\softmax}{\mathrm{softmax}}
\newcommand{\ind}{\mathbb{I}} % función indicadora

\setlength{\LTpre}{0pt}
\setlength{\LTpost}{6pt}
\setlength{\emergencystretch}{3em}
\captionsetup[table]{justification=raggedright,singlelinecheck=false}
\captionsetup[longtable]{justification=raggedright,singlelinecheck=false}

\begin{document}

% ============================================================================
% CARÁTULA
% ============================================================================
\begin{titlepage}
\thispagestyle{empty}
\begin{center}
{\Large Universidad Nacional de Asunción}\\[0.25cm]
{\Large Facultad Politécnica}\\[0.25cm]
{\Large Ingeniería en Informática}\\[1.2cm]

\IfFileExists{logoFpuna.png}{\includegraphics[width=4cm]{logoFpuna.png}}{\vspace{4cm}}\par\vspace{1.5cm}

\rule{0.9\textwidth}{0.4pt}\\[0.9cm]
{\Large Trabajo Final de Grado}\\[0.6cm]
{\LARGE\bfseries Detección de anomalías y clasificación de ataques en tráfico HTTP: comparación entre enfoques one-class, multiclase y multietiqueta}\\[0.9cm]
\rule{0.9\textwidth}{0.4pt}\\[1.2cm]

\begin{minipage}[t]{0.45\textwidth}
\textit{Autores:}\\[0.3cm]
Andrés Román\\
Juan González
\end{minipage}
\hfill
\begin{minipage}[t]{0.45\textwidth}
\raggedleft
\textit{Tutor:}\\[0.3cm]
Dr. Cristian Cappo
\end{minipage}

\vfill
SAN LORENZO - PARAGUAY\\
ENERO - 2026
\end{center}
\end{titlepage}

% ============================================================================
% PRELIMINARES
% ============================================================================
\pagenumbering{roman}

\chapter*{Resumen}
Las aplicaciones web están expuestas de manera continua a ataques como inyección SQL, \textit{Cross-Site Scripting} y otras variantes de manipulación de peticiones HTTP. Detectar estos ataques de forma efectiva es un problema complejo: el tráfico legítimo presenta alta variabilidad, existe un fuerte desbalance entre peticiones normales y maliciosas, y los atacantes generan constantemente variantes o patrones de ataque para los cuales todavía no existen reglas, firmas o ejemplos suficientes en los sistemas de protección tradicionales. Los \textit{Web Application Firewalls} (WAF) basados únicamente en reglas o firmas requieren mantenimiento continuo y tienden a degradarse frente a estas variaciones, lo que limita su eficacia en escenarios cambiantes.

Ante estas limitaciones, el aprendizaje automático ofrece una alternativa de análisis que permite capturar patrones en la estructura del tráfico HTTP sin depender exclusivamente de un catálogo fijo de ataques conocidos. En este trabajo se presenta un flujo reproducible de aprendizaje automático orientado a la detección de anomalías y a la clasificación de ataques en tráfico HTTP. La comparación experimental se organiza sobre tres formulaciones complementarias del problema: detección \textit{one-class}, clasificación multiclase y clasificación multietiqueta. El alcance es experimental y comparativo. No se implementa ni se despliega un \textit{Web Application Firewall} completo; en su lugar, se examina la factibilidad teórica de utilizar los modelos evaluados en un contexto tipo WAF a partir de métricas de efectividad y de costo operativo.

La representación de entrada se construye sobre una base común de peticiones HTTP (\textit{requests}), definida mediante variables interpretables extraídas del método, la URI, el cuerpo y, cuando está disponible, un campo derivado de encabezados. Esa base se materializa en un conjunto fijo de veinticinco características (\textit{features}) que describen longitud, estructura, composición de caracteres, codificación, señales léxicas y atributos básicos del protocolo.

La comparación experimental se distribuye sobre tres fuentes de datos con distinta granularidad de etiquetado: CSIC 2010, útil principalmente para detección binaria; CSIC/TorpEda 2012, apto para detección binaria y clasificación multiclase; y SR-BH 2020, un dataset de tráfico real de honeypot con anotación CAPEC multietiqueta. Para preservar comparabilidad entre escenarios, la implementación normaliza los campos de entrada y deriva las variables objetivo \texttt{label\_binary}, \texttt{label\_multiclass} y \texttt{label\_multilabel} según la información disponible en cada dataset.

Los modelos seleccionados responden a criterios de escalabilidad, trazabilidad e interpretabilidad. Para detección \textit{one-class} se utiliza un \texttt{Pipeline} compuesto por \texttt{StandardScaler}, \texttt{Nystroem} y \texttt{SGDOneClassSVM}, entrenado únicamente con tráfico normal. Para las tareas supervisadas se emplea regresión logística lineal, mientras que para clasificación multietiqueta se adopta la descomposición \textit{One-vs-Rest}. El protocolo experimental aísla el conjunto de prueba externo, restringe el ajuste de hiperparámetros al bloque de entrenamiento y contempla tanto métricas de clasificación como métricas operativas, entre ellas latencia, \textit{throughput}, consumo de CPU y memoria, tamaño de modelo y tiempo de entrenamiento.

La metodología incorpora además una comparación directa entre corridas con el conjunto completo de variables y corridas en las que se eliminan únicamente las variables basadas en tokens sospechosos. La comparación permite distinguir con mayor precisión cuánto del rendimiento observado proviene de señales estructurales del tráfico HTTP y cuánto depende de indicadores léxicos más cercanos a reglas o firmas.

\chapter*{Abstract}
Web applications are continuously exposed to attacks such as SQL Injection, \textit{Cross-Site Scripting}, path traversal, and other forms of HTTP request manipulation. Detecting these attacks effectively is a complex problem because legitimate traffic is highly variable, malicious requests are usually a minority with respect to normal traffic, and attackers constantly introduce new variants for which traditional protection systems may not yet have rules, signatures, or sufficient labeled examples. Web Application Firewalls (WAFs) based only on rules or signatures require continuous maintenance and may lose effectiveness when payloads change or when attack patterns evolve.

In response to these limitations, machine learning provides an alternative analysis strategy by learning patterns from the structure and content of HTTP traffic without depending exclusively on a fixed catalog of known attacks. This work presents a reproducible machine learning workflow for anomaly detection and attack classification in HTTP traffic. The experimental comparison is organized around three complementary problem formulations: \textit{one-class} detection, multiclass classification, and multilabel classification. The scope is experimental and comparative. A full Web Application Firewall is neither implemented nor deployed; instead, the theoretical feasibility of using the evaluated models in a WAF-like setting is examined through effectiveness metrics and operational-cost metrics.

The input representation is built from a common description of HTTP requests based on interpretable variables extracted from the method, URI, body, and, when available, a header-derived field. This design is materialized through a fixed set of twenty-five features describing length, structure, character composition, encoding, lexical cues, and basic protocol attributes.

The experimental comparison uses three data sources with different labeling granularity: CSIC 2010, mainly useful for binary detection; CSIC/TorpEda 2012, suitable for binary detection and multiclass classification; and SR-BH 2020, a real honeypot traffic dataset with CAPEC multilabel annotations. To preserve comparability across scenarios, the implementation normalizes the input fields and derives the target variables \texttt{label\_binary}, \texttt{label\_multiclass}, and \texttt{label\_multilabel} according to the information available in each dataset.

The selected models follow the criteria of scalability, traceability, and interpretability. For \textit{one-class} detection, the implementation uses a \texttt{Pipeline} composed of \texttt{StandardScaler}, \texttt{Nystroem}, and \texttt{SGDOneClassSVM}, trained only with normal traffic. For supervised tasks, linear logistic regression is used, while multilabel classification is handled through the \textit{One-vs-Rest} decomposition. The experimental protocol isolates an external test set, restricts hyperparameter tuning to the training block, and considers both classification metrics and operational metrics such as latency, throughput, CPU and memory usage, model size, and training time.

The methodology also includes a direct comparison between runs using the complete feature set and runs in which only suspicious-token-based variables are removed. This comparison makes it possible to distinguish how much of the observed performance comes from structural HTTP signals and how much depends on lexical indicators that are closer to rules or signatures.

\chapter*{Agradecimientos}
Agradecemos primeramente a Dios por la oportunidad, la fortaleza y la perseverancia necesarias para completar este Trabajo Final de Grado. También expresamos nuestro agradecimiento a nuestras familias, por su acompañamiento constante, paciencia y apoyo durante el desarrollo de este proyecto.

Agradecemos especialmente al Dr. Cristian Cappo, tutor de este trabajo, por su orientación académica, sus observaciones y el acompañamiento brindado durante el proceso de investigación. Sus recomendaciones permitieron mejorar la delimitación del problema, el diseño experimental y la presentación de los resultados.

Extendemos nuestro reconocimiento a la Facultad Politécnica de la Universidad Nacional de Asunción y a los docentes de la carrera de Ingeniería en Informática, por la formación recibida a lo largo de la carrera. Finalmente, agradecemos a todas las personas que, directa o indirectamente, colaboraron con comentarios, revisión, apoyo técnico o motivación para la culminación de este trabajo.

\tableofcontents
\listoffigures
\clearpage
\pagenumbering{arabic}

% ============================================================================
% CAPÍTULO 1
% ============================================================================
\chapter{Introducción}

\section{Contexto y motivación}
Las aplicaciones web son hoy uno de los principales vectores de ataque en sistemas de información. Ataques como inyección SQL, \textit{Cross-Site Scripting} (XSS) o manipulación de parámetros se ejecutan mediante peticiones HTTP (\textit{requests}) especialmente construidas para explotar vulnerabilidades en el servidor o en la aplicación. Este tipo de tráfico malicioso es difícil de separar del tráfico legítimo por varias razones simultáneas: el tráfico normal presenta alta variabilidad según el contexto de uso, existe un fuerte desbalance numérico entre peticiones normales y maliciosas, y los atacantes generan constantemente variantes o patrones de ataque para los cuales todavía no existen reglas, firmas o ejemplos suficientes en los sistemas de protección existentes.

Frente a este problema, una solución ampliamente adoptada es el \textit{Web Application Firewall} (WAF), un componente que inspecciona el tráfico HTTP entrante y aplica políticas para mitigar ataques conocidos \cite{owasp_waf}. La variante más extendida de WAF opera mediante firmas o reglas: listas de patrones predefinidos que, al detectarse en una petición, desencadenan un bloqueo o una alerta. ModSecurity es un ejemplo representativo de este enfoque, ampliamente utilizado en producción como motor de reglas configurable para monitoreo y bloqueo de tráfico \cite{modsecurity_site}. Sin embargo, los WAF basados exclusivamente en firmas o reglas presentan limitaciones estructurales: requieren mantenimiento continuo a medida que aparecen nuevas técnicas de ataque, y tienden a degradarse frente a variaciones de la carga útil (\textit{payload}) o frente a ataques para los cuales todavía no existe una firma actualizada.

Como alternativa, los enfoques basados en anomalías no dependen de un catálogo fijo de ataques conocidos. En lugar de buscar patrones explícitos, aprenden una representación del tráfico ``normal'' y marcan como sospechosas las peticiones que se desvían de esa representación. Esto resulta útil cuando la cobertura de ataques es incompleta o cambiante, ya que el modelo puede señalar actividad inusual aun sin haber visto antes un ataque de ese tipo. Trabajos clásicos como el de Kruegel y Vigna muestran que atributos estadísticos y de estructura de las peticiones HTTP permiten puntuar anomalías con buenos resultados \cite{kruegel_vigna_2003}.

Sin embargo, la detección de anomalías por sí sola tiene un límite: señala que algo es inusual, pero no identifica qué tipo de ataque es. Para respuesta efectiva e investigación de incidentes, resulta valioso poder también clasificar el tipo o patrón de ataque. Esto introduce las formulaciones supervisadas, donde se dispone de etiquetas que identifican clases o categorías de ataque. El aprendizaje automático permite, en este contexto, entrenar modelos que asignan una o varias etiquetas a cada petición, aprovechando la información semántica disponible en datasets con anotación por tipo de ataque.

El planteamiento adoptado en este trabajo toma como base el enfoque por anomalías y lo extiende con un dataset de tráfico HTTP real con etiquetas de ataque estructuradas según la taxonomía CAPEC, habilitando comparaciones entre la formulación \textit{one-class}, la clasificación supervisada multiclase y la clasificación multietiqueta.

\section{Conceptos fundamentales}

A continuación se presentan las definiciones de los conceptos centrales utilizados en este trabajo, orientadas a lectores sin experiencia previa en seguridad web o aprendizaje automático.

\begin{description}[leftmargin=0pt, labelindent=0pt]

  \item[\textbf{Petición HTTP (\textit{request}).}]
  Una petición HTTP es el mensaje que un cliente (por ejemplo, un navegador web) envía a un servidor para solicitar un recurso o ejecutar una acción. Contiene al menos un método (como \texttt{GET} o \texttt{POST}), una URI que identifica el recurso, la versión del protocolo y, opcionalmente, encabezados y un cuerpo. En este trabajo, el alcance de ``petición HTTP'' se limita a los campos disponibles y normalizados en los datasets: principalmente método, URI, cuerpo y, cuando está disponible, información derivada de encabezados.

  \item[\textbf{WAF (\textit{Web Application Firewall}).}]
  Un WAF es un componente de seguridad que se interpone entre el cliente y la aplicación web para inspeccionar el tráfico HTTP y aplicar políticas de filtrado. Su objetivo es detectar y bloquear peticiones maliciosas antes de que lleguen a la aplicación.

  \item[\textbf{Carga útil (\textit{payload}).}]
  En el contexto de este trabajo, la carga útil o \textit{payload} es la parte de la petición HTTP donde pueden aparecer los datos enviados por el cliente: parámetros, cuerpo de la solicitud o cadenas utilizadas por un atacante para explotar una vulnerabilidad. Es en la carga útil donde con frecuencia se insertan los patrones maliciosos.

  \item[\textbf{Firma o regla.}]
  Una firma es un patrón predefinido que describe cómo se manifiesta un tipo de ataque conocido. Un WAF basado en firmas compara cada petición contra una lista de esos patrones y bloquea las que coinciden. Su ventaja es la alta precisión en ataques conocidos; su desventaja es que no detecta variantes para las cuales todavía no existe una firma.

  \item[\textbf{Anomalía.}]
  Una anomalía es una petición que se desvía estadística o estructuralmente del comportamiento habitual del tráfico normal. No requiere coincidir con un patrón de ataque conocido: basta con que difiera suficientemente de lo que el modelo aprendió como ``normal''.

  \item[\textbf{Clasificación binaria.}]
  Tarea de aprendizaje automático en la que el modelo asigna a cada petición una de dos etiquetas posibles: normal o ataque. Es la formulación más simple y directa para detección de intrusiones.

  \item[\textbf{Clasificación multiclase.}]
  Tarea en la que el modelo asigna a cada petición exactamente una clase entre varias posibles, por ejemplo: normal, inyección SQL, XSS, traversal de ruta, etc. Permite tipificar el ataque, pero supone que cada petición pertenece a una única categoría.

  \item[\textbf{CAPEC (\textit{Common Attack Pattern Enumeration and Classification}).}]
  CAPEC es una taxonomía de patrones de ataque que organiza y describe formas comunes en que un atacante puede explotar debilidades de un sistema. Cada patrón CAPEC representa una técnica o estrategia de ataque con un identificador propio, por ejemplo CAPEC-66 para inyección SQL o CAPEC-63 para XSS. En este trabajo, CAPEC se utiliza como referencia para interpretar y agrupar tipos de ataque presentes en los datasets, especialmente cuando una petición puede estar asociada a una o más categorías de ataque.

  \item[\textbf{Clasificación multietiqueta.}]
  Tarea en la que el modelo puede asignar simultáneamente más de una etiqueta a una misma petición. Es útil cuando una petición combina técnicas de varios tipos de ataque, como puede ocurrir en tráfico real anotado con la taxonomía CAPEC.

  \item[\textbf{Modelo de \textit{Machine Learning} (aprendizaje automático).}]
  A partir de un conjunto de ejemplos con características de entrada y etiquetas conocidas, un modelo de aprendizaje automático aprende una función de decisión para asignar una clase o conjunto de etiquetas a nuevos ejemplos no vistos durante el entrenamiento.

\end{description}

\section{Formulaciones del problema}

La detección de ataques en tráfico HTTP puede formularse de maneras distintas según la información disponible y el objetivo operativo. En este trabajo se contemplan tres formulaciones complementarias:

\begin{itemize}[leftmargin=*]
  \item \textbf{Detección binaria u \textit{one-class}.} El objetivo es distinguir tráfico normal de tráfico anómalo o malicioso, sin identificar el tipo de ataque. La variante \textit{one-class} es especialmente relevante cuando no se dispone de ejemplos etiquetados de ataque: el modelo se entrena solo con tráfico normal y aprende a detectar desviaciones. Es útil frente a ataques novedosos o cuando el catálogo de ataques conocidos es incompleto.

  \item \textbf{Clasificación multiclase.} Cuando se dispone de etiquetas que identifican el tipo de ataque, es posible entrenar un modelo que asigne una única clase o tipo principal a cada petición. Esto permite no solo detectar, sino también tipificar el ataque, lo que facilita la respuesta y el análisis forense. La limitación es que supone que cada petición pertenece a un único tipo de ataque.

  \item \textbf{Clasificación multietiqueta.} En tráfico real, una misma petición puede participar de más de un patrón de ataque simultáneamente. La formulación multietiqueta permite que el modelo asigne varias etiquetas a una misma petición, conservando la riqueza semántica de la anotación original. En este trabajo, esta formulación se habilita sobre el dataset SR-BH~2020, que incorpora etiquetas según la taxonomía CAPEC.
\end{itemize}

Las tres formulaciones no son excluyentes: este trabajo las evalúa comparativamente sobre los mismos datasets y con la misma representación de entrada, para estudiar qué formulación resulta más adecuada según el tipo de datos disponible y el objetivo operativo.

\section{Planteamiento del problema}
En entornos reales existe: (i) alta variabilidad del tráfico legítimo, (ii) fuerte desbalance entre normales y ataques,
y (iii) ataques emergentes, es decir, variantes o patrones de ataque para los cuales todavía no existen reglas, firmas o ejemplos suficientes en el sistema. Un sistema útil debe:
\begin{itemize}[leftmargin=*]
  \item detectar anomalías cuando hay pocos datos rotulados (o cambian los ataques), y
  \item cuando hay etiquetas, clasificar el tipo o patrón de ataque para respuesta e investigación.
\end{itemize}

\section{Objetivos}
\subsection{Objetivo general}

El objetivo general es diseñar, implementar y evaluar un \textit{pipeline} reproducible de aprendizaje automático para
(i) detección de anomalías/ataques en tráfico HTTP y (ii) clasificación de ataques por tipo (multiclase y, cuando aplique, multietiqueta),
comparando un dataset de tráfico \textbf{real} recolectado en honeypot (SR-BH 2020, Harvard Dataverse) con datasets \textbf{sintéticos/benchmark}
(CSIC 2010 y CSIC/TorpEda 2012) para estudiar \textit{domain shift}.
Además, se analizará la \textbf{viabilidad de implementación en tiempo real} del pipeline (extracción de features + inferencia) en un entorno tipo WAF,
midiendo eficiencia (latencia P50/P95/P99, throughput, CPU/RAM, tamaño del modelo, tiempo de entrenamiento) y funcionamiento (TPR/FPR y trade-offs FP/FN),
y se compararán los resultados con trabajos relacionados usando un conjunto común de métricas \cite{srbh_paper,ocswaf_zenodo,riverol_2024,fang_2024,zhou_2024,saito_pr_auc,chicco_mcc}.

En términos experimentales, la comparación se construye sobre una misma representación de peticiones HTTP aplicada a tres datasets y a tres formulaciones del problema. El contraste no se limita al rendimiento predictivo: incorpora también costo computacional, estabilidad y trazabilidad de las decisiones, de modo que la discusión final distinga con claridad entre el aporte de cada modelo y las exigencias propias de cada escenario de datos.


\subsection{Objetivos específicos}
\begin{itemize}[leftmargin=*]
  \item Implementar un pipeline de extracción de \textbf{25 características} desde método, URI, cuerpo HTTP y, cuando esté disponible, un campo derivado de \textit{headers} (\texttt{req\_content\_length}), alineado con literatura de detección en requests.
  \item Unificar y evaluar tres fuentes de datos: SR-BH 2020 (tráfico real de honeypot con etiquetas CAPEC) \cite{srbh_dataverse,srbh_paper}, y dos datasets sintéticos/benchmark (CSIC 2010 y CSIC/TorpEda 2012) \cite{csic2010_impact,csic2010_doc,torpeda_2012}.
  \item Entrenar y comparar \textbf{tres configuraciones/modelos} sobre los \textbf{tres datasets}, \textbf{adaptando la tarea a la disponibilidad de etiquetas} de cada dataset: (i) un detector one-class escalable para detección binaria (normal/anómalo) entrenado solo con normales, implementado como \texttt{StandardScaler + Nystroem + SGDOneClassSVM} \cite{sklearn_sgd_ocsvm,sklearn_kernel_approx}; (ii) regresión logística supervisada lineal como baseline eficiente \cite{islr_logreg,sklearn_logreg} (multiclase cuando el dataset provea tipo de ataque o etiqueta primaria; en CSIC~2010 se restringe a binario); (iii) One-vs-Rest (OvR) como baseline multi-salida \cite{tsoumakas_multilabel_overview,rifkin_ova} (multietiqueta CAPEC en SR-BH; en TorpEda se usa en modo multiclase o binario cuando corresponda; en CSIC~2010 se restringe a binario).
  \item Definir una metodología experimental \textbf{sin contaminación entre entrenamiento, validación y prueba}: en supervisado, usar particionado externo \textbf{80/20} del total del dataset; y en el detector one-class, entrenar \textbf{solo con el 80\% del tráfico normal} y evaluar sobre un conjunto mixto compuesto por el \textbf{20\% de normales en holdout} y el \textbf{total de anomalías/ataques disponibles}, realizando la selección de hiperparámetros \textbf{solo sobre el conjunto habilitado para entrenamiento} y con presupuesto computacional fijo.
  \item Reportar \textbf{métricas de efectividad} alineadas con trabajos relacionados: Accuracy, Precision/Recall/F1 (micro/macro/weighted), TPR/FPR, Specificity, ROC-AUC y PR-AUC \cite{saito_pr_auc}, además de MCC como métrica robusta al desbalance \cite{chicco_mcc}. En multietiqueta: Hamming Loss, Jaccard Similarity y Exact Match Ratio \cite{srbh_paper,tsoumakas_multilabel_overview}.
  \item Medir \textbf{métricas de eficiencia/tiempo real}: latencia P50/P95/P99 del pipeline (feature extraction + inferencia), throughput (req/s), CPU/RAM, tamaño del modelo y tiempo de entrenamiento; y estabilidad bajo carga \cite{ocswaf_zenodo}.
  \item Analizar el aporte de las \textit{features} mediante una estrategia eficiente: coeficientes estandarizados en los modelos lineales, ablaciones por grupos y comparación \textbf{con y sin tokens sospechosos}, priorizando interpretabilidad con costo computacional controlado.
  \item Comparar resultados con los reportados en literatura sobre CSIC/TorpEda/SR-BH, alineando tareas y métricas, y documentando amenazas a validez cuando difieran \textit{splits}, preprocesamiento o definición de etiquetas \cite{riverol_2024,fang_2024,zhou_2024}.
\end{itemize}

\section{Contribuciones}
\begin{itemize}[leftmargin=*]
  \item Pipeline reproducible de preprocesamiento y extracción de características HTTP (25 features) a partir de método, URI, body y un campo derivado de headers cuando esté disponible.
  \item Definición e implementación de etiquetas \texttt{label\_binary}, \texttt{label\_multiclass} y \texttt{label\_multilabel} para tareas comparables entre datasets.
  \item Política explícita para derivar \texttt{label\_multiclass} desde la anotación CAPEC multietiqueta de SR-BH mediante una regla determinista basada en severidad, conservando columnas auxiliares de auditoría para revisar la etiqueta primaria elegida.
  \item Adopción de un detector \textit{one-class} escalable basado en \texttt{SGDOneClassSVM} con aproximación de kernel, manteniendo el criterio de entrenar exclusivamente con tráfico normal.
  \item Baselines supervisados lineales implementados en \texttt{scikit-learn}: regresión logística para tareas multiclase y \textit{One-vs-Rest} para salidas multietiqueta, con configuraciones compatibles con el tamaño de los datasets y con análisis posterior por coeficientes.
  \item Selección opcional de hiperparámetros mediante búsquedas acotadas dentro del bloque de entrenamiento, usando el esquema que corresponda a cada script y guardando los resultados de búsqueda para trazabilidad \cite{sklearn_randomsearch,sklearn_halving}.
  \item Importancia de características basada en coeficientes estandarizados y ablaciones por grupos, manteniendo el análisis interpretativo dentro del presupuesto computacional del estudio.
  \item Documento técnico con análisis crítico de decisiones de modelado, compromisos entre costo y rendimiento, y amenazas a la validez.
\end{itemize}

\chapter{Fundamentación teórica}

Este capítulo presenta la fundamentación teórica que sustenta el trabajo. Se describe el protocolo HTTP como superficie de ataque, se caracterizan los enfoques de protección basados en firmas y en anomalías, y se exponen los fundamentos de los modelos de aprendizaje automático utilizados.

\section{El protocolo HTTP y la superficie de ataque}

El protocolo HTTP (\textit{HyperText Transfer Protocol}) es un protocolo de aplicación sin estado utilizado para la transferencia de información en la Web. La semántica común de HTTP, incluyendo recursos, métodos, campos de cabecera, códigos de estado y conceptos centrales del protocolo, se define en RFC~9110; la sintaxis y el procesamiento de mensajes HTTP/1.1 se especifican en RFC~9112 \cite{rfc9110,rfc9112}. En términos generales, una petición HTTP (\textit{request}) está compuesta por los siguientes elementos: un \textbf{método} que indica la acción solicitada (por ejemplo, \texttt{GET} para recuperar un recurso o \texttt{POST} para enviar datos), una \textbf{URI} (\textit{Uniform Resource Identifier}) que identifica el recurso objetivo incluyendo el \textit{path} y, opcionalmente, una cadena de parámetros (\textit{query string}), la \textbf{versión del protocolo}, un conjunto de \textbf{cabeceras} (\textit{headers}) que transportan metadatos de la comunicación y, en algunos métodos como \texttt{POST} o \texttt{PUT}, un \textbf{cuerpo} (\textit{body}) opcional que contiene los datos enviados por el cliente.

En el alcance de este trabajo, no se procesa necesariamente el mensaje HTTP completo tal como circula por la red. En cambio, se utilizan los campos disponibles y normalizados en los datasets analizados: principalmente el método, la URI, el cuerpo y, cuando está disponible, información derivada de cabeceras. Esta aclaración es importante para interpretar correctamente la representación de entrada utilizada por los modelos.

Los ataques web explotan los campos de la petición HTTP para manipular el comportamiento de la aplicación. Los patrones maliciosos tienden a manifestarse como: valores de longitud atípica en la URI o el cuerpo, estructuras inusuales en los parámetros, secuencias de codificación (\textit{percent-encoding}) empleadas para ofuscar instrucciones maliciosas, y presencia de tokens específicos asociados a familias de ataque como inyección SQL o XSS. Los WAF inspeccionan precisamente esos campos para aplicar políticas de filtrado \cite{owasp_waf,modsecurity_site}.

\section{WAF: enfoque por firmas frente a enfoque por anomalías}

Existen dos grandes aproximaciones para la detección de ataques en tráfico HTTP: los enfoques basados en firmas o reglas, y los enfoques basados en anomalías. Comprender sus ventajas y limitaciones es fundamental para justificar el uso de aprendizaje automático en este dominio.

\subsection{Enfoque basado en firmas o reglas}

En un WAF basado en firmas, el sistema mantiene un conjunto de patrones predefinidos que describen cómo se manifiestan los ataques conocidos. Cada petición entrante se compara contra esos patrones; si coincide con alguno, se bloquea o se genera una alerta. ModSecurity es un ejemplo representativo y ampliamente utilizado de este enfoque: actúa como motor de reglas configurable que permite definir, actualizar y aplicar políticas de filtrado sobre tráfico HTTP real \cite{modsecurity_site}. Se menciona aquí porque ilustra con claridad el paradigma de protección basado en firmas, útil para contrastar con los enfoques basados en anomalías y en aprendizaje automático que se estudian en este trabajo.

\textbf{Ventajas:} alta precisión para ataques conocidos y cubiertos por las reglas; bajo índice de falsos positivos cuando las reglas están bien ajustadas; razonamiento directo sobre la causa del bloqueo.

\textbf{Limitaciones:} requiere mantenimiento continuo a medida que surgen nuevas técnicas de ataque; puede fallar ante variantes o mutaciones de ataques conocidos para los cuales no existe todavía una regla; el catálogo de firmas puede volverse extenso y costoso de gestionar.

\subsection{Enfoque basado en anomalías}

En un WAF basado en anomalías, el sistema no depende de patrones predefinidos de ataque. En cambio, aprende una representación del tráfico ``normal'' y detecta como sospechosas las peticiones que se desvían de esa representación. Esto lo hace útil cuando la cobertura de ataques conocidos es incompleta o cambiante, ya que puede señalar actividad inusual sin necesidad de haber visto antes ese tipo de ataque específico.

En el dominio web, trabajos clásicos como el de Kruegel y Vigna demostraron que atributos estadísticos y estructurales de las peticiones HTTP —como longitudes, distribución de caracteres o frecuencia de parámetros— permiten construir modelos de normalidad efectivos para puntuar anomalías \cite{kruegel_vigna_2003}.

\textbf{Ventajas:} capacidad de detectar ataques novedosos o variantes no contempladas en un catálogo; no requiere actualización manual de reglas; puede adaptarse a cambios en el perfil del tráfico normal.

\textbf{Limitaciones:} puede generar mayor tasa de falsos positivos si el tráfico normal es muy diverso o variable; las decisiones son más difíciles de explicar que en el enfoque de firmas; requiere datos de entrenamiento representativos del tráfico legítimo.

\section{CAPEC como taxonomía de patrones de ataque}

CAPEC (\textit{Common Attack Pattern Enumeration and Classification}) es una taxonomía de patrones de ataque mantenida por MITRE que describe, de forma estructurada, los mecanismos mediante los cuales se explotan vulnerabilidades en sistemas de software \cite{capec_home,capec_about}. A diferencia de una lista de vulnerabilidades, CAPEC no enumera vulnerabilidades concretas de un producto; enumera \textit{patrones de ataque}, es decir, formas recurrentes en que un atacante puede organizar sus acciones para comprometer un sistema.

La utilidad de CAPEC para este trabajo está en que aporta una semántica común para nombrar ataques web. Por ejemplo, una petición con una carga como \texttt{' OR 1=1 --} puede asociarse a un patrón de inyección SQL; una petición que incluye \texttt{../} puede relacionarse con traversal de rutas; y una petición con \texttt{<script>} u otros manejadores de eventos puede relacionarse con \textit{cross-site scripting}. Estos ejemplos no significan que la presencia literal de esos textos sea suficiente para clasificar una petición: solo ilustran cómo ciertos patrones observables del tráfico HTTP pueden vincularse con familias de ataque.

CAPEC también debe leerse con cuidado porque es jerárquico. El esquema incluye niveles más generales y más específicos: algunas entradas funcionan como categorías o patrones abstractos, mientras que otras describen patrones más concretos \cite{capec_schema_docs,capec_schema_desc}. Por eso, dos etiquetas CAPEC pueden no estar al mismo nivel de granularidad. Esta característica es importante para la evaluación: en clasificación multietiqueta se conserva el conjunto de etiquetas activas por petición, pero en clasificación multiclase se debe seleccionar una sola etiqueta primaria. Esa reducción simplifica el problema y permite comparar modelos multiclase, aunque pierde parte de la estructura jerárquica y semántica del ataque.

En este trabajo, CAPEC se utiliza como esquema semántico para clasificar los ataques presentes en el dataset SR-BH~2020. Cada petición maliciosa del dataset puede estar anotada con una o varias etiquetas CAPEC, lo que habilita la formulación multietiqueta del problema. Antes de usar estas etiquetas en las secciones matemáticas y metodológicas que siguen, conviene tener presente que una etiqueta CAPEC identifica un patrón de ataque, no una vulnerabilidad ni un vector técnico específico.

\section{Fundamentos matemáticos de los modelos}

Esta sección presenta la formulación matemática de los modelos utilizados y fija una notación única
para que las variables de las ecuaciones correspondan de manera directa con la implementación.

\subsection{Tarea de clasificación en aprendizaje automático}

En aprendizaje automático, una tarea de clasificación consiste en lo siguiente: dado un conjunto de ejemplos etiquetados, donde cada ejemplo tiene un vector de características de entrada y una o varias etiquetas conocidas como salida, el objetivo es aprender una función de decisión que permita asignar una clase o conjunto de etiquetas a nuevos ejemplos no vistos durante el entrenamiento. La calidad del modelo se evalúa comparando sus predicciones con las etiquetas reales en un conjunto de prueba separado.

En el contexto de este trabajo, cada ejemplo es una petición HTTP (\textit{request}), el vector de características (\textit{features}) se extrae de los campos del protocolo, y las etiquetas pueden ser binarias (normal/ataque), multiclase (tipo de ataque) o multietiqueta (conjunto de patrones CAPEC activos).


\subsection{Notación y definición exacta de variables}

\subsubsection{Dataset, features y objetivos}
Sea un conjunto de $n$ peticiones HTTP (requests). Cada request se transforma a un vector numérico de $d$ características (features)
extraídas de método, URI y cuerpo. Definimos:

\begin{itemize}[leftmargin=*]
  \item $x_i \in \R^d$: vector de features de la request $i$ (una fila del dataset).
  \item $X \in \R^{n \times d}$: matriz de diseño; su fila $i$ es $x_i^\top$ (todas las requests juntas).
\end{itemize}

La variable objetivo depende de la tarea:

\begin{itemize}[leftmargin=*]
  \item \textbf{Binaria:} $y_i \in \{0,1\}$, donde $y_i=0$ indica tráfico normal y $y_i=1$ indica ataque/anómalo.
        El vector completo es $y \in \{0,1\}^n$.
  \item \textbf{Multiclase:} $y_i \in \{1,\dots,K\}$, donde $K$ es el número de clases (por ejemplo, \texttt{NORMAL} + clases CAPEC o un agrupamiento).
        El vector completo es $y \in \{1,\dots,K\}^n$.
  \item \textbf{Multietiqueta:} $y_i \in \{0,1\}^L$, donde $L$ es el número de etiquetas posibles.
        El componente $y_{i\ell}=1$ indica que la request $i$ pertenece a la etiqueta $\ell$.
        El conjunto completo se representa como $Y \in \{0,1\}^{n \times L}$.
\end{itemize}

\noindent\textbf{Nota (disponibilidad de etiquetas):} no todos los datasets soportan todas las tareas. CSIC~2010 ofrece solo binario (normal/anómalo); TorpEda~2012 ofrece multiclase (tipo de ataque) y su versión binaria por colapso; SR-BH~2020 ofrece binario, multiclase (vía etiqueta primaria) y multietiqueta CAPEC. En consecuencia, al ejecutar los modelos sobre los tres datasets, las tareas multiclase/multietiqueta se evalúan en su versión reducida cuando el dataset no provee la granularidad requerida.

Además, usamos:

\begin{itemize}[leftmargin=*]
  \item $\hat{y}$ o $\hat{Y}$: predicciones del modelo (mismo tipo/forma que $y$ o $Y$).
  \item $\ind[\cdot]$: función indicadora; vale 1 si la condición es verdadera y 0 si es falsa.
\end{itemize}

\subsubsection{Tabla de símbolos (resumen)}
{\footnotesize
\renewcommand{\arraystretch}{1.1}
\begin{longtable}{@{}L{3.6cm}L{11.6cm}@{}}
\toprule
\textbf{Símbolo} & \textbf{Significado exacto} \\
\midrule
\endfirsthead
\toprule
\textbf{Símbolo} & \textbf{Significado exacto} \\
\midrule
\endhead
\multicolumn{2}{@{}l}{\textit{Dataset y dimensiones}}\\
$n$ & cantidad total de requests del dataset \\
$d$ & cantidad de features numéricas por request \\
$K$ & cantidad de clases en multiclase \\
$L$ & cantidad de etiquetas en multietiqueta \\
$x_i \in \R^d$ & features de la request $i$ \\
$X \in \R^{n\times d}$ & matriz con todas las requests (filas) y features (columnas) \\
$y_i$ & etiqueta (binaria o multiclase) de la request $i$ \\
$y \in \{0,1\}^n$ & vector de etiquetas binarias de todas las requests \\
$y \in \{1,\dots,K\}^n$ & vector de etiquetas multiclase de todas las requests \\
$y_i \in \{0,1\}^L$ & vector de etiquetas (multietiqueta) para la request $i$ \\
$y_{i\ell}$ & componente $\ell$ del vector multietiqueta de la request $i$ \\
$Y \in \{0,1\}^{n\times L}$ & matriz multietiqueta para todo el dataset \\
$\hat{y}, \hat{Y}$ & predicciones del modelo \\
$\ind[\cdot]$ & función indicadora \\
\midrule
\multicolumn{2}{@{}l}{\textit{Detector one-class}}\\
$f(x)$ & función de decisión del detector one-class \\
$z(x)$ & transformación de $x$ al espacio inducido por el kernel \\
$w$ & vector de pesos que define la orientación de la frontera \\
$\rho$ & umbral de separación de la frontera one-class \\
$\operatorname{sign}(\cdot)$ & función signo aplicada a la puntuación de normalidad \\
\midrule
\multicolumn{2}{@{}l}{\textit{Regresión logística (binaria y multiclase)}}\\
$z$ & puntaje lineal $z=w^\top x + b$ \\
$b$ & sesgo (intercepto) del modelo lineal \\
$\sigma(z)$ & función sigmoide, $\sigma(z)=1/(1+e^{-z})$ \\
$P(y=1\mid x)$ & probabilidad estimada de la clase positiva \\
$p_i$ & probabilidad estimada $P(y=1\mid x_i)$ para la request $i$ \\
$\mathcal{L}(w,b)$ & función de pérdida (log-loss / entropía cruzada) \\
$\lambda$ & coeficiente de regularización $L_2$ \\
$\lVert w\rVert^2$ & norma euclídea al cuadrado del vector de pesos \\
$s_k(x)$ & puntaje lineal de la clase $k$: $s_k(x)=w_k^\top x + b_k$ \\
$w_k,\ b_k$ & pesos y sesgo asociados a la clase $k$ \\
$P(y=k\mid x)$ & probabilidad estimada de pertenecer a la clase $k$ (softmax) \\
\midrule
\multicolumn{2}{@{}l}{\textit{One-vs-Rest (multietiqueta)}}\\
$w_\ell,\ b_\ell$ & pesos y sesgo del clasificador binario para la etiqueta $\ell$ \\
$P(y_\ell=1\mid x)$ & probabilidad estimada de la etiqueta $\ell$ \\
$\tau_\ell$ & umbral de decisión binaria para la etiqueta $\ell$ \\
$\hat{y}_\ell$ & predicción binaria final para la etiqueta $\ell$ \\
\midrule
\multicolumn{2}{@{}l}{\textit{Métricas binarias}}\\
$TP,\ FP,\ FN,\ TN$ & verdaderos/falsos positivos y verdaderos/falsos negativos \\
$\mathrm{TPR}$ & tasa de verdaderos positivos (recall), $TP/(TP+FN)$ \\
$\mathrm{TNR}$ & tasa de verdaderos negativos (specificity), $TN/(TN+FP)$ \\
$\mathrm{FPR}$ & tasa de falsos positivos, $FP/(FP+TN)=1-\mathrm{TNR}$ \\
$\mathrm{FNR}$ & tasa de falsos negativos, $FN/(FN+TP)=1-\mathrm{TPR}$ \\
$F_1$ & media armónica de precision y recall \\
$\mathrm{MCC}$ & coeficiente de correlación de Matthews \\
\midrule
\multicolumn{2}{@{}l}{\textit{Métricas multiclase}}\\
$TP_k,\ FP_k,\ FN_k$ & conteos one-vs-rest correspondientes a la clase $k$ \\
$\mathrm{Precision}_k,\ \mathrm{Recall}_k,\ F_{1,k}$ & métricas calculadas sobre la clase $k$ \\
$s_k$ & soporte de la clase $k$, con $\sum_k s_k = n$ \\
$F_{1,micro},\ F_{1,macro},\ F_{1,weighted}$ & promedios micro, macro y ponderado del $F_1$ \\
$\mathrm{AUC}_k$ & ROC-AUC one-vs-rest de la clase $k$ \\
\midrule
\multicolumn{2}{@{}l}{\textit{Métricas multietiqueta}}\\
$TP_\ell,\ FP_\ell,\ FN_\ell$ & conteos por etiqueta $\ell$ sobre las celdas $(i,\ell)$ \\
$Y_i,\ \hat{Y}_i$ & conjuntos de etiquetas verdaderas y predichas de la request $i$ \\
$J_i$ & Jaccard por instancia, $|Y_i\cap\hat{Y}_i|/|Y_i\cup\hat{Y}_i|$ \\
$J_{samples},\ J_{micro},\ J_{macro}$ & agregaciones del Jaccard por muestra, micro y macro \\
$\mathrm{HL}$ & Hamming Loss \\
$\mathrm{EMR}$ & Exact Match Ratio (subset accuracy) \\
\midrule
\multicolumn{2}{@{}l}{\textit{Métricas temporales y operativas}}\\
$t_{opt}$ & tiempo total del bloque de búsqueda de hiperparámetros \\
$t_{build}$ & tiempo de ajuste final del modelo (\texttt{train\_time\_seconds}) \\
$t_{apply}$ & tiempo total del benchmark operativo (\texttt{wall\_time\_seconds}) \\
$\mathrm{Throughput}$ & requests procesados por segundo durante el benchmark \\
\bottomrule
\end{longtable}
}

\subsection{Detector one-class escalable}

\subsubsection{definición}

Un detector \textit{one-class} es un modelo diseñado para aprender el comportamiento habitual de una sola clase de datos. En problemas de detección de anomalías, esa clase suele corresponder al comportamiento considerado normal. La idea central no es aprender todas las formas posibles de ataque, sino construir una representación del tráfico esperado y detectar como anómalas aquellas peticiones que se alejan de ese patrón.

A diferencia de un clasificador binario tradicional, que necesita ejemplos de tráfico normal y de tráfico malicioso para aprender una separación entre ambas clases, el enfoque \textit{one-class} parte principalmente de una pregunta distinta: ¿qué características definen a las peticiones normales? Una vez aprendida esa región de normalidad, cualquier observación suficientemente diferente puede ser tratada como sospechosa.

Esta lógica es especialmente útil en escenarios donde las anomalías pueden ser variadas, poco frecuentes o difíciles de anticipar. En tráfico HTTP, por ejemplo, una petición puede diferir del comportamiento habitual por tener una estructura de URL inusual, una cantidad anormal de parámetros, longitudes atípicas, codificaciones extrañas o combinaciones de características poco frecuentes. El detector no necesita conocer de antemano el nombre exacto del ataque; basta con que la petición se aparte del patrón aprendido como normal.

\subsubsection{Modelo conceptual}

Sea $x \in \R^d$ el vector de características que representa una petición. Cada componente de $x$ describe una propiedad observable de esa petición, como una longitud, un conteo, una proporción o un indicador binario. El objetivo del detector \textit{one-class} es aprender una región del espacio de características donde se concentran las observaciones normales.

El modelo construye una frontera alrededor de los datos normales. Las peticiones que quedan dentro de esa frontera se consideran compatibles con el comportamiento esperado, mientras que las que quedan fuera se consideran anomalías.

La Ecuación~\eqref{eq:ocsvm-decision} expresa una forma general de la decisión del modelo:

\begin{equation}
\label{eq:ocsvm-decision}
f(x)=\operatorname{sign}\big(w^\top z(x)-\rho\big).
\end{equation}

En esta expresión, $z(x)$ representa una transformación del vector original $x$ hacia otro espacio de características. Esta transformación permite describir relaciones más complejas que las que podrían observarse directamente en el espacio original. El vector $w$ define la orientación de la frontera de decisión y $\rho$ actúa como umbral de separación.

El término interno de la Ecuación~\eqref{eq:ocsvm-decision}, mostrado por separado en la Ecuación~\eqref{eq:ocsvm-score}, puede interpretarse como una puntuación de normalidad:

\begin{equation}
\label{eq:ocsvm-score}
w^\top z(x)-\rho
\end{equation} Si el valor queda del lado aceptado por la frontera, la petición se considera normal. Si queda del lado opuesto, se considera anómala. Por eso, el modelo no solo produce una etiqueta final, sino también una noción de distancia respecto al comportamiento aprendido.

\subsubsection{Frontera de decisión y anomalías}

La frontera de decisión separa la región donde se concentra el tráfico normal del resto del espacio de características. En un caso simple, esa frontera podría ser aproximadamente lineal. Sin embargo, en datos reales, el comportamiento normal no siempre se distribuye de forma simple. Puede haber varias zonas de normalidad, relaciones no lineales entre variables o combinaciones de características que solo son normales bajo ciertos contextos.

Por ese motivo, los detectores \textit{one-class} suelen apoyarse en transformaciones de características o funciones de kernel. Estas permiten que una frontera aparentemente simple en un espacio transformado represente una separación más flexible en el espacio original. Así, el modelo puede adaptarse mejor a patrones donde la normalidad no forma una nube compacta y lineal.

La anomalía, en este contexto, no significa necesariamente ataque confirmado. Significa que la petición se aleja del comportamiento aprendido como habitual. Por lo tanto, la salida del modelo debe interpretarse como una señal de sospecha o desviación, no como una identificación directa del tipo de ataque.

\subsubsection{Escalabilidad del enfoque}

El enfoque \textit{one-class} puede volverse costoso cuando se trabaja con grandes volúmenes de datos, especialmente si se utilizan métodos basados en kernels completos. Esto ocurre porque algunos modelos necesitan comparar muchas observaciones entre sí para construir la frontera de decisión, lo que incrementa el costo computacional a medida que crece el número de peticiones.

Una alternativa escalable consiste en aproximar el efecto de un kernel mediante una representación explícita de menor costo. En lugar de trabajar directamente con todas las comparaciones entre ejemplos, se transforma cada petición a un espacio donde el modelo pueda operar de manera más eficiente. De esta forma, se busca mantener parte de la flexibilidad de los métodos no lineales, pero reduciendo el costo de entrenamiento y aplicación.

La escalabilidad, entonces, no cambia la idea central del detector \textit{one-class}. El modelo sigue aprendiendo una región de normalidad y marcando como anómalas las observaciones que quedan fuera de ella. Lo que cambia es la forma computacional de aproximar esa frontera para que el enfoque pueda aplicarse a conjuntos de datos más grandes.

\subsubsection{Interpretación de la salida}

La salida de un detector \textit{one-class} suele interpretarse en dos niveles. El primero es la etiqueta final: normal o anómala. El segundo es la puntuación asociada a la decisión, que indica qué tan cerca o lejos se encuentra una petición respecto de la frontera aprendida.

Una petición muy cercana a la región normal tendrá una puntuación compatible con el comportamiento esperado. En cambio, una petición muy alejada de esa región tendrá una puntuación más anómala. Esta diferencia es importante porque no todas las anomalías tienen la misma intensidad: algunas pueden estar apenas fuera de la frontera, mientras que otras pueden ser claramente atípicas.

En términos prácticos, esta puntuación permite ordenar las peticiones según su grado de sospecha. Sin embargo, el modelo no explica por sí solo cuál es la causa exacta de la anomalía. Para interpretar el resultado, es necesario analizar las características de la petición y observar qué aspectos la separan del patrón normal.

\subsubsection{Fortalezas y limitaciones}

Una fortaleza del enfoque \textit{one-class} es que no depende de tener ejemplos rotulados de todos los ataques posibles. Esto lo hace útil cuando las anomalías son raras, cambiantes o difíciles de etiquetar. Además, puede detectar comportamientos novedosos, siempre que estos se aparten suficientemente de la normalidad aprendida.

Sin embargo, también presenta limitaciones. Si el tráfico normal usado como referencia es incompleto o contiene mucha variabilidad, la frontera aprendida puede ser demasiado restrictiva o demasiado amplia. Una frontera muy restrictiva puede generar falsos positivos, marcando como anómalas peticiones legítimas. Una frontera demasiado amplia puede producir falsos negativos, aceptando como normales peticiones que en realidad son sospechosas.

Otra limitación importante es que el modelo detecta desviaciones, pero no clasifica el tipo específico de ataque. Por ello, un detector \textit{one-class} es adecuado para identificar comportamientos anómalos, pero no reemplaza a un clasificador multiclase o multietiqueta cuando el objetivo es asignar una categoría concreta de ataque.

\subsection{Regresión logística}

\subsubsection{Caso binario y función sigmoide}

La regresión logística es un modelo supervisado de clasificación. En el caso
binario, la tarea consiste en decidir entre dos categorías: por ejemplo,
si una petición HTTP corresponde a tráfico normal o a un ataque.

El modelo procede en dos pasos. Primero, a partir de una petición
representada por un vector de características $x$, calcula el puntaje lineal de la Ecuación~\eqref{eq:logreg-linear-score}:
\begin{equation}
\label{eq:logreg-linear-score}
z=w^\top x+b,
\end{equation}
donde $w$ es el vector de pesos y $b$ es el sesgo. Este puntaje no es más
que una combinación ponderada de las características de la petición. A modo
de ejemplo, si una característica indica la presencia de la palabra clave
\texttt{SELECT} en la URL y su peso asociado es positivo, esa característica
empuja el puntaje hacia arriba; si otra característica indica que el método
HTTP es \texttt{GET} y su peso es negativo, la empuja hacia abajo.

El puntaje $z$ puede tomar cualquier valor real: puede ser muy negativo,
cercano a cero o muy positivo. Para interpretarlo como una probabilidad
entre 0 y 1, se aplica la \textit{función sigmoide} de la Ecuación~\eqref{eq:logreg-sigmoid}:
\begin{equation}
\label{eq:logreg-sigmoid}
P(y=1\mid x)=\sigma(z)=\frac{1}{1+e^{-z}}.
\end{equation}
La sigmoide es una curva en forma de S que comprime cualquier número real
al intervalo $(0,1)$. Algunos valores concretos muestran su comportamiento:
$\sigma(-3)\approx 0{,}05$, $\sigma(0)=0{,}5$ y $\sigma(3)\approx 0{,}95$.
En la práctica, un puntaje muy negativo se traduce en ``el modelo está casi
seguro de que la petición es normal'', un puntaje cercano a cero en ``el modelo
no está seguro'' y un puntaje muy positivo en ``el modelo está casi seguro de
que es un ataque''.

En este trabajo, cuando la tarea es binaria, la clase positiva ($y=1$)
representa ataque o anomalía y la clase negativa ($y=0$) representa tráfico
normal. La Figura~\ref{fig:sigmoid} ilustra esta transformación de puntaje a
probabilidad. Es importante notar que la salida no es solo una etiqueta, sino
también un grado de confianza interpretable directamente como probabilidad
estimada.

\begin{figure}[H]
\centering
\begin{tikzpicture}
\begin{axis}[
  width=0.70\textwidth,
  height=0.35\textwidth,
  xlabel={$z$},
  ylabel={$\sigma(z)$},
  domain=-6:6,
  samples=200,
  grid=both
]
\addplot {1/(1+exp(-x))};
\end{axis}
\end{tikzpicture}
\caption{Función sigmoide usada por regresión logística para transformar un
puntaje lineal en una probabilidad \cite{islr_logreg}.}
\label{fig:sigmoid}
\end{figure}

\subsubsection{Función de pérdida binaria}

Una vez que el modelo produce una probabilidad, hace falta una manera de
medir qué tan buenas son sus predicciones, de forma que los pesos $w$ y el
sesgo $b$ puedan ajustarse. Esa medida se llama \textit{función de pérdida}.

La idea es simple: el modelo debe pagar un costo alto cuando se
equivoca con seguridad, y un costo bajo cuando acierta con seguridad.
Considérese una petición que realmente es un ataque ($y=1$):
\begin{itemize}
  \item si el modelo le asigna una probabilidad de $0{,}99$ a la clase
        ataque, la pérdida es muy pequeña, porque acertó con confianza;
  \item si le asigna $0{,}5$, la pérdida es intermedia, porque dudó;
  \item si le asigna $0{,}01$, la pérdida es muy grande, porque estaba
        seguro y se equivocó.
\end{itemize}

Esta intuición se formaliza en la pérdida logística o \textit{log-loss}.
Dado un conjunto $\{(x_i,y_i)\}_{i=1}^n$ con $y_i\in\{0,1\}$ y
$p_i=P(y=1\mid x_i)$, su forma regularizada se presenta en la Ecuación~\eqref{eq:logreg-binary-loss}:
\begin{equation}
\label{eq:logreg-binary-loss}
\mathcal{L}(w,b)=
-\frac{1}{n}\sum_{i=1}^n
\Big(y_i\log(p_i) + (1-y_i)\log(1-p_i)\Big)
+ \frac{\lambda}{2}\lVert w\rVert^2.
\end{equation}
En el primer término, para cada ejemplo solo uno de los dos sumandos queda
activo: si $y_i=1$ se evalúa $-\log(p_i)$, que crece cuando $p_i$ es bajo;
si $y_i=0$ se evalúa $-\log(1-p_i)$, que crece cuando $p_i$ es alto. El
segundo término, conocido como regularización $L2$, penaliza pesos muy
grandes. Esto evita que el modelo se ajuste demasiado a los datos de
entrenamiento y favorece soluciones más estables, sobre todo cuando hay
muchas características o las clases están desbalanceadas
\cite{islr_logreg,sklearn_logreg}.

\subsubsection{Extensión multiclase mediante softmax}

Cuando hay más de dos clases, no alcanza con estimar una sola probabilidad.
El modelo debe asignar una probabilidad a cada clase posible, y esas
probabilidades deben sumar 1.

Para lograrlo, la regresión logística multiclase calcula un puntaje lineal
por cada clase, como se indica en la Ecuación~\eqref{eq:softmax-score}:
\begin{equation}
\label{eq:softmax-score}
s_k(x)=w_k^\top x + b_k, \qquad k=1,\dots,K.
\end{equation}
Cada clase tiene su propio vector de pesos $w_k$ y su propio sesgo $b_k$.
Los $K$ puntajes resultantes se convierten luego en una distribución de
probabilidades mediante la \textit{función softmax} definida en la Ecuación~\eqref{eq:softmax-probability}:
\begin{equation}
\label{eq:softmax-probability}
P(y=k\mid x)=\frac{\exp(s_k(x))}{\sum_{j=1}^K \exp(s_j(x))}.
\end{equation}
La softmax puede verse como una generalización de la sigmoide al caso de
varias clases: toma todos los puntajes y los normaliza para que sean
positivos y sumen 1. La exponencial cumple dos funciones simultáneas:
garantiza positividad y amplifica las diferencias, de modo que la clase con
mayor puntaje concentra una porción notablemente más grande de la
probabilidad total.

A modo de ejemplo, supóngase un problema con tres clases: \texttt{NORMAL},
\texttt{SQLi} y \texttt{XSS}. Para una petición particular, el modelo podría
producir los puntajes $s_{\texttt{NORMAL}}=0{,}5$, $s_{\texttt{SQLi}}=2{,}5$
y $s_{\texttt{XSS}}=1{,}0$. Tras aplicar softmax, esos puntajes se traducen
en probabilidades aproximadas de $0{,}10$, $0{,}70$ y $0{,}20$,
respectivamente. La clase predicha es \texttt{SQLi}, porque concentra la
mayor probabilidad; el modelo también informa que existe un $20\%$ de
posibilidad de que la petición sea \texttt{XSS} y un $10\%$ de que sea
\texttt{NORMAL}.

Existen otras formas de abordar problemas con varias clases. Una alternativa
frecuente es One-vs-Rest, que entrena un clasificador binario independiente
por cada clase. Sin embargo, cuando cada petición pertenece exactamente a
una sola clase, la formulación multinomial con softmax resulta más natural,
porque modela todas las clases de manera conjunta y garantiza que la salida
sea una única distribución sobre las $K$ clases
\cite{sklearn_logreg,islr_logreg}.

La pérdida asociada al caso multiclase es la \textit{entropía cruzada}, formulada en la Ecuación~\eqref{eq:softmax-loss}:
\begin{equation}
\label{eq:softmax-loss}
\mathcal{L}=
-\frac{1}{n}\sum_{i=1}^{n}\log P(y=y_i\mid x_i)
+ \frac{\lambda}{2}\sum_{k=1}^{K}\lVert w_k\rVert^2.
\end{equation}
La interpretación es análoga al caso binario: para cada ejemplo solo
interesa la probabilidad que el modelo asignó a la clase verdadera; si esa
probabilidad es alta, la pérdida es pequeña, y si es baja, la pérdida es
grande. La regularización cumple el mismo papel que antes: limitar pesos
extremos y favorecer modelos más estables.

\subsection{Clasificación multietiqueta con One-vs-Rest (OvR)}

En clasificación multiclase, cada petición recibe una sola clase final. En clasificación multietiqueta, en cambio, una misma petición puede activar varias etiquetas al mismo tiempo. Esto es importante para SR-BH~2020 porque una petición HTTP puede estar asociada a más de un patrón CAPEC.

One-vs-Rest resuelve el problema multietiqueta entrenando un clasificador binario por etiqueta. Para cada etiqueta $\ell\in\{1,\dots,L\}$, el modelo aprende si esa etiqueta está presente o ausente mediante la probabilidad de la Ecuación~\eqref{eq:ovr-label-probability}:
\begin{equation}
\label{eq:ovr-label-probability}
P(y_\ell=1\mid x)=\sigma(w_\ell^\top x+b_\ell).
\end{equation}

Después, cada probabilidad se convierte en una decisión binaria mediante un umbral $\tau_\ell$, como se define en la Ecuación~\eqref{eq:ovr-threshold}:
\begin{equation}
\label{eq:ovr-threshold}
\hat{y}_{\ell}=\ind\big[P(y_\ell=1\mid x)\ge \tau_\ell\big].
\end{equation}

La idea puede leerse con un ejemplo. Si las etiquetas posibles fueran \texttt{SQL Injection}, \texttt{Path Traversal} y \texttt{Command Injection}, OvR entrenaría tres clasificadores: uno pregunta si la petición es o no \texttt{SQL Injection}, otro si es o no \texttt{Path Traversal}, y otro si es o no \texttt{Command Injection}. Como las predicciones son independientes por etiqueta, la salida final puede contener una, varias o ninguna etiqueta.

\begin{figure}[H]
\centering
\resizebox{0.95\textwidth}{!}{%
\begin{tikzpicture}[
  block/.style={draw, rounded corners, align=center, minimum width=3.2cm, minimum height=0.9cm},
  arrow/.style={-Latex, thick}
]
\node[block] (x) {$x$\\(características HTTP)};
\node[block, right=1.8cm of x] (c1) {Clasificador\\etiqueta $1$};
\node[block, above=0.7cm of c1] (c2) {Clasificador\\etiqueta $2$};
\node[block, below=0.7cm of c1] (c3) {Clasificador\\etiqueta $L$};

\node[block, right=2.1cm of c1] (thr) {Umbrales\\por etiqueta};
\node[block, right=1.8cm of thr] (y) {$\hat{y}\in\{0,1\}^L$\\salida multietiqueta};

\draw[arrow] (x) -- (c1);
\draw[arrow] (x) -- (c2);
\draw[arrow] (x) -- (c3);

\draw[arrow] (c2) -- (thr);
\draw[arrow] (c1) -- (thr);
\draw[arrow] (c3) -- (thr);

\draw[arrow] (thr) -- (y);
\end{tikzpicture}%
}
\caption{Esquema One-vs-Rest para clasificación multietiqueta: una misma petición se evalúa contra varios clasificadores binarios, uno por etiqueta \cite{tsoumakas_multilabel_overview,rifkin_ova}.}
\label{fig:ovr}
\end{figure}

La Figura~\ref{fig:ovr} resume este flujo. El vector $x$ entra a varios clasificadores binarios; cada clasificador produce una probabilidad asociada a una etiqueta; luego se aplican umbrales y se obtiene un vector de salida. Esta salida no es una única clase, sino una lista de predicciones sí/no, una por etiqueta.

\subsection{Métricas y base matemática}

Esta sección presenta las métricas usadas para evaluar los modelos de
detección. Antes de definir cada una conviene introducir los bloques
básicos que las componen, ya que casi todas se construyen sobre los
mismos conceptos: primero el caso binario, luego la generalización a
multiclase y finalmente el caso multietiqueta. Las fórmulas incluidas no
se proponen como contribuciones nuevas de este trabajo: corresponden a
medidas estándar de evaluación de clasificadores, sistematizadas en la
literatura para tareas binarias, multiclase y multietiqueta
\cite{sokolova_lapalme_2009,powers2011}. Dentro de cada bloque
reutilizamos un mismo ejemplo numérico a lo largo de las distintas
métricas.

\subsubsection{Clasificación binaria}

\paragraph{La matriz de confusión}

Cuando un clasificador binario distingue entre dos clases (por ejemplo,
\emph{ataque} vs.\ \emph{normal}), cada predicción cae en una de cuatro
casillas:
\begin{itemize}[leftmargin=*]
  \item $TP$ (\emph{true positives}): ataques correctamente predichos como ataque.
  \item $FP$ (\emph{false positives}): tráfico normal predicho erróneamente como ataque.
  \item $FN$ (\emph{false negatives}): ataques predichos erróneamente como normal.
  \item $TN$ (\emph{true negatives}): tráfico normal correctamente predicho como normal.
\end{itemize}

\paragraph{Ejemplo de referencia.} Supongamos un conjunto de prueba con
100 requests, de los cuales 20 son ataques. El clasificador detecta
correctamente 16 ataques ($TP=16$), pasa por alto 4 ($FN=4$), confunde 5
normales con ataques ($FP=5$) y clasifica bien los 75 restantes
($TN=75$). Este ejemplo se reutiliza en las métricas siguientes.

\paragraph{Accuracy}

\textbf{Fórmula.} La Ecuación~\eqref{eq:binary-accuracy} define la accuracy binaria \cite{sokolova_lapalme_2009}:
\begin{equation}
\label{eq:binary-accuracy}
\mathrm{Accuracy}=\frac{TP+TN}{TP+TN+FP+FN}.
\end{equation}

\textbf{Utilidad.} Mide la proporción global de aciertos. Es la métrica
más intuitiva, pero engaña cuando las clases están desbalanceadas: si el
99\% de las requests son normales, predecir siempre ``normal'' alcanza
un 99\% de accuracy sin detectar ni un solo ataque.

\textbf{Ejemplo.} El cálculo con los conteos del ejemplo se muestra en la Ecuación~\eqref{eq:binary-accuracy-example}.
\begin{equation}
\label{eq:binary-accuracy-example}
\mathrm{Accuracy}=\frac{16+75}{100}=0{,}91.
\end{equation}

\paragraph{Precision}

\textbf{Fórmula.} La Ecuación~\eqref{eq:binary-precision} define la precision \cite{sokolova_lapalme_2009,powers2011}:
\begin{equation}
\label{eq:binary-precision}
\mathrm{Precision}=\frac{TP}{TP+FP}.
\end{equation}

\textbf{Utilidad.} De todas las requests que el modelo etiquetó como
ataque, qué proporción \emph{realmente} lo era. Una precision alta
significa pocas falsas alarmas, lo cual es deseable cuando cada alerta
tiene un costo (por ejemplo, una intervención humana o el bloqueo de un
usuario legítimo).

\textbf{Ejemplo.} El modelo marcó como ataque a $TP+FP=21$ requests, de
las cuales 16 lo eran realmente. El cálculo aparece en la Ecuación~\eqref{eq:binary-precision-example}:
\begin{equation}
\label{eq:binary-precision-example}
\mathrm{Precision}=\frac{16}{21}\approx 0{,}76.
\end{equation}
Aproximadamente 3 de cada 4 alertas son legítimas.

\paragraph{Recall (sensibilidad o TPR)}

\textbf{Fórmula.} La Ecuación~\eqref{eq:binary-recall} define recall o TPR \cite{sokolova_lapalme_2009,powers2011}:
\begin{equation}
\label{eq:binary-recall}
\mathrm{Recall}=\mathrm{TPR}=\frac{TP}{TP+FN}.
\end{equation}

\textbf{Utilidad.} De todos los ataques que existían realmente, qué
proporción detectó el modelo. En seguridad suele ser la métrica más
crítica: un ataque no detectado puede tener consecuencias graves.

\textbf{Ejemplo.} Existían 20 ataques y se detectaron 16. La Ecuación~\eqref{eq:binary-recall-example} muestra el cálculo:
\begin{equation}
\label{eq:binary-recall-example}
\mathrm{Recall}=\frac{16}{20}=0{,}80.
\end{equation}
El 80\% de los ataques fue detectado; el 20\% restante pasó inadvertido.

\paragraph{$F_1$-score}

Para entender $F_1$ es necesario haber visto antes precision y recall,
porque $F_1$ los combina. La Ecuación~\eqref{eq:binary-f1} muestra esa combinación, usada de forma estándar como $F$-measure \cite{powers2011}.

\textbf{Fórmula:}
\begin{equation}
\label{eq:binary-f1}
F_1=\frac{2\cdot \mathrm{Precision}\cdot \mathrm{Recall}}{\mathrm{Precision}+\mathrm{Recall}}.
\end{equation}

\textbf{Utilidad.} Es la \emph{media armónica} de precision y recall. Se
usa la media armónica (en lugar de la aritmética) porque penaliza
fuertemente cuando una de las dos es baja: si una es 1 y la otra es 0,
el promedio aritmético sería 0{,}5, pero $F_1$ es 0. Esto evita inflar
el resultado optimizando solo una de las dos.

\textbf{Ejemplo.} Con $\mathrm{Precision}=0{,}76$ y $\mathrm{Recall}=0{,}80$, la Ecuación~\eqref{eq:binary-f1-example} muestra el resultado:
\begin{equation}
\label{eq:binary-f1-example}
F_1=\frac{2\cdot 0{,}76\cdot 0{,}80}{0{,}76+0{,}80}\approx 0{,}78.
\end{equation}

\paragraph{Specificity (TNR), FPR y FNR}

Estas tres métricas son ``primas'' del recall, vistas desde la otra clase. Las Ecuaciones~\eqref{eq:binary-specificity-fpr} y~\eqref{eq:binary-fnr} las definen según las convenciones habituales de matriz de confusión \cite{sokolova_lapalme_2009}.

\textbf{Fórmulas:}
\begin{equation}
\label{eq:binary-specificity-fpr}
\mathrm{Specificity}=\mathrm{TNR}=\frac{TN}{TN+FP},\qquad
\mathrm{FPR}=\frac{FP}{FP+TN}=1-\mathrm{Specificity},
\end{equation}
\begin{equation}
\label{eq:binary-fnr}
\mathrm{FNR}=\frac{FN}{FN+TP}=1-\mathrm{Recall}.
\end{equation}

\textbf{Utilidad.}
\begin{itemize}[leftmargin=*]
  \item \textbf{Specificity (TNR):} de todas las requests normales, qué
  proporción se clasificó bien como normal. Es el ``recall de la clase
  negativa''.
  \item \textbf{FPR:} qué porcentaje de las requests normales generó una
  falsa alarma.
  \item \textbf{FNR:} qué porcentaje de los ataques pasó desapercibido.
\end{itemize}

\textbf{Ejemplo.} Con $TN=75$, $FP=5$, $FN=4$, $TP=16$, la Ecuación~\eqref{eq:binary-specificity-fpr-fnr-example} muestra el cálculo conjunto:
\begin{equation}
\label{eq:binary-specificity-fpr-fnr-example}
\mathrm{Specificity}=\frac{75}{80}=0{,}9375,\qquad
\mathrm{FPR}=\frac{5}{80}=0{,}0625,\qquad
\mathrm{FNR}=\frac{4}{20}=0{,}20.
\end{equation}

\paragraph{Balanced Accuracy}

\textbf{Fórmula.} La Ecuación~\eqref{eq:binary-balanced-accuracy} define la balanced accuracy \cite{brodersen2010}:
\begin{equation}
\label{eq:binary-balanced-accuracy}
\mathrm{Balanced\ Accuracy}=\frac{1}{2}\left(\mathrm{TPR}+\mathrm{TNR}\right)
=\frac{1}{2}\left(\frac{TP}{TP+FN}+\frac{TN}{TN+FP}\right).
\end{equation}

\textbf{Utilidad.} Promedia el recall de cada clase. Corrige el problema
de la accuracy en datasets desbalanceados: aunque haya muchos más
normales que ataques, ambas clases pesan lo mismo.

\textbf{Ejemplo.} Con $\mathrm{TPR}=0{,}80$ y $\mathrm{TNR}=0{,}9375$, la Ecuación~\eqref{eq:binary-balanced-accuracy-example} muestra el cálculo:
\begin{equation}
\label{eq:binary-balanced-accuracy-example}
\mathrm{Balanced\ Accuracy}=\frac{0{,}80+0{,}9375}{2}\approx 0{,}87.
\end{equation}
Es algo menor que la accuracy ordinaria (0{,}91): refleja que la clase
minoritaria (ataques) se detecta peor que la mayoritaria.

\paragraph{Matthews Correlation Coefficient (MCC)}

\textbf{Fórmula.} La Ecuación~\eqref{eq:binary-mcc} define el MCC, introducido originalmente como coeficiente de correlación para comparar predicción y observación \cite{matthews1975}:
\begin{equation}
\label{eq:binary-mcc}
\mathrm{MCC}=\frac{TP\cdot TN - FP\cdot FN}{\sqrt{(TP+FP)(TP+FN)(TN+FP)(TN+FN)}}.
\end{equation}

\textbf{Utilidad.} El MCC condensa las cuatro celdas de la matriz de
confusión en un único número en el rango $[-1,+1]$:
\begin{itemize}[leftmargin=*]
  \item $+1$: predicción perfecta.
  \item $0$: el modelo no es mejor que adivinar al azar.
  \item $-1$: predicción inversa (siempre se equivoca).
\end{itemize}
Se considera más robusto que accuracy o $F_1$ en escenarios
desbalanceados, porque solo toma valores altos cuando el modelo acierta
bien en \emph{ambas} clases \cite{chicco_mcc}.

\textbf{Ejemplo.} La Ecuación~\eqref{eq:binary-mcc-example} aplica la fórmula del MCC al ejemplo de referencia.
\begin{equation}
\label{eq:binary-mcc-example}
\mathrm{MCC}=\frac{16\cdot 75 - 5\cdot 4}{\sqrt{21\cdot 20\cdot 80\cdot 79}}
=\frac{1180}{\sqrt{2{,}6544\times 10^{6}}}\approx 0{,}72.
\end{equation}

\paragraph{ROC-AUC y PR-AUC}

Las curvas ROC se usan para visualizar y analizar clasificadores a través de distintos umbrales \cite{fawcett2006}; las curvas Precision--Recall son especialmente informativas cuando la clase positiva es minoritaria \cite{saito_pr_auc}. Cuando el clasificador no entrega solo una etiqueta sino un \emph{score
continuo} (típicamente una probabilidad entre 0 y 1), hace falta un
\emph{umbral} para decidir si una request es ataque. Cada umbral produce
una matriz de confusión distinta y, por lo tanto, un par diferente de
$(\mathrm{Precision},\mathrm{Recall})$ o de $(\mathrm{FPR},\mathrm{TPR})$.
Las curvas ROC y PR resumen el comportamiento del modelo a lo largo de
\emph{todos} los umbrales posibles, y sus áreas bajo la curva (AUC) lo
sintetizan en un solo número.

\textbf{ROC-AUC.} La curva ROC grafica TPR (eje $y$) frente a FPR (eje
$x$). El ROC-AUC vive en $[0,1]$ y se interpreta como la probabilidad de
que el modelo asigne un score más alto a un ataque elegido al azar que a
un normal elegido al azar. $0{,}5$ corresponde a un modelo aleatorio;
$1{,}0$ a uno perfecto.

\textbf{PR-AUC.} La curva PR grafica Precision (eje $y$) frente a Recall
(eje $x$). Cuando la clase positiva es minoritaria, PR-AUC suele ser más
informativa que ROC-AUC, porque la curva ROC puede verse demasiado
optimista al estar dominada por el gran número de verdaderos negativos
\cite{saito_pr_auc}.

\textbf{Ejemplo} Si el modelo asigna scores
$\{0{,}9;\ 0{,}8;\ 0{,}3\}$ a tres ataques y
$\{0{,}7;\ 0{,}2;\ 0{,}1\}$ a tres normales, los ataques superan a los
normales en 8 de las 9 comparaciones posibles, así que
$\mathrm{ROC\text{-}AUC}=8/9\approx 0{,}89$.

\subsubsection{Clasificación multiclase}

En multiclase hay $K$ clases posibles (por ejemplo, varias categorías de
ataque más la clase ``normal'') y cada instancia pertenece a
\emph{exactamente una}. Para reutilizar las métricas binarias se usa el
esquema \emph{one-vs-rest}: para cada clase $k$ se considera ``positiva''
la clase $k$ y ``negativas'' todas las demás, obteniéndose $TP_k$,
$FP_k$ y $FN_k$ por clase.

\paragraph{Ejemplo de referencia.} Trabajaremos con tres clases ($A$,
$B$, $C$) sobre 100 instancias, con la matriz de confusión de la Ecuación~\eqref{eq:multiclass-confusion-example}:
\begin{equation}
\label{eq:multiclass-confusion-example}
\begin{array}{l|ccc|c}
 & \text{Pred }A & \text{Pred }B & \text{Pred }C & \text{Total} \\
\hline
\text{True }A & 85 & 3 & 2 & 90 \\
\text{True }B & 2  & 4 & 1 & 7  \\
\text{True }C & 1  & 1 & 1 & 3  \\
\hline
\text{Total} & 88 & 8 & 4 & 100
\end{array}
\end{equation}
La clase $A$ es mayoritaria (soporte 90), $B$ y $C$ son minoritarias.
A partir de esta matriz se obtienen los conteos por clase de la Ecuación~\eqref{eq:multiclass-counts-example}, que
reutilizamos en los ejemplos:
\begin{equation}
\label{eq:multiclass-counts-example}
\begin{array}{l|ccc}
 & TP_k & FP_k & FN_k \\
\hline
A & 85 & 3 & 5 \\
B & 4  & 4 & 3 \\
C & 1  & 3 & 2
\end{array}
\end{equation}

\paragraph{Accuracy multiclase}

\textbf{Fórmula.} La Ecuación~\eqref{eq:multiclass-accuracy} define la accuracy multiclase \cite{sokolova_lapalme_2009}:
\begin{equation}
\label{eq:multiclass-accuracy}
\mathrm{Accuracy}=\frac{1}{n}\sum_{i=1}^{n}\ind[\hat y_i=y_i].
\end{equation}

\textbf{Utilidad.} Proporción de instancias correctamente clasificadas.
Sigue siendo intuitiva pero hereda el problema del desbalance: si una
clase domina, puede ocultar el desempeño en clases minoritarias.

\textbf{Ejemplo.} En la diagonal de la matriz hay $85+4+1=90$ aciertos
sobre 100, así que $\mathrm{Accuracy}=0{,}90$.

\paragraph{Balanced Accuracy multiclase}

\textbf{Fórmula.} La Ecuación~\eqref{eq:multiclass-balanced-accuracy} define la balanced accuracy multiclase \cite{brodersen2010,sokolova_lapalme_2009}:
\begin{equation}
\label{eq:multiclass-balanced-accuracy}
\mathrm{Balanced\ Accuracy}=\frac{1}{K}\sum_{k=1}^{K}\frac{TP_k}{TP_k+FN_k}.
\end{equation}

\textbf{Utilidad.} Promedio del recall por clase: cada clase contribuye
igual sin importar cuántos ejemplos tenga. Es útil cuando las clases
minoritarias son tan importantes como las mayoritarias.

\textbf{Ejemplo.} Los recalls por clase son $85/90\approx 0{,}944$,
$4/7\approx 0{,}571$ y $1/3\approx 0{,}333$; el cálculo agregado se muestra en la Ecuación~\eqref{eq:multiclass-balanced-accuracy-example}:
\begin{equation}
\label{eq:multiclass-balanced-accuracy-example}
\mathrm{Balanced\ Accuracy}=\frac{0{,}944+0{,}571+0{,}333}{3}\approx 0{,}616.
\end{equation}
Mucho menor que la accuracy ordinaria (0{,}90), porque las clases
minoritarias se detectan bastante peor.

\paragraph{Promedios micro, macro y weighted}

Cuando hay más de dos clases hace falta agregar las métricas por clase
en un solo número. Existen tres formas estándar; cada una responde a una
pregunta diferente y se usa habitualmente en evaluaciones multiclase \cite{sokolova_lapalme_2009}. Primero se definen las métricas por clase. Para cada clase $k$, tratada como ``clase positiva'' frente al resto, las Ecuaciones~\eqref{eq:multiclass-per-class-formulas} y~\eqref{eq:multiclass-per-class-f1-formula} definen precisión, exhaustividad y $F_1$ por clase:
\begin{equation}
\label{eq:multiclass-per-class-formulas}
\mathrm{Precision}_k=\frac{TP_k}{TP_k+FP_k},\qquad
\mathrm{Recall}_k=\frac{TP_k}{TP_k+FN_k},
\end{equation}
\begin{equation}
\label{eq:multiclass-per-class-f1-formula}
F_{1,k}=\frac{2\cdot \mathrm{Precision}_k\cdot \mathrm{Recall}_k}{\mathrm{Precision}_k+\mathrm{Recall}_k}.
\end{equation}
La precisión por clase responde cuántas predicciones de la clase $k$ fueron correctas; la exhaustividad por clase responde cuántos ejemplos reales de la clase $k$ fueron recuperados; y $F_{1,k}$ resume ambas. A partir del ejemplo de referencia, la Ecuación~\eqref{eq:multiclass-per-class-example} muestra los valores por clase:
\begin{equation}
\label{eq:multiclass-per-class-example}
\begin{array}{l|ccc}
 & \mathrm{Precision}_k & \mathrm{Recall}_k & F_{1,k} \\
\hline
A & 85/88\approx 0{,}966 & 85/90\approx 0{,}944 & 0{,}955 \\
B & 4/8 = 0{,}500        & 4/7 \approx 0{,}571  & 0{,}533 \\
C & 1/4 = 0{,}250        & 1/3 \approx 0{,}333  & 0{,}286
\end{array}
\end{equation}

\textbf{Micro.} Se suman primero los conteos de todas las clases y luego
se calcula la métrica, como se define en la Ecuación~\eqref{eq:multiclass-micro}:
\begin{equation}
\label{eq:multiclass-micro}
\mathrm{Precision}_{micro}=\frac{\sum_{k=1}^{K} TP_k}{\sum_{k=1}^{K}(TP_k+FP_k)},\qquad
\mathrm{Recall}_{micro}=\frac{\sum_{k=1}^{K} TP_k}{\sum_{k=1}^{K}(TP_k+FN_k)},
\end{equation}
\begin{equation}
\label{eq:multiclass-f1-micro}
F_{1,micro}=\frac{2\cdot \mathrm{Precision}_{micro}\cdot \mathrm{Recall}_{micro}}{\mathrm{Precision}_{micro}+\mathrm{Recall}_{micro}}.
\end{equation}
Cada \emph{instancia} pesa igual. Esta agregación es útil para resumir el rendimiento global cuando se quiere que las clases con más ejemplos tengan más influencia, pero puede ocultar un mal desempeño en clases raras.

\emph{Ejemplo.} Con $\sum_k TP_k=90$, $\sum_k FP_k=10$ y $\sum_k FN_k=10$, la Ecuación~\eqref{eq:multiclass-micro-example} muestra el cálculo:
\begin{equation}
\label{eq:multiclass-micro-example}
\mathrm{Precision}_{micro}=\frac{90}{100}=0{,}90,\quad
\mathrm{Recall}_{micro}=\frac{90}{100}=0{,}90,\quad
F_{1,micro}=0{,}90.
\end{equation}
En problemas multiclase de etiqueta única las tres métricas micro coinciden con la exactitud (accuracy), como ocurre aquí ($0{,}90$).

\textbf{Macro.} Se calcula la métrica por clase y se promedia sin
ponderar, como indica la Ecuación~\eqref{eq:multiclass-macro}:
\begin{equation}
\label{eq:multiclass-macro}
\mathrm{Precision}_{macro}=\frac{1}{K}\sum_{k=1}^{K}\mathrm{Precision}_k,\qquad
\mathrm{Recall}_{macro}=\frac{1}{K}\sum_{k=1}^{K}\mathrm{Recall}_k,\qquad
F_{1,macro}=\frac{1}{K}\sum_{k=1}^{K}F_{1,k}.
\end{equation}
Cada \emph{clase} pesa igual. Esta versión es especialmente útil en seguridad cuando las clases minoritarias representan ataques importantes: si una clase rara se clasifica mal, la métrica macro lo refleja con claridad.

\emph{Ejemplo.} Las Ecuaciones~\eqref{eq:multiclass-precision-recall-macro-example} y~\eqref{eq:multiclass-f1-macro-example} calculan las versiones macro del ejemplo.
\begin{equation}
\label{eq:multiclass-precision-recall-macro-example}
\mathrm{Precision}_{macro}=\frac{0{,}966+0{,}500+0{,}250}{3}\approx 0{,}572,\qquad
\mathrm{Recall}_{macro}=\frac{0{,}944+0{,}571+0{,}333}{3}\approx 0{,}616,
\end{equation}
\begin{equation}
\label{eq:multiclass-f1-macro-example}
F_{1,macro}=\frac{0{,}955+0{,}533+0{,}286}{3}\approx 0{,}591.
\end{equation}

\textbf{Ponderado por soporte (weighted).} Se calcula la métrica por clase y se promedia ponderando por el soporte $s_k$ (cantidad de ejemplos verdaderos de la clase $k$, con $\sum_k s_k = n$), según la Ecuación~\eqref{eq:multiclass-weighted}:
\begin{equation}
\label{eq:multiclass-weighted}
\mathrm{Precision}_{weighted}=\sum_{k=1}^{K}\frac{s_k}{n}\mathrm{Precision}_k,\qquad
\mathrm{Recall}_{weighted}=\sum_{k=1}^{K}\frac{s_k}{n}\mathrm{Recall}_k,\qquad
F_{1,weighted}=\sum_{k=1}^{K}\frac{s_k}{n}F_{1,k}.
\end{equation}
Las clases mayoritarias pesan más. Es útil cuando se busca un valor representativo del comportamiento global respetando el tamaño real de cada clase, pero puede suavizar problemas en clases poco frecuentes.

\emph{Ejemplo.} Con soportes $s_A=90$, $s_B=7$, $s_C=3$, las Ecuaciones~\eqref{eq:multiclass-precision-recall-weighted-example} y~\eqref{eq:multiclass-f1-weighted-example} muestran el cálculo:
\begin{equation}
\label{eq:multiclass-precision-recall-weighted-example}
\mathrm{Precision}_{weighted}=0{,}90\cdot0{,}966+0{,}07\cdot0{,}500+0{,}03\cdot0{,}250\approx0{,}912,
\end{equation}
\begin{equation}
\label{eq:multiclass-f1-weighted-example}
F_{1,weighted}=0{,}90\cdot 0{,}955+0{,}07\cdot 0{,}533+0{,}03\cdot 0{,}286\approx 0{,}906.
\end{equation}

\textbf{Comparación.} Sobre el mismo modelo, la Ecuación~\eqref{eq:multiclass-f1-comparison-example} resume las tres agregaciones:
\begin{equation}
\label{eq:multiclass-f1-comparison-example}
F_{1,micro}=0{,}90,\qquad F_{1,macro}\approx 0{,}59,\qquad F_{1,weighted}\approx 0{,}91.
\end{equation}
Las versiones micro y ponderada dan valores parecidos y altos porque la clase mayoritaria ($A$) domina el cálculo. La macro, en cambio, revela que el modelo trata mal a las clases minoritarias: ese 0{,}59 es la señal honesta de que $B$ y $C$ están lejos de funcionar bien.

\paragraph{MCC multiclase}

El coeficiente de correlación de Matthews también puede calcularse para $K$ clases usando la matriz de confusión completa. Sea $C\in\mathbb{N}^{K\times K}$ una matriz donde $C_{ij}$ indica cuántas instancias de clase real $i$ fueron predichas como clase $j$. Definimos:
\begin{equation}
\label{eq:multiclass-mcc-defs}
c=\sum_{k=1}^{K} C_{kk},\qquad
s=\sum_{i=1}^{K}\sum_{j=1}^{K} C_{ij},\qquad
p_k=\sum_{i=1}^{K} C_{ik},\qquad
t_k=\sum_{j=1}^{K} C_{kj},
\end{equation}
donde $c$ es la suma de la diagonal, $s$ es el total de instancias, $p_k$ es el total predicho como clase $k$ y $t_k$ es el soporte real de la clase $k$. La generalización multiclase se define en la Ecuación~\eqref{eq:multiclass-mcc} \cite{gorodkin2004,chicco_mcc}:
\begin{equation}
\label{eq:multiclass-mcc}
\mathrm{MCC}_{multi}=\frac{c\,s-\sum_{k=1}^{K} p_k t_k}{\sqrt{\left(s^2-\sum_{k=1}^{K}p_k^2\right)\left(s^2-\sum_{k=1}^{K}t_k^2\right)}}.
\end{equation}

La interpretación se mantiene igual que en binario: $+1$ indica predicción perfecta, $0$ indica comportamiento equivalente al azar y valores negativos indican desacuerdo sistemático. Su utilidad en este trabajo es que resume la matriz completa sin quedar dominado únicamente por la clase mayoritaria.

\textbf{Ejemplo.} En la matriz de la Ecuación~\eqref{eq:multiclass-confusion-example}, $c=90$, $s=100$, los totales predichos son $p_A=88$, $p_B=8$, $p_C=4$, y los soportes reales son $t_A=90$, $t_B=7$, $t_C=3$. La Ecuación~\eqref{eq:multiclass-mcc-example} muestra el cálculo:
\begin{equation}
\label{eq:multiclass-mcc-example}
\mathrm{MCC}_{multi}=\frac{90\cdot100-(88\cdot90+8\cdot7+4\cdot3)}{\sqrt{(100^2-88^2-8^2-4^2)(100^2-90^2-7^2-3^2)}}\approx 0{,}506.
\end{equation}
Este valor es bastante menor que la exactitud de $0{,}90$ porque el modelo acierta mucho en la clase mayoritaria, pero falla de forma importante en las clases minoritarias.

\paragraph{ROC-AUC y PR-AUC multiclase}

Cuando $K>2$ y el modelo entrega probabilidades, ROC-AUC y PR-AUC se
calculan sobre una binarización \emph{one-vs-rest}: para cada clase $k$
se construye una curva tratando ``$k$ vs.\ todas las demás'' y se
obtiene un AUC por clase. Luego se agregan: las versiones \emph{macro}
promedian sin ponderar, las \emph{weighted} promedian ponderando por el
soporte $s_k$.

\textbf{Ejemplo.} Supongamos que, sobre el ejemplo de referencia, las
ROC-AUC one-vs-rest por clase son $\mathrm{AUC}_A=0{,}97$,
$\mathrm{AUC}_B=0{,}80$, $\mathrm{AUC}_C=0{,}60$. Entonces las Ecuaciones~\eqref{eq:multiclass-roc-auc-macro-example} y~\eqref{eq:multiclass-roc-auc-weighted-example} muestran las agregaciones macro y ponderada:
\begin{equation}
\label{eq:multiclass-roc-auc-macro-example}
\mathrm{ROC\text{-}AUC}_{macro}=\frac{0{,}97+0{,}80+0{,}60}{3}\approx 0{,}79,
\end{equation}
\begin{equation}
\label{eq:multiclass-roc-auc-weighted-example}
\mathrm{ROC\text{-}AUC}_{weighted}=0{,}90\cdot 0{,}97+0{,}07\cdot 0{,}80+0{,}03\cdot 0{,}60\approx 0{,}94.
\end{equation}
La versión macro deja ver que la clase $C$ se ordena casi al azar; la
weighted lo enmascara al estar dominada por $A$.

\subsubsection{Clasificación multietiqueta}

A diferencia de multiclase, donde cada instancia pertenece a
\emph{exactamente una} clase, en multietiqueta cada instancia puede
tener \emph{varias etiquetas a la vez} o ninguna, formulación ampliamente tratada en la literatura multietiqueta \cite{tsoumakas_multilabel_overview,read2011}. Por ejemplo, una
request maliciosa puede combinar varias técnicas de ataque (varios
CAPECs simultáneos).

Formalmente, $Y\in\{0,1\}^{n\times L}$ y $\hat{Y}\in\{0,1\}^{n\times L}$,
donde $L$ es la cantidad total de etiquetas posibles y $y_{i\ell}=1$
indica que la instancia $i$ tiene la etiqueta $\ell$.

\paragraph{Ejemplo de referencia.} Usaremos $n=3$ instancias y $L=3$
etiquetas ($L_1$, $L_2$, $L_3$), con la verdad y predicción de la Ecuación~\eqref{eq:multilabel-truth-pred-example}:
\begin{equation}
\label{eq:multilabel-truth-pred-example}
\begin{array}{c|ccc|ccc}
 & \multicolumn{3}{c|}{Y_i\ \text{(verdad)}} & \multicolumn{3}{c}{\hat Y_i\ \text{(predicción)}} \\
i & L_1 & L_2 & L_3 & L_1 & L_2 & L_3 \\
\hline
1 & 1 & 1 & 0 & 1 & 1 & 0 \\
2 & 0 & 1 & 1 & 1 & 1 & 0 \\
3 & 1 & 0 & 0 & 1 & 0 & 1
\end{array}
\end{equation}
Hay $nL=9$ predicciones etiqueta-instancia en total. La Ecuación~\eqref{eq:multilabel-counts-example} resume los conteos por etiqueta:
\begin{equation}
\label{eq:multilabel-counts-example}
\begin{array}{l|ccc}
 & TP_\ell & FP_\ell & FN_\ell \\
\hline
L_1 & 2 & 1 & 0 \\
L_2 & 2 & 0 & 0 \\
L_3 & 0 & 1 & 1
\end{array}
\end{equation}

\paragraph{Hamming Loss}

\textbf{Fórmula.} La Ecuación~\eqref{eq:multilabel-hamming-loss} define Hamming Loss, métrica estándar de evaluación multietiqueta \cite{tsoumakas_multilabel_overview,read2011}:
\begin{equation}
\label{eq:multilabel-hamming-loss}
\mathrm{HL}=\frac{1}{nL}\sum_{i=1}^n\sum_{\ell=1}^L \ind\big[\hat{y}_{i\ell}\ne y_{i\ell}\big].
\end{equation}

\textbf{Utilidad.} Mide el error promedio por predicción
etiqueta-por-instancia: equivale a la fracción de celdas equivocadas en
la matriz $\hat{Y}$ comparada con $Y$. Cuanto menor, mejor;
$\mathrm{HL}=0$ es predicción perfecta. Su rango es $0 \leq \mathrm{HL} \leq 1$; valores cercanos a 1 indican una proporción alta de predicciones etiqueta-petición incorrectas.

\textbf{Ejemplo.} En el ejemplo de referencia las celdas mal predichas
son tres: $(i{=}2,L_1)$, $(i{=}2,L_3)$ y $(i{=}3,L_3)$. Entonces la Ecuación~\eqref{eq:multilabel-hamming-loss-example} muestra el cálculo:
\begin{equation}
\label{eq:multilabel-hamming-loss-example}
\mathrm{HL}=\frac{3}{9}\approx 0{,}33.
\end{equation}

\paragraph{Precision, Recall y $F_1$ micro/macro multietiqueta}

Las definiciones por etiqueta son las mismas que en multiclase, pero
ahora cada etiqueta $\ell$ aporta sus propios $TP_\ell$, $FP_\ell$,
$FN_\ell$ basándose en las celdas correspondientes de $Y$ y $\hat Y$.
\begin{itemize}[leftmargin=*]
  \item \textbf{Micro:} agrega todas las predicciones positivas de todas
  las etiquetas antes de calcular la métrica. Cada predicción
  etiqueta-instancia pesa igual.
  \item \textbf{Macro:} calcula la métrica por etiqueta y luego promedia
  sin ponderar. Cada etiqueta pesa igual, lo que penaliza el mal
  desempeño en etiquetas raras.
\end{itemize}

\textbf{Ejemplo.} La Ecuación~\eqref{eq:multilabel-per-label-example} muestra primero los valores por etiqueta:
\begin{equation}
\label{eq:multilabel-per-label-example}
\begin{array}{l|ccc}
 & \mathrm{Precision}_\ell & \mathrm{Recall}_\ell & F_{1,\ell} \\
\hline
L_1 & 2/3\approx 0{,}667 & 2/2 = 1{,}000 & 0{,}800 \\
L_2 & 2/2 = 1{,}000      & 2/2 = 1{,}000 & 1{,}000 \\
L_3 & 0/1 = 0            & 0/1 = 0       & 0
\end{array}
\end{equation}

\emph{Macro.} La Ecuación~\eqref{eq:multilabel-f1-macro-example} muestra el promedio macro.
\begin{equation}
\label{eq:multilabel-f1-macro-example}
F_{1,macro}=\frac{0{,}800+1{,}000+0}{3}=0{,}600.
\end{equation}

\emph{Micro.} Con $\sum_\ell TP_\ell=4$, $\sum_\ell FP_\ell=2$ y
$\sum_\ell FN_\ell=1$, la Ecuación~\eqref{eq:multilabel-micro-example} muestra el cálculo micro:
\begin{equation}
\label{eq:multilabel-micro-example}
\mathrm{Precision}_{micro}=\frac{4}{4+2}\approx 0{,}667,\quad
\mathrm{Recall}_{micro}=\frac{4}{4+1}=0{,}800,\quad
F_{1,micro}\approx 0{,}727.
\end{equation}

La macro castiga fuerte el desastre en $L_3$ (que vale tanto como las
otras dos etiquetas), mientras que la micro lo diluye porque $L_3$
aporta pocas predicciones positivas en términos absolutos.

\paragraph{Jaccard Similarity}

\textbf{Fórmula por instancia.} La Ecuación~\eqref{eq:jaccard-instance} define Jaccard por instancia, también usada como medida de similitud o accuracy multietiqueta basada en conjuntos \cite{tsoumakas_multilabel_overview,read2011}:
\begin{equation}
\label{eq:jaccard-instance}
J_i = \frac{|Y_i \cap \hat{Y}_i|}{|Y_i \cup \hat{Y}_i|},
\end{equation}
donde $Y_i$ y $\hat{Y}_i$ son los conjuntos de etiquetas verdaderas y
predichas de la instancia $i$.

\textbf{Utilidad.} Mide cuánto se parecen los dos conjuntos de
etiquetas: $J_i=1$ si coinciden exactamente, $J_i=0$ si son disjuntos.
Es una forma natural de comparar predicciones con etiquetas múltiples,
porque considera simultáneamente aciertos, omisiones y excesos.

\textbf{Aclaración sobre el caso binario.} En este trabajo Jaccard se
interpreta principalmente en la formulación multietiqueta, donde cada
petición puede tener varias etiquetas CAPEC activas. Si una tarea queda
reducida a binaria, el Jaccard se puede calcular algebraicamente como
$TP/(TP+FP+FN)$, que es la forma binaria equivalente derivada de la Ecuación~\eqref{eq:jaccard-instance}; sin embargo, esa lectura aporta menos información que las
métricas binarias principales —recall/TPR, specificity/TNR, FPR, F1,
PR-AUC y MCC— porque ya no se comparan conjuntos ricos de etiquetas,
sino una sola decisión positiva/negativa. Por eso, cuando un script
lo exporta en un modo binario reducido, se usa solo como indicador
secundario de consistencia y no como métrica central de la detección
binaria.

Existen tres variantes según cómo se agregan los $J_i$:
\begin{itemize}[leftmargin=*]
  \item \textbf{Jaccard samples:} promedia $J_i$ sobre las instancias.
  \item \textbf{Jaccard micro:} agrega todas las celdas positivas y
  calcula un único cociente global, $J_{micro}=\sum_\ell TP_\ell\,/\,\sum_\ell(TP_\ell+FP_\ell+FN_\ell)$.
  \item \textbf{Jaccard macro:} calcula un Jaccard por etiqueta,
  $J_\ell=TP_\ell/(TP_\ell+FP_\ell+FN_\ell)$, y luego promedia.
\end{itemize}

\textbf{Ejemplo.} Sobre el ejemplo de referencia.

\emph{Por instancia.} Las Ecuaciones~\eqref{eq:jaccard-j1-example} y~\eqref{eq:jaccard-j2-j3-example} muestran los cálculos por instancia.
\begin{equation}
\label{eq:jaccard-j1-example}
J_1=\tfrac{|\{L_1,L_2\}\cap\{L_1,L_2\}|}{|\{L_1,L_2\}\cup\{L_1,L_2\}|}=\tfrac{2}{2}=1,
\end{equation}
\begin{equation}
\label{eq:jaccard-j2-j3-example}
J_2=\tfrac{|\{L_2,L_3\}\cap\{L_1,L_2\}|}{|\{L_2,L_3\}\cup\{L_1,L_2\}|}=\tfrac{1}{3},\qquad
J_3=\tfrac{|\{L_1\}\cap\{L_1,L_3\}|}{|\{L_1\}\cup\{L_1,L_3\}|}=\tfrac{1}{2}.
\end{equation}

\emph{Samples (promedio sobre instancias).} La Ecuación~\eqref{eq:jaccard-samples-example} promedia los valores por instancia.
\begin{equation}
\label{eq:jaccard-samples-example}
J_{samples}=\frac{1+\tfrac{1}{3}+\tfrac{1}{2}}{3}=\frac{11}{18}\approx 0{,}611.
\end{equation}

\emph{Macro (promedio sobre etiquetas).} Con $J_{L_1}=2/3$,
$J_{L_2}=1$, $J_{L_3}=0$, la Ecuación~\eqref{eq:jaccard-macro-example} calcula el promedio macro:
\begin{equation}
\label{eq:jaccard-macro-example}
J_{macro}=\frac{2/3+1+0}{3}\approx 0{,}556.
\end{equation}

\emph{Micro (cociente global).} La Ecuación~\eqref{eq:jaccard-micro-example} muestra el cociente global.
\begin{equation}
\label{eq:jaccard-micro-example}
J_{micro}=\frac{4}{4+2+1}=\frac{4}{7}\approx 0{,}571.
\end{equation}

\paragraph{Exact Match Ratio}

\textbf{Fórmula.} La Ecuación~\eqref{eq:exact-match-ratio} define Exact Match Ratio, también llamado \emph{subset accuracy} en evaluación multietiqueta \cite{tsoumakas_multilabel_overview,read2011}:
\begin{equation}
\label{eq:exact-match-ratio}
\mathrm{EMR}=\frac{1}{n}\sum_{i=1}^n \ind[\hat{Y}_i = Y_i].
\end{equation}

\textbf{Utilidad.} También llamada \emph{subset accuracy}. Mide la
proporción de instancias en las que la predicción multilabel coincide
\emph{exactamente} con el conjunto verdadero. Es muy estricta: basta una
etiqueta de más o de menos para que la instancia cuente como error.

\textbf{Ejemplo.} En el ejemplo de referencia solo $i=1$ tiene
predicción idéntica a la verdad; $i=2$ e $i=3$ difieren en alguna
etiqueta. La Ecuación~\eqref{eq:exact-match-ratio-example} muestra el cálculo:
\begin{equation}
\label{eq:exact-match-ratio-example}
\mathrm{EMR}=\frac{1}{3}\approx 0{,}33.
\end{equation}
Notar que las instancias $i=2$ e $i=3$ aportan información parcial al
Jaccard (0{,}33 y 0{,}50) pero cero al EMR.

\paragraph{ROC-AUC y PR-AUC multietiqueta}

Cuando se dispone de scores continuos por etiqueta, se calculan las
versiones \emph{micro} y \emph{macro} de ROC-AUC y PR-AUC tratando cada
etiqueta como un problema binario independiente y luego agregando: la
versión micro fusiona los scores de todas las etiquetas en una única
curva global, mientras que la macro promedia las áreas calculadas por
etiqueta.

\textbf{Ejemplo.} Si las ROC-AUC por etiqueta fueran $\mathrm{AUC}_{L_1}=0{,}90$,
$\mathrm{AUC}_{L_2}=0{,}95$ y $\mathrm{AUC}_{L_3}=0{,}55$, la versión
macro sería la de la Ecuación~\eqref{eq:multilabel-roc-auc-macro-example}:
\begin{equation}
\label{eq:multilabel-roc-auc-macro-example}
\mathrm{ROC\text{-}AUC}_{macro}=\frac{0{,}90+0{,}95+0{,}55}{3}=0{,}80.
\end{equation}
La versión micro no se obtiene como promedio de estos tres números:
requiere reagrupar los scores de las tres etiquetas en una única curva
global y calcular el área de esa curva, así que en general
$\mathrm{ROC\text{-}AUC}_{micro}\ne \mathrm{ROC\text{-}AUC}_{macro}$.

\subsubsection{Métricas temporales y operativas}

Además de medir efectividad predictiva, el trabajo considera métricas operativas porque un modelo útil en un entorno tipo WAF no solo debe clasificar bien: también debe responder con bajo costo y con tiempos compatibles con la inspección de tráfico HTTP. Estas métricas se relacionan con la viabilidad práctica del flujo de procesamiento y se interpretan como mediciones de tiempo, capacidad de procesamiento, estabilidad y consumo de recursos. A diferencia de las métricas predictivas anteriores, no evalúan si la etiqueta predicha es correcta; evalúan cuánto cuesta aplicar el modelo y si la medición se completa de manera estable.

\textbf{Alcance de la medición.} Se distinguen tres momentos del flujo experimental. Primero, la búsqueda de hiperparámetros, cuando existe. Segundo, la construcción o entrenamiento final del modelo. Tercero, la aplicación del flujo ya entrenado sobre peticiones medidas en una prueba operativa o \emph{benchmark}. Esta última etapa no incluye entrenamiento ni búsqueda de hiperparámetros: mide la capacidad de procesamiento de peticiones por unidad de tiempo. Cuando la medición desagrega tiempos, la latencia total corresponde a extracción de características más inferencia; la latencia de inferencia corresponde solo al paso del vector de características por el modelo. Los nombres técnicos de campos, como \texttt{wall\_time\_seconds}, se retienen únicamente para mantener trazabilidad con los artefactos de resultados; la interpretación de la métrica se da aquí en términos conceptuales.

\paragraph{Tiempo de optimización de parámetros.}
La optimización de hiperparámetros consiste en probar configuraciones del modelo antes de elegir la variante final. Es importante medirla porque puede mejorar el rendimiento, pero también aumenta el costo experimental. Si una mejora pequeña requiere una búsqueda demasiado larga, esa configuración puede ser poco conveniente para un flujo reproducible o para reentrenamientos frecuentes. La Ecuación~\eqref{eq:topt} define este intervalo.
\begin{equation}
\label{eq:topt}
t_{opt}=t_{fin\_busqueda}-t_{inicio\_busqueda}.
\end{equation}

Cuando una corrida no realiza búsqueda de hiperparámetros, este tiempo no aplica. En ese caso no debe interpretarse como que optimizar cuesta cero, sino como que esa etapa no fue ejecutada en esa corrida.

\paragraph{Tiempo de construcción del modelo.}
Mide el tiempo usado por el ajuste final del modelo, según la Ecuación~\eqref{eq:tbuild}:
\begin{equation}
\label{eq:tbuild}
t_{build}=t_{fin\_ajuste}-t_{inicio\_ajuste}.
\end{equation}

Esta métrica permite comparar cuánto cuesta construir cada modelo una vez fijada su configuración. En las tablas de resultados corresponde al tiempo de entrenamiento registrado para la corrida final.

\paragraph{Tiempo de aplicación.}
Mide el tiempo total empleado por la prueba operativa para procesar un conjunto de peticiones con el modelo ya entrenado, según la Ecuación~\eqref{eq:tapply}:
\begin{equation}
\label{eq:tapply}
t_{apply}=t_{fin\_aplicacion}-t_{inicio\_aplicacion}.
\end{equation}

En los artefactos de medición este intervalo se registra como tiempo de reloj total, equivalente al campo \texttt{wall\_time\_seconds}. Este valor se interpreta junto con las latencias por petición, porque el tiempo total depende de cuántas peticiones fueron medidas y de si el flujo incluyó extracción de características, inferencia o ambas etapas.

\paragraph{Capacidad de procesamiento (throughput).}
Indica cuántas peticiones se procesan por segundo durante la prueba operativa, como define la Ecuación~\eqref{eq:throughput}:
\begin{equation}
\label{eq:throughput}
\mathrm{Throughput}=\frac{\mathrm{peticiones\ completadas}}{t_{apply}}.
\end{equation}

Un valor mayor sugiere mayor capacidad de procesamiento, aunque debe leerse junto con latencia y estabilidad. Por ejemplo, un modelo puede procesar muchas peticiones por segundo en promedio y aun así presentar una latencia P95 alta si algunas peticiones tardan mucho más que la mayoría.

\paragraph{Latencia.}
La latencia mide cuánto tarda una petición individual en atravesar el flujo medido. Se reportan percentiles como P50, P95 y P99. P50 resume el caso típico; P95 indica que el 95\% de las peticiones medidas terminó por debajo de ese tiempo; y P99 aproxima la cola más exigente de la distribución. En un escenario tipo WAF, P95 y P99 son más útiles que el promedio para detectar demoras que podrían afectar la experiencia de uso o la capacidad de inspección en línea.

\paragraph{Estabilidad de la prueba operativa.}
La tasa de solicitudes fallidas durante la prueba operativa no describe errores de clasificación ni un ``fallo'' lógico del modelo. Describe peticiones que no pudieron completarse correctamente durante la medición, por ejemplo por errores de ejecución, excepción del flujo, entrada inválida no controlada o interrupción de la medición. La Ecuación~\eqref{eq:benchmark-failure-rate} formaliza esta tasa:
\begin{equation}
\label{eq:benchmark-failure-rate}
\mathrm{Tasa\ de\ fallos}=\frac{\mathrm{peticiones\ fallidas}}{\mathrm{peticiones\ medidas}}.
\end{equation}

Por tanto, esta métrica se interpreta como estabilidad de la corrida de medición. No implica que el modelo pueda quedar en un ciclo infinito ni que una predicción incorrecta cuente como fallo operativo.

\paragraph{Uso de CPU.}
El uso aproximado de CPU puede estimarse con la Ecuación~\eqref{eq:cpu-utilization}:
\begin{equation}
\label{eq:cpu-utilization}
\mathrm{CPU\ util\ \%}\approx 100\cdot \frac{\Delta(\mathrm{user}+\mathrm{system})}{t_{apply}},
\end{equation}
donde $\Delta(\mathrm{user}+\mathrm{system})$ representa el aumento del tiempo de CPU consumido en modo usuario y sistema durante la ventana de medición. El operador $\Delta$ significa ``cambio entre el valor final y el valor inicial''; por ejemplo, si al inicio se registran 12 segundos de CPU acumulada y al final 18 segundos, entonces $\Delta=6$ segundos de CPU.

\paragraph{Memoria y tamaño del modelo.}
La memoria se evalúa comparando el consumo de memoria residente antes y después de la medición, y registrando el pico aproximado cuando el sistema de medición lo permite. La Ecuación~\eqref{eq:memory-delta} define el incremento observado:
\begin{equation}
\label{eq:memory-delta}
\Delta \mathrm{RSS}=\mathrm{RSS}_{despues}-\mathrm{RSS}_{antes}.
\end{equation}
El pico aproximado, $\mathrm{RSS}_{pico}$, resume la mayor memoria residente observada durante la prueba. El tamaño del modelo se calcula como el tamaño en bytes del artefacto serializado, según la Ecuación~\eqref{eq:model-size}:
\begin{equation}
\label{eq:model-size}
\text{Tamaño del modelo}=\text{bytes}(\text{artefacto serializado}).
\end{equation}
Estas métricas permiten evaluar si una solución con buen rendimiento predictivo también es razonable desde el punto de vista operativo: un modelo puede clasificar bien, pero ser poco conveniente si requiere demasiada memoria, tarda demasiado en aplicarse o genera una prueba inestable.

\subsubsection{Diccionario resumido de métricas retenidas}

\begin{longtable}{@{}L{0.24\textwidth}L{0.43\textwidth}L{0.25\textwidth}@{}}
\caption{Diccionario resumido de métricas utilizadas en la evaluación.}\label{tab:diccionario-metricas}\\
\toprule
\textbf{Métrica} & \textbf{Qué mide} & \textbf{Uso principal} \\
\midrule
\endfirsthead
\toprule
\textbf{Métrica} & \textbf{Qué mide} & \textbf{Uso principal} \\
\midrule
\endhead
Exactitud (accuracy) & Proporción de predicciones correctas: $\frac{1}{n}\sum_i \ind[\hat y_i=y_i]$. & Multiclase y casos reducidos a dos clases. \\
Exactitud balanceada (balanced accuracy) & Promedio de la exhaustividad por clase; en binario equivale a $(\mathrm{TPR}+\mathrm{TNR})/2$. & Detección binaria y \textit{one-class}, útil con desbalance. \\
Precisión (precision) & $TP/(TP+FP)$ en binario; en agregaciones micro/macro/ponderadas sigue las fórmulas definidas para cada caso. & Calidad de predicciones positivas. \\
Exhaustividad / sensibilidad (recall / TPR) & $TP/(TP+FN)$. En binario coincide con la tasa de verdaderos positivos. & Capacidad de recuperación de ataques o clases reales. \\
$F_1$ & Media armónica entre precisión y exhaustividad. & Resumen sintético en detección binaria. \\
$F_1$ macro & Promedio simple del $F_1$ por clase o etiqueta. & Multiclase y multietiqueta, sensible a clases minoritarias. \\
$F_1$ ponderado / micro & En multiclase se retiene $F_1$ ponderado por soporte; en multietiqueta se retiene $F_1$ micro, agregado globalmente. & Comparación global de modelos supervisados. \\
MCC & Coeficiente de correlación de Matthews; en multiclase se usa la generalización de la Ecuación~\eqref{eq:multiclass-mcc}. & Resumen robusto con desbalance. \\
Hamming loss & Fracción de predicciones etiqueta-por-petición incorrectas. & Multietiqueta. \\
Jaccard micro & Intersección sobre unión calculada globalmente sobre todas las etiquetas; en binario reducido se lee solo como indicador secundario. & Multietiqueta. \\
Exact match ratio & Proporción de peticiones cuya salida multietiqueta coincide exactamente con la verdad. & Multietiqueta. \\
Tiempo de optimización & Tiempo total del bloque de búsqueda de hiperparámetros. & Ajuste de hiperparámetros supervisado y \textit{one-class}. \\
Tiempo de construcción & Tiempo del ajuste final del modelo.  & Todos los modelos. \\
Tiempo de aplicación & Tiempo total de la prueba operativa con el modelo ya entrenado. & Todos los modelos. \\
Capacidad de procesamiento (throughput) & Peticiones completadas por segundo: peticiones completadas/$t_{apply}$. & Viabilidad operativa. \\
Latencia P50/P95/P99 & Percentiles de tiempo por petición en el flujo medido. & Viabilidad operativa. \\
Tasa de fallos & Proporción de solicitudes no completadas durante la prueba operativa, no errores de clasificación. & Estabilidad operativa. \\
CPU/memoria/tamaño & Recursos consumidos durante la ejecución y tamaño en bytes del artefacto serializado. & Costo de despliegue. \\
\bottomrule
\end{longtable}


% ============================================================================
% CAPÍTULO 4
% ============================================================================


\chapter{Estado del arte y trabajos relacionados}

Este capítulo resume los trabajos relacionados que sirven como contexto para la comparación experimental. La revisión se organiza en torno a cuatro líneas: detección de anomalías en tráfico web, datasets de ataques HTTP, clasificación multietiqueta con taxonomías de ataque y evaluación operativa de modelos aplicables a un entorno tipo WAF.

\section{Detección de anomalías en tráfico HTTP}
Los primeros trabajos sobre detección de ataques web basada en anomalías mostraron que las peticiones HTTP contienen señales útiles para identificar desviaciones respecto del tráfico legítimo. Kruegel y Vigna propusieron analizar atributos estadísticos y estructurales de las peticiones para detectar ataques web sin depender exclusivamente de firmas predefinidas \cite{kruegel_vigna_2003}. En la misma línea, Ingham e Inoue compararon técnicas de detección de anomalías para HTTP y destacaron la importancia de modelar características específicas del protocolo \cite{ingham_inoue_2007}. También existen aproximaciones basadas en el contenido de la carga útil, como la detección de anomalías en payloads propuesta por Wang y Stolfo \cite{wang_stolfo_2004}.

Estos trabajos justifican el uso de variables interpretables extraídas de la URI, parámetros, cuerpo y estructura de la petición. Sin embargo, muchos enfoques clásicos se concentran en detección binaria o en escenarios controlados, por lo que no siempre permiten estudiar simultáneamente clasificación multiclase, clasificación multietiqueta y costo operativo.

\section{WAF basado en aprendizaje automático y métricas operativas}
El trabajo de Epp, Funk y Cappo propone un WAF basado en anomalías con características específicas de HTTP y One-Class SVM, incluyendo mediciones de efectividad y sobrecosto operacional \cite{ocswaf_zenodo}. Esa línea resulta cercana al objetivo de este estudio, porque combina detección de anomalías con una preocupación explícita por la viabilidad de integración en un entorno de inspección HTTP.

A diferencia de un despliegue completo de WAF, este trabajo conserva un alcance experimental: no implementa un proxy de filtrado ni reglas de bloqueo en producción. La comparación se enfoca en determinar si el flujo de extracción de características e inferencia puede sostener métricas razonables de latencia, rendimiento, memoria y tamaño de modelo.

\section{Datasets y clasificación de ataques web}
El dataset CSIC 2010 ha sido utilizado de forma recurrente como referencia para evaluar detección de ataques en aplicaciones web, aunque su carácter sintético limita la generalización directa a tráfico real \cite{csic2010_impact,csic2010_doc}. TorpEda 2012 aporta una especificación abierta de conjuntos de datos de ataques a aplicaciones web y permite trabajar con etiquetas de ataque más explícitas \cite{torpeda_2012}. Por su parte, SR-BH 2020 introduce tráfico real de honeypot con anotaciones CAPEC multietiqueta, lo que habilita una evaluación más rica desde el punto de vista semántico \cite{srbh_dataverse,srbh_paper}.

La literatura reciente también destaca que la selección de características puede capturar conocimiento experto de seguridad. Riverol et al. analizan el papel de la selección de variables en detección de ataques web, reforzando la necesidad de discutir qué señales usa realmente el modelo \cite{riverol_2024}. Esa preocupación se incorpora en este trabajo mediante una comparación con y sin tokens sospechosos.

\section{Modelos supervisados y multietiqueta}
La clasificación supervisada permite pasar de la pregunta ``normal o ataque'' a la identificación del tipo de ataque. La regresión logística se mantiene como una línea base robusta, eficiente e interpretable para problemas de clasificación \cite{islr_logreg,sklearn_logreg}. Para clasificación multietiqueta, la descomposición One-vs-Rest es una estrategia estándar que entrena un clasificador binario por etiqueta y permite conservar trazabilidad por categoría \cite{rifkin_ova,tsoumakas_multilabel_overview}.

En estudios más recientes de detección de intrusiones y ataques web se observan modelos de mayor complejidad, incluyendo enfoques generativos, ensambles y métodos de aprendizaje profundo \cite{fang_2024,zhou_2024}. Estos modelos son relevantes como contexto de literatura; sin embargo, el presente trabajo prioriza modelos lineales y one-class escalables para mantener interpretabilidad, reproducibilidad y costo controlado.

\section{Síntesis de la brecha abordada}
La revisión muestra tres brechas que motivan este trabajo. Primero, no basta con reportar exactitud predictiva: en un contexto tipo WAF también importan latencia, rendimiento, memoria y tamaño del modelo. Segundo, la comparación entre datasets sintéticos y tráfico real permite discutir el posible cambio de dominio. Tercero, la clasificación multietiqueta basada en CAPEC permite representar mejor peticiones asociadas a múltiples patrones de ataque, algo que se pierde cuando se fuerza una única clase por instancia. Por ello, este estudio compara de manera integrada enfoques one-class, multiclase y multietiqueta sobre una representación común de peticiones HTTP.

\chapter{Conjuntos de datos para experimentos}

Este capítulo describe las fuentes de datos utilizadas en los experimentos. La selección no busca reunir cualquier dataset de seguridad, sino conjuntos de datos que permitan trabajar con peticiones HTTP y comparar tres formulaciones del problema: detección binaria, clasificación multiclase y clasificación multietiqueta. Por esa razón, se distinguen tres criterios: pertinencia al dominio web, disponibilidad de etiquetas y posibilidad de transformar los campos de entrada a una representación común.

\section{SR-BH 2020: tráfico real de honeypot con etiquetas CAPEC}
SR-BH~2020 fue publicado en Harvard Dataverse y descrito en el artículo asociado como un dataset multietiqueta para clasificación de ataques web mediante CAPEC \cite{srbh_dataverse,srbh_paper}. Su interés principal para este trabajo es que no se limita a una etiqueta binaria de ataque, sino que permite representar una petición maliciosa con una o varias etiquetas CAPEC. Esa propiedad lo convierte en el dataset central para evaluar la formulación multietiqueta y para derivar, con una regla explícita, una versión multiclase basada en etiqueta primaria.

El dataset está alineado con un caso de uso tipo WAF porque el registro se realiza a nivel de petición HTTP y conserva campos directamente inspeccionables, como método, URI y cuerpo \cite{srbh_paper,owasp_waf}. Además, al provenir de un honeypot expuesto a Internet, aporta variabilidad de tráfico real y reduce el riesgo de evaluar únicamente sobre patrones sintéticos o demasiado controlados. Esta ventaja también introduce una limitación: el tráfico refleja el contexto específico del honeypot y no necesariamente todas las aplicaciones web de producción.

\subsection{SR-BH 2020 como referencia de tráfico real}
\begin{itemize}[leftmargin=*]
  \item \textbf{Tráfico real}: reduce el riesgo de sobre-ajuste a patrones sintéticos y captura variabilidad de producción \cite{srbh_paper}.
  \item \textbf{Etiquetas CAPEC multietiqueta}: permite ir más allá del binario y evaluar clasificación semántica por patrones \cite{srbh_paper,capec_about}.
  \item \textbf{Pertinencia a WAF}: modela directamente campos inspeccionados por un WAF, principalmente método, URI y cuerpo \cite{owasp_waf,modsecurity_site}.
  \item \textbf{Reproducibilidad}: al estar publicado con DOI persistente en Harvard Dataverse, facilita citación, trazabilidad y replicación de experimentos \cite{srbh_dataverse}.
\end{itemize}

\section{Datasets alternativos y unificación a un esquema común}
Además de SR-BH, se contemplan dos datasets alternativos para comparar desempeño y \textit{domain shift}. \textbf{CSIC~2010} es un conjunto ampliamente usado para evaluar detección de ataques web sobre tráfico HTTP normal y anómalo. Su aporte principal es la comparabilidad con trabajos previos, aunque su naturaleza sintética obliga a interpretarlo con cautela frente a tráfico real \cite{csic2010_impact,csic2010_doc}. \textbf{TorpEda~2012} propone una especificación abierta de conjuntos de datos de ataques a aplicaciones web y resulta útil como punto intermedio, porque permite trabajar con clases de ataque y no solo con una decisión normal/anómalo \cite{torpeda_2012}.

La inclusión de estos datasets permite evaluar si un modelo mantiene su comportamiento cuando cambia el origen de los datos, la granularidad de las etiquetas y el grado de realismo del tráfico. En consecuencia, la comparación no se interpreta como una competencia directa entre datasets, sino como una forma de observar estabilidad metodológica bajo escenarios de datos distintos.

\subsection{Estrategia de unificación del esquema de datos}
Para poder entrenar/evaluar con el mismo pipeline, CSIC y TorpEda se transforman a un esquema compatible con SR-BH:
\begin{itemize}[leftmargin=*]
  \item columnas: \texttt{request\_http\_method}, \texttt{request\_http\_request}, \texttt{request\_body}, y opcionalmente \texttt{request\_headers\_json}.
  \item variables objetivo: \texttt{label\_binary}, \texttt{label\_multiclass}, \texttt{label\_multilabel}.
\end{itemize}

En todos los datasets, para el caso \textit{one-class} se trata todo lo no-normal como anomalía (tarea \textbf{binaria}). Para las tareas supervisadas, la granularidad depende del dataset: \textbf{CSIC~2010} provee únicamente \textit{normal/anómalo}, por lo que las configuraciones multiclase/multietiqueta se evalúan en su \textit{versión reducida} (binaria); \textbf{TorpEda~2012} preserva sus etiquetas \textit{multiclase} nativas (tipo de ataque) y además permite el colapso a binario cuando corresponda; \textbf{SR-BH~2020} habilita binario, multiclase (vía etiqueta CAPEC primaria) y multietiqueta CAPEC.

\subsection{Comparativa de la tríada de datasets}
El diseño experimental conserva los tres datasets porque cada uno aporta una pieza distinta de evidencia metodológica. \textbf{SR-BH~2020} se usa como referencia principal de \textbf{realismo}, ya que contiene tráfico real de honeypot y etiquetas CAPEC multietiqueta \cite{srbh_dataverse,srbh_paper}. \textbf{CSIC~2010} y \textbf{TorpEda~2012} se mantienen porque siguen siendo puntos de comparación frecuentes en detección de ataques web, ofrecen escenarios más controlados y permiten contrastar cuánto del rendimiento observado depende del dataset y cuánto del modelo \cite{csic2010_impact,csic2010_doc,torpeda_2012}. En otras palabras, la tríada se elige para cubrir simultáneamente \textbf{realismo}, \textbf{comparabilidad con la literatura} y \textbf{variación en granularidad de etiquetas}; quitar cualquiera de los tres debilitaría alguna de esas dimensiones.

\subsection{Reducción multietiqueta a multiclase}
En SR-BH una petición puede tener múltiples etiquetas CAPEC. Para definir una tarea multiclase se colapsa la multietiqueta a una sola clase primaria, seleccionando la etiqueta con mayor severidad dentro del conjunto activo \cite{capec_schema_docs,capec_schema_desc}. Esta regla permite obtener una única clase por petición, necesaria para entrenar y evaluar modelos multiclase, pero debe interpretarse con cuidado por dos motivos. Primero, una petición puede combinar patrones de ataque distintos y la clase primaria no conserva toda esa información. Segundo, CAPEC es jerárquico: dos etiquetas activas pueden estar en niveles diferentes de abstracción. Por tanto, la reducción multietiqueta a multiclase se usa como recurso experimental, no como afirmación de que el ataque real tenga una sola interpretación semántica.

\textbf{Ejemplo concreto de la reducción.} Supóngase una petición anotada con dos patrones activos: \texttt{CAPEC-A} y \texttt{CAPEC-B}. Si el mapa de severidad usado por el experimento asigna puntaje 4 a \texttt{CAPEC-A} y puntaje 3 a \texttt{CAPEC-B}, la versión multiclase conserva solo \texttt{CAPEC-A} como etiqueta primaria. La anotación multietiqueta original, en cambio, conserva el conjunto \{\texttt{CAPEC-A}, \texttt{CAPEC-B}\}. La diferencia es metodológicamente importante: en multiclase se mide si el modelo predice la etiqueta primaria elegida por la regla; en multietiqueta se mide si recupera el conjunto de patrones activos. Si dos etiquetas pertenecen a niveles distintos de la jerarquía CAPEC, por ejemplo una categoría más general y un patrón más específico, la reducción a una sola clase puede ocultar esa relación jerárquica.

% ============================================================================
% CAPÍTULO 5
% ============================================================================
\chapter{Metodología propuesta}

\section{Principio general}
La elección de features busca capturar:
\begin{itemize}[leftmargin=*]
  \item \textbf{Estructura}: profundidad del path, cantidad de parámetros, longitudes.
  \item \textbf{Codificación/anormalidad}: percent-encoding, proporción de caracteres no alfanuméricos.
  \item \textbf{Indicadores de payload}: presencia (como feature) de tokens comúnmente asociados a familias de ataque.
  \item \textbf{Método}: one-hot de métodos y bandera de método poco común.
\end{itemize}
En detección de anomalías web, atributos estadísticos y de estructura se han mostrado útiles para puntuar requests anómalos \cite{kruegel_vigna_2003}. En particular, Kruegel y Vigna modelan características de parámetros de consultas web, como longitud, estructura y distribución de caracteres; Ingham e Inoue comparan técnicas de detección de anomalías para HTTP; y Wang y Stolfo muestran la utilidad de modelar contenido/payload a nivel de aplicación \cite{kruegel_vigna_2003,ingham_inoue_2007,wang_stolfo_2004}. Además, trabajos más recientes sobre detección de ataques web y selección de variables muestran que señales extraídas directamente del request HTTP pueden capturar conocimiento experto útil sin abandonar la trazabilidad del modelo \cite{riverol_2024,srbh_paper}. Por ello, se priorizan features \textbf{baratas de extraer}, compatibles con la inspección de un WAF y suficientemente interpretables para discutir después qué señales aportan rendimiento real y cuáles se acercan demasiado a firmas \cite{owasp_waf,modsecurity_site}.


\section{Extracción de características HTTP}
La representación utilizada por los modelos se obtiene directamente a partir de cuatro insumos del request: el método HTTP, la URI, los encabezados y el cuerpo. En la implementación, la extracción se realiza sobre cada petición en forma individual. Primero se normaliza el método a mayúsculas, luego se descompone la URI con \texttt{urlsplit} para separar \textit{path} y \textit{query}, se procesan los parámetros con \texttt{parse\_qsl}, se inspecciona la presencia de secuencias de \textit{percent-encoding} mediante una expresión regular y, finalmente, se revisan URI y body en minúsculas para contabilizar ocurrencias de tokens heurísticos. El resultado es un vector fijo de 25 variables numéricas y binarias. La lista exacta de 25 variables es una decisión de diseño de esta implementación, pero sus familias de señales se apoyan en literatura previa sobre anomalías en HTTP: longitudes y estructura de parámetros, distribución/composición de caracteres y análisis de payload \cite{kruegel_vigna_2003,ingham_inoue_2007,wang_stolfo_2004}.

\subsection{Cómo se construye cada feature}

\subsubsection{Features de longitud}
Las variables de longitud siguen la idea de que las peticiones y sus parámetros pueden modelarse mediante perfiles estadísticos, y que longitudes atípicas pueden señalar payloads o abusos \cite{kruegel_vigna_2003,ingham_inoue_2007}.
\begin{itemize}[leftmargin=*]
  \item \textbf{\texttt{uri\_len}}: representa la longitud total de la URI original. Se calcula como la cantidad de caracteres de la cadena completa recibida en \texttt{uri}. Si la URI está vacía, el valor es 0.

  \item \textbf{\texttt{query\_len}}: representa la longitud del \textit{query string}. Después de separar la URI con \texttt{urlsplit}, se toma el componente \texttt{query} y se cuenta su cantidad de caracteres. Si la URI no tiene parámetros, el valor es 0.

  \item \textbf{\texttt{body\_len}}: representa el tamaño del cuerpo del request en bytes. Se obtiene con \texttt{len(body)} cuando el cuerpo existe. Si el request no trae body, se asigna 0.

  \item \textbf{\texttt{req\_content\_length}}: representa el tamaño declarado del cuerpo. Se intenta leer primero el encabezado \texttt{Content-Length} (aceptando tanto \texttt{content-length} como \texttt{Content-Length}) y convertirlo a entero. Si el encabezado no existe, no es numérico o es negativo, se usa 0. Además, cuando ese valor queda en 0 pero el body existe, la implementación lo reemplaza por \texttt{body\_len} para no perder completamente la señal de tamaño.
\end{itemize}

\subsubsection{Features de estructura}
Las variables de estructura resumen cómo se organiza el recurso solicitado y sus parámetros, siguiendo la misma lógica de describir regularidades del request HTTP antes de evaluar desvíos \cite{kruegel_vigna_2003,ingham_inoue_2007}.
\begin{itemize}[leftmargin=*]
  \item \textbf{\texttt{path\_depth}}: aproxima la profundidad del recurso solicitado. Tras separar la URI, se toma el componente \texttt{path} y se cuenta cuántas veces aparece el carácter \texttt{/}. No mide niveles semánticos perfectos, pero sí ofrece una señal consistente sobre cuán profunda o anidada es la ruta.

  \item \textbf{\texttt{n\_query\_params}}: representa la cantidad de parámetros presentes en la query. Se calcula aplicando \texttt{parse\_qsl(query, keep\_blank\_values=True)} y contando cuántos pares clave--valor fueron extraídos. Esto permite contabilizar también parámetros con valor vacío.

  \item \textbf{\texttt{max\_param\_value\_len}}: representa la mayor longitud observada entre los valores de los parámetros de la query. Una vez obtenidos los pares clave--valor, se calcula el máximo de \texttt{len(valor)}. Si no hay parámetros, el valor es 0.
\end{itemize}

\subsubsection{Features de composición de caracteres y codificación}
Las variables de composición y codificación reflejan que los ataques web suelen alterar la distribución de caracteres, usar símbolos de control o codificar payloads para evadir reglas o filtros \cite{kruegel_vigna_2003,wang_stolfo_2004}.
\begin{itemize}[leftmargin=*]
  \item \textbf{\texttt{uri\_pct\_non\_alnum\_ratio}}: mide la proporción de caracteres no alfanuméricos dentro de la URI completa. Para obtenerla, se recorre la URI carácter por carácter, se cuenta cuántos no satisfacen \texttt{isalnum()}, y luego se divide esa cantidad por \texttt{uri\_len}. Si la URI está vacía, el cociente se define como 0.

  \item \textbf{\texttt{encoded}}: aproxima el grado de \textit{percent-encoding} en la URI. La implementación busca coincidencias del patrón \texttt{\%[0-9a-fA-F]\{2\}}. Cada coincidencia representa una secuencia de tres caracteres codificados, por lo que se multiplica la cantidad de coincidencias por 3 y se divide por \texttt{uri\_len}. El resultado es una razón entre 0 y 1 cuando la longitud es positiva.

  \item \textbf{\texttt{body\_encoded}}: aplica la misma idea sobre el body. Si existe cuerpo, este se decodifica en UTF-8 ignorando errores; sobre ese texto se buscan las mismas secuencias de \textit{percent-encoding}. Luego se divide la longitud total ocupada por esas coincidencias entre la longitud del texto decodificado. Si no hay body, el valor es 0.
\end{itemize}

\subsubsection{Features heurísticas basadas en tokens sospechosos}
En este estudio, los tokens sospechosos no se usan como reglas de bloqueo, sino como señales numéricas adicionales dentro de la representación. Se trata de una lista fija de cadenas elegida para capturar indicios frecuentes de traversal, \textit{cross-site scripting}, inyección SQL y manipulación de payloads. Esta decisión se relaciona con el conocimiento experto que históricamente aparece en reglas y firmas de seguridad web, pero aquí se transforma en variables numéricas para que el modelo aprenda combinaciones con otras señales \cite{owasp_waf,modsecurity_site,riverol_2024}. La lista implementada es la siguiente:

\begin{verbatim}
../, <script, onerror=, onload=, ', ", --, ;,
 union ,  select , sleep(,  or 1=1, benchmark(, information_schema
\end{verbatim}

Antes de contarlos, la URI y el body se convierten a minúsculas. Luego, para cada token de la lista, la implementación aplica una búsqueda por ocurrencias usando \texttt{count()} sobre la cadena correspondiente. En el código, cada token se normaliza con \texttt{strip().lower()} antes de contar, con lo cual los espacios laterales definidos en la lista no forman parte de la coincidencia final. Eso hace que la señal sea simple y reproducible, aunque también más sensible a apariciones parciales dentro del texto.

A partir de esa lista se construyen cuatro variables:

\begin{itemize}[leftmargin=*]
  \item \textbf{\texttt{suspicious\_tokens\_count}}: suma total de ocurrencias de todos los tokens heurísticos en la URI ya pasada a minúsculas.

  \item \textbf{\texttt{has\_suspicious\_tokens}}: indicador binario derivado de la variable anterior. Toma valor 1 cuando \texttt{suspicious\_tokens\_count} es mayor que 0, y 0 en caso contrario.

  \item \textbf{\texttt{body\_suspicious\_tokens\_count}}: suma total de ocurrencias de los mismos tokens, pero en el body decodificado y pasado a minúsculas. Si el request no tiene body, el valor es 0.

  \item \textbf{\texttt{body\_has\_suspicious\_tokens}}: indicador binario derivado de \texttt{body\_suspicious\_tokens\_count}. Vale 1 cuando el conteo del body es mayor que 0, y 0 cuando no se detecta ninguna ocurrencia.
\end{itemize}

Estas variables se describen como heurísticas porque no provienen de un parser semántico ni de una taxonomía aprendida automáticamente. Son señales manuales, sencillas de calcular y cercanas a conocimiento experto del dominio web. Precisamente por eso, más adelante se consideran en forma crítica mediante ablaciones con y sin este grupo de variables; esa lectura evita presentarlas como una solución general basada en firmas y permite medir cuánto aportan frente a señales estructurales más generales.

\subsubsection{Features del método HTTP}
\begin{itemize}[leftmargin=*]
  \item \textbf{\texttt{uncommon\_method}}: indica si el método observado cae fuera del conjunto más común \{\texttt{GET}, \texttt{POST}, \texttt{HEAD}\}. Después de convertir el método a mayúsculas, la variable toma valor 1 si pertenece a otro método distinto de esos tres; en caso contrario toma 0.

  \item \textbf{\texttt{method\_GET}, \texttt{method\_POST}, \texttt{method\_HEAD}, \texttt{method\_PUT}, \texttt{method\_DELETE}, \texttt{method\_PATCH}, \texttt{method\_OPTIONS}, \texttt{method\_TRACE}, \texttt{method\_CONNECT}}: corresponden a una codificación one-hot del método HTTP. Cada una vale 1 únicamente cuando el método observado coincide exactamente con esa opción, y 0 en los demás casos.

  \item \textbf{\texttt{method\_OTHER}}: complementa la codificación anterior. Vale 1 cuando el método no coincide con ninguno de los métodos conocidos contemplados por la implementación; si el método pertenece al conjunto conocido, esta variable vale 0.
\end{itemize}

\subsection{Lectura metodológica de estas variables}
El conjunto completo combina señales estructurales, de tamaño, de codificación, de contenido superficial y de protocolo. Algunas de ellas buscan capturar desvíos estadísticos del request, mientras que otras aproximan indicios más cercanos a firmas. Mantener esa mezcla dentro de un esquema fijo permite comparar modelos sobre una misma representación y, al mismo tiempo, analizar luego qué parte del rendimiento proviene de rasgos generales del tráfico HTTP y qué parte depende de heurísticas léxicas explícitas.

\subsection{Ejemplos de cálculo e interpretación por característica}
El siguiente cuadro resume, para cada característica, un ejemplo concreto de lectura y la clase de señal que puede aportar. Los ejemplos son ilustrativos: una característica aislada no determina por sí sola que exista un ataque; el modelo aprende combinaciones de señales sobre el conjunto de entrenamiento.

{%
\small
\setlength{\tabcolsep}{4pt}%
\renewcommand{\arraystretch}{1.12}%
\begin{longtable}{@{}L{0.25\textwidth}L{0.35\textwidth}L{0.32\textwidth}@{}}
\caption{Ejemplos e interpretación de las 25 características HTTP.}\label{tab:features-ejemplos}\\
\toprule
\textbf{Característica} & \makecell[l]{\textbf{Ejemplo de cálculo}\\\textbf{o activación}} & \textbf{Señal asociada} \\
\midrule
\endfirsthead
\toprule
\textbf{Característica} & \makecell[l]{\textbf{Ejemplo de cálculo}\\\textbf{o activación}} & \textbf{Señal asociada} \\
\midrule
\endhead
\makecell[tl]{\texttt{uri\_len}} & En \texttt{/login?u=admin}, cuenta todos los caracteres de la URI. & URIs muy largas pueden aparecer en inyecciones, fuzzing o payloads extensos. \\
\makecell[tl]{\texttt{query\_len}} & En \texttt{/buscar?q=test}, mide la longitud de \texttt{q=test}. & Queries largas pueden indicar parámetros manipulados o payloads en URL. \\
\makecell[tl]{\texttt{body\_len}} & En un \texttt{POST} con formulario, mide el tamaño del cuerpo recibido. & Cuerpos atípicamente grandes pueden indicar carga maliciosa o abuso de endpoint. \\
\makecell[tl]{\texttt{req\_content\_length}} & Si el header \texttt{Content-Length: 120} existe, toma 120; si no, puede usar \texttt{body\_len}. & Contrasta tamaño declarado y cuerpo; útil para describir requests con body. \\
\makecell[tl]{\texttt{path\_depth}} & \texttt{/a/b/c} contiene tres barras. & Paths profundos o inusuales pueden aparecer en crawling agresivo o traversal. \\
\makecell[tl]{\texttt{n\_query\_params}} & \texttt{/x?a=1\&b=2} tiene dos parámetros. & Muchos parámetros pueden sugerir manipulación de entrada o exploración de superficie. \\
\makecell[tl]{\texttt{max\_param\_value\_len}} & En \texttt{a=1\&q=abcdef}, el máximo es 6. & Valores muy largos pueden transportar payloads de SQLi, XSS u ofuscación. \\
\makecell[tl]{\texttt{uri\_pct\_non\_}\\\texttt{alnum\_ratio}} & En una URI con muchos símbolos como \texttt{?}, \texttt{=}, \texttt{\%} o \texttt{/}, aumenta la razón. & Una proporción alta puede reflejar codificación, operadores o estructura anómala. \\
\makecell[tl]{\texttt{encoded}} & En \texttt{/x?u=..\%2F..\%2Fetc}, las secuencias \texttt{\%2F} aumentan la razón. & El percent-encoding puede usarse para evasión u ofuscación. \\
\makecell[tl]{\texttt{body\_encoded}} & Un cuerpo con \texttt{\%3Cscript\%3E} incrementa esta razón. & Señal de payload codificado en el cuerpo. \\
\makecell[tl]{\texttt{suspicious\_tokens\_}\\\texttt{count}} & En \texttt{/x?id=1 union select}, cuenta ocurrencias de \texttt{union} y \texttt{select}. & Indicios léxicos de SQLi u otras familias conocidas. \\
\makecell[tl]{\texttt{has\_suspicious\_}\\\texttt{tokens}} & Vale 1 si el conteo anterior es mayor que 0. & Resume presencia de señales léxicas sospechosas en la URI. \\
\makecell[tl]{\texttt{body\_suspicious\_}\\\texttt{tokens\_count}} & Un body con \texttt{<script>} o \texttt{or 1=1} suma ocurrencias. & Indicios de XSS o SQLi enviados por formulario o JSON. \\
\makecell[tl]{\texttt{body\_has\_}\\\texttt{suspicious\_tokens}} & Vale 1 si el body contiene al menos un token sospechoso. & Resume presencia de señales sospechosas en el cuerpo. \\
\makecell[tl]{\texttt{uncommon\_method}} & Vale 1 para métodos fuera de \texttt{GET}, \texttt{POST} y \texttt{HEAD}. & Métodos poco habituales pueden indicar pruebas, escaneo o abuso de protocolo. \\
\makecell[tl]{\texttt{method\_GET}} & Vale 1 si el método es \texttt{GET}. & Distingue tráfico de consulta o navegación. \\
\makecell[tl]{\texttt{method\_POST}} & Vale 1 si el método es \texttt{POST}. & Distingue envíos con cuerpo, formularios o APIs. \\
\makecell[tl]{\texttt{method\_HEAD}} & Vale 1 si el método es \texttt{HEAD}. & Puede aparecer en verificaciones livianas o sondeos. \\
\makecell[tl]{\texttt{method\_PUT}} & Vale 1 si el método es \texttt{PUT}. & En endpoints no esperados puede sugerir intento de escritura. \\
\makecell[tl]{\texttt{method\_DELETE}} & Vale 1 si el método es \texttt{DELETE}. & En APIs puede ser legítimo; fuera de contexto puede indicar abuso. \\
\makecell[tl]{\texttt{method\_PATCH}} & Vale 1 si el método es \texttt{PATCH}. & Señal de modificación parcial; depende del endpoint observado. \\
\makecell[tl]{\texttt{method\_OPTIONS}} & Vale 1 si el método es \texttt{OPTIONS}. & Puede indicar exploración de métodos permitidos. \\
\makecell[tl]{\texttt{method\_TRACE}} & Vale 1 si el método es \texttt{TRACE}. & Método sensible en seguridad web; puede indicar prueba de configuración. \\
\makecell[tl]{\texttt{method\_CONNECT}} & Vale 1 si el método es \texttt{CONNECT}. & Inusual en muchas aplicaciones web; puede indicar intento de túnel. \\
\makecell[tl]{\texttt{method\_OTHER}} & Vale 1 si el método no coincide con los contemplados. & Captura métodos desconocidos, malformados o específicos de algún cliente. \\
\bottomrule
\end{longtable}
}

\section{Criterios de diseño de la representación}
\begin{itemize}[leftmargin=*]
  \item \textbf{Longitudes y estructura}: ataques suelen introducir payloads largos o patrones estructurales (query con muchos parámetros, paths profundos). Esto está alineado con enfoques basados en atributos del request para detección de anomalías y con trabajos de selección de features en tráfico web \cite{kruegel_vigna_2003,riverol_2024}.
  \item \textbf{Percent-encoding y no-alfanuméricos}: evasión y ofuscación frecuentemente se apoyan en \textit{encoding} y caracteres especiales; modelarlos ayuda a capturar desviaciones que un WAF también inspecciona \cite{owasp_waf,modsecurity_site}.
  \item \textbf{Tokens sospechosos como features (no como reglas)}: se usan como señales numéricas que pueden ayudar a modelos lineales/no lineales, pero deben discutirse críticamente porque se aproximan a firmas; por eso se justifican solo si luego se auditan con ablaciones \cite{kruegel_vigna_2003,riverol_2024}.
  \item \textbf{One-hot de métodos}: el método HTTP condiciona la forma del request; modelarlo evita confundir variaciones legítimas del protocolo con anomalías y mantiene coherencia con campos efectivamente inspeccionados en entornos WAF \cite{owasp_waf,modsecurity_site}.
  \item \textbf{\texttt{req\_content\_length}}: cuando está disponible, aporta una señal adicional de tamaño esperable del body con costo de extracción prácticamente nulo, útil para mantener compatibilidad entre datasets y enriquecer la descripción estructural del request \cite{srbh_paper,owasp_waf}.
\end{itemize}

% ============================================================================
% CAPÍTULO 6
% ============================================================================

\section{Implementación del pipeline}

\subsection{Arquitectura del pipeline}
\begin{figure}[H]
\centering
\resizebox{0.98\textwidth}{!}{%
\begin{tikzpicture}[
  block/.style={draw, rounded corners, align=center, minimum width=3.8cm, minimum height=1.0cm},
  arrow/.style={-Latex, thick},
  small/.style={font=\small}
]
\node[block] (raw) {Datasets crudos\\SR-BH / CSIC / TorpEda};
\node[block, right=1.2cm of raw] (prep) {Parsing HTTP\\método/URI/body};
\node[block, right=1.2cm of prep] (feat) {Extracción\\25 features};
\node[block, right=1.2cm of feat] (ds) {Dataset final\\(parquet/csv)};

\node[block, below=1.2cm of ds] (train1) {Entrenamiento\\one-class escalable};
\node[block, below=1.2cm of feat] (train2) {Entrenamiento\\supervisado lineal};
\node[block, below=1.2cm of raw] (eval) {Evaluación\\métricas + FI};

\draw[arrow] (raw) -- (prep);
\draw[arrow] (prep) -- (feat);
\draw[arrow] (feat) -- (ds);
\draw[arrow] (ds) -- (train1);
\draw[arrow] (ds) -- (train2);
\draw[arrow] (train1) -- (eval);
\draw[arrow] (train2) -- (eval);

\node[font=\scriptsize, align=center, anchor=south, yshift=4mm, fill=white, inner sep=1pt]
  at (raw.north) {srbh\_loader\\csic\_loader\\torpeda\_loader};

\node[font=\scriptsize, anchor=south, yshift=4mm, fill=white, inner sep=1pt]
  at (feat.north) {http\_features.py};

\node[small] at ($(raw)+(0,-1.35)$) {build\_features*.py};
\node[small] at ($(train1)+(0,-0.95)$) {train\_oneclass.py};
\node[small] at ($(train2)+(0,-0.95)$) {train\_supervised.py};
\end{tikzpicture}%
}
\caption{Pipeline propuesto: unificación de datasets, extracción de features, entrenamiento escalable, tuning con presupuesto fijo y evaluación.}
\label{fig:pipeline}
\end{figure}

\subsection{Artefactos generados}

El pipeline produce artefactos intermedios y finales que permiten reproducir las corridas y auditar los resultados. Los principales artefactos son los siguientes:

\begin{itemize}[leftmargin=*]
  \item \textbf{Datasets procesados:} archivos tabulares en formato CSV o Parquet con los campos HTTP normalizados, las 25 características calculadas y las variables objetivo \texttt{label\_binary}, \texttt{label\_multiclass} y \texttt{label\_multilabel} cuando corresponda.
  \item \textbf{Particiones experimentales:} registros de tamaños de entrenamiento, validación y prueba, junto con la semilla utilizada y la distribución de clases o etiquetas. Estos artefactos permiten verificar que el conjunto de prueba externo no fue utilizado durante el ajuste.
  \item \textbf{Modelos entrenados:} modelos serializados para las tres familias evaluadas: detector \textit{one-class}, regresión logística multiclase y One-vs-Rest multietiqueta. Cuando corresponde, también se conserva el escalador y cualquier transformación necesaria para reproducir la inferencia.
  \item \textbf{Resultados de búsqueda de hiperparámetros:} tablas con las configuraciones evaluadas, métricas internas y configuración seleccionada. Estos archivos documentan el costo y el criterio de selección del ajuste.
  \item \textbf{Métricas de efectividad:} tablas con accuracy, balanced accuracy, precision, recall, $F_1$, MCC, ROC-AUC, PR-AUC y métricas específicas de multietiqueta según la tarea ejecutada.
  \item \textbf{Métricas operativas:} archivos con latencias, rendimiento, consumo aproximado de CPU, memoria, tamaño del modelo, tiempo de entrenamiento y tiempo de aplicación del flujo de extracción de características más inferencia.
  \item \textbf{Artefactos de interpretabilidad:} coeficientes estandarizados de los modelos lineales, resúmenes por clase o etiqueta y resultados de las ablaciones con y sin variables basadas en tokens sospechosos.
\end{itemize}

Estos artefactos cumplen una doble función. En primer lugar, permiten reconstruir el flujo experimental desde los datos procesados hasta las tablas finales. En segundo lugar, facilitan la trazabilidad de las decisiones metodológicas, especialmente en los puntos donde se reduce una salida multietiqueta a multiclase, se eliminan variables por ablación o se seleccionan hiperparámetros dentro del bloque de entrenamiento.

\subsection{Derivación de \texttt{label\_multiclass} a partir de CAPEC (SR-BH)}
En SR-BH, una misma request puede activar varias etiquetas CAPEC al mismo tiempo. Esa es precisamente la razón por la que el dataset resulta valioso para estudiar una formulación multietiqueta. Sin embargo, para poder compararlo con tareas multiclase sobre un mismo pipeline, la implementación también necesita una única etiqueta principal por instancia.

La reducción a multiclase se resuelve mediante una regla determinista basada en severidad. Primero se construye \texttt{label\_multilabel} como la lista de etiquetas CAPEC activas en cada fila. Luego, cuando la estrategia seleccionada es \texttt{severity}, se elige como \texttt{label\_multiclass} la etiqueta con mayor severidad según el mapa definido en el loader. Si no hay etiquetas activas, la instancia queda como \texttt{NORMAL}. Esta decisión no elimina la salida multietiqueta original: ambas conviven para que la comparación entre formulaciones no obligue a perder información semántica.

La implementación conserva además una columna auxiliar \texttt{label\_multiclass\_first}, útil para auditoría, y una columna \texttt{label\_primary\_severity} que deja trazado el puntaje de severidad asociado a la etiqueta elegida. De esta manera, la reducción a multiclase no queda como una operación opaca, sino como una transformación explícita y revisable dentro del mismo dataset procesado.

\textbf{Ejemplo operativo de la regla.} Si una fila tiene \texttt{label\_multilabel=\{CAPEC-A, CAPEC-B\}} y el mapa de severidad asigna \texttt{CAPEC-A}=4 y \texttt{CAPEC-B}=3, entonces \texttt{label\_multiclass} toma el valor \texttt{CAPEC-A} y \texttt{label\_primary\_severity} toma el valor 4. Si la fila no tiene etiquetas CAPEC activas, \texttt{label\_multiclass} toma el valor \texttt{NORMAL}. Esta regla permite reproducir exactamente la conversión y auditar qué información se pierde al pasar de una salida multietiqueta a una única clase.

Desde el punto de vista metodológico, esta política se usa solo para habilitar una tarea adicional de comparación. La salida multietiqueta sigue siendo la representación más fiel del fenómeno cuando una misma petición participa de más de un patrón CAPEC. Por eso, en la discusión de resultados, la lectura multiclase se interpreta como una simplificación útil para contraste experimental, no como sustituto de la anotación original.


\chapter{Diseño experimental}

\section{Tareas definidas}
\begin{itemize}[leftmargin=*]
  \item \textbf{Detección binaria}: \texttt{label\_binary} (0 normal, 1 ataque/anómalo).
  \item \textbf{Multiclase}: \texttt{label\_multiclass}. En SR-BH: etiqueta CAPEC primaria derivada desde la multietiqueta con una política explícita; en TorpEda: clases nativas por tipo de ataque; en CSIC~2010 (sin taxonomía por tipo), se usa la versión reducida $K=2$ (normal vs anómalo).
  \item \textbf{Multietiqueta}: \texttt{label\_multilabel}. Disponible plenamente en SR-BH (vector/lista de etiquetas CAPEC por request). Para ejecutar el baseline OvR sobre los tres datasets, se parametriza la salida según etiquetas disponibles: CSIC~2010 se restringe a binario y TorpEda se usa en modo multiclase o binario cuando corresponda.
\end{itemize}

\section{Matriz experimental (3 modelos $\times$ 3 datasets)}
En términos operativos, el diseño experimental es una matriz completa: \textbf{cada familia de modelos se entrena y evalúa sobre cada uno de los tres datasets}, adaptando la granularidad de la tarea cuando el dataset no provee todas las etiquetas.

\begin{table}[H]
\centering
\caption{Matriz experimental (modelo $\times$ dataset) y tarea evaluada.}
\label{tab:matriz_experimental_libro}
\footnotesize
\setlength{\tabcolsep}{4pt}
\renewcommand{\arraystretch}{1.15}
\begin{tabularx}{\textwidth}{@{}>{\RaggedRight\arraybackslash}p{4.8cm}>{\RaggedRight\arraybackslash}X>{\RaggedRight\arraybackslash}X>{\RaggedRight\arraybackslash}X@{}}
\toprule
\textbf{Modelo / configuración} & \textbf{CSIC~2010} & \textbf{TorpEda~2012} & \textbf{SR-BH~2020} \\
\midrule
\parbox[t]{\linewidth}{Detector one-class escalable\\(\texttt{SGDOneClassSVM} +\\aproximación de kernel)} &
Binario (normal vs.\ anómalo) &
Binario (normal vs.\ anómalo) &
Binario (normal vs.\ anómalo) \\
Regresión logística lineal &
Binario (versión reducida) &
Multiclase (tipo de ataque) &
Multiclase (CAPEC primaria) \\
OvR lineal &
Binario (versión reducida) &
Multiclase o binario, según etiqueta disponible &
Multietiqueta (CAPEC) \\
\bottomrule
\end{tabularx}
\end{table}

\section{Política general de particionado y reproducibilidad}
La comparación entre modelos dentro de un mismo dataset debe hacerse sobre una \textbf{misma política de particionado}. Por ello, se adopta como protocolo principal un \textbf{particionado externo fijo 80/20} para los experimentos supervisados, con \textbf{semilla explícita y reproducible}. En esos casos, el \textbf{20\% externo} queda reservado exclusivamente para \textbf{prueba final}; no se reutiliza para ajuste de hiperparámetros, selección de variables, calibración de umbrales ni ajuste de transformaciones. En consecuencia, bajo este protocolo no hay contaminación entre entrenamiento y prueba: el conjunto de prueba externo permanece aislado hasta la evaluación final. Para el \textit{detector one-class}, la política específica es entrenar con el \textbf{80\% de instancias normales} y evaluar sobre un conjunto mixto compuesto por el \textbf{20\% de normales en holdout} y el \textbf{total de anomalías/ataques disponibles}.

Como criterio adicional de trazabilidad:
\begin{itemize}[leftmargin=*]
  \item se registran las semillas, tamaños de partición y prevalencias de clases/etiquetas;
  \item cualquier transformación (\texttt{StandardScaler}, \texttt{MultiLabelBinarizer}, codificación u otras) se ajusta únicamente con datos de entrenamiento;
  \item cuando un dataset disponga de un \textit{split} canónico propio (por ejemplo, CSIC~2010), este podrá reportarse como \textbf{análisis secundario}, pero la comparación principal entre modelos/datasets se basará en el protocolo 80/20 unificado.
\end{itemize}

\section{Particionado por tarea}
\subsection{Supervisado (multiclase y multietiqueta)}
En los modelos supervisados, el \textbf{80\% de entrenamiento} se utiliza para ajuste y selección de hiperparámetros; el \textbf{20\% externo} se usa una única vez para evaluación final.

\paragraph{Multiclase.}
Se intentará \texttt{train\_test\_split} estratificado por clase. Si alguna clase tiene cardinalidad insuficiente para sostener la estratificación o la validación cruzada, se aplicará una política explícita de clases raras previa al \textit{split}: \texttt{keep}, \texttt{merge} (por ejemplo, a \texttt{OTHER\_RARE}) o \texttt{drop}. Si aun así la estratificación no fuera viable, se utilizará un \textit{split} no estratificado con semilla fija y se reportará la prevalencia final por clase.

\paragraph{Multietiqueta.}
Cuando sea técnicamente viable, se preferirá una estratificación aproximada que preserve la distribución conjunta de etiquetas. Si la esparsidad o las restricciones de implementación lo impiden, se utilizará como baseline un \textit{split} fijo no estratificado y se reportarán las prevalencias por etiqueta antes y después del particionado. Esa condición se documenta como limitación y se complementa con análisis por etiqueta para evitar lecturas engañosas.

\subsection{Detector one-class sin contaminación}
Para el detector one-class, la política debe respetar dos condiciones: (i) entrenar solo con tráfico normal y (ii) mantener un conjunto de prueba verdaderamente no visto. Por ello, la estrategia adoptada es la siguiente:
\begin{enumerate}[leftmargin=*]
  \item Se realiza un \textbf{particionado 80/20 únicamente sobre las instancias normales}. El \textbf{20\% de normales en holdout} se reserva para prueba final.
  \item El modelo final se ajusta \emph{únicamente} con el \textbf{80\% de normales}. El \textbf{total de anomalías/ataques disponibles} se reserva para evaluación final y no participa del ajuste del modelo.
  \item Si hay \textit{tuning}, se realiza \textbf{solo dentro del bloque de entrenamiento normal}: se separa una validación interna de normales desde el 80\% de normales y se seleccionan hiperparámetros con un criterio coherente con el escenario \textit{one-class} (por ejemplo, minimizar rechazo de tráfico normal o controlar una tasa objetivo de outliers), sin tocar el conjunto final de ataques.
\end{enumerate}
Con ello se elimina el principal riesgo metodológico de este tipo de detector: usar el mismo conjunto mixto tanto para escoger hiperparámetros como para reportar métricas finales, y además se aprovechan \textbf{todas las anomalías/ataques disponibles} en la evaluación final.

\section{Modelos seleccionados}
Se mantienen las tres familias metodológicas definidas para el estudio, con una implementación ajustada para conservar compatibilidad con datasets muy grandes. La lógica de selección no es ``usar el modelo más complejo disponible'', sino elegir un conjunto de referencias que permita responder las preguntas del estudio con \textbf{trazabilidad}, \textbf{comparabilidad} y \textbf{factibilidad computacional}. Modelos más complejos, incluyendo enfoques recientes basados en redes profundas, son útiles como contexto de literatura, pero no necesariamente como \textit{baseline} principal cuando el objetivo incluye interpretabilidad, evaluación cruzada entre datasets y ejecución sobre volúmenes del orden de $10^6$ peticiones \cite{fang_2024,zhou_2024}.

\subsection{Detector one-class escalable}

\subsubsection{Justificación del enfoque one-class}
El uso de un detector \textit{one-class} se justifica en este trabajo por dos razones complementarias. En primer lugar, permite entrenar el modelo exclusivamente con tráfico normal y detectar como anómalas las peticiones que se desvían de ese soporte aprendido. Esto resulta útil cuando no se dispone de un catálogo completo de ataques etiquetados, situación frecuente en entornos reales donde los atacantes generan constantemente variantes nuevas. En segundo lugar, este enfoque mantiene coherencia con la motivación original de los WAF basados en anomalías: aprender qué aspecto tiene el tráfico legítimo y señalar desviaciones, sin depender de un inventario fijo de patrones maliciosos \cite{kruegel_vigna_2003}.

\subsubsection{Justificación de la variante escalable}
La formulación clásica del One-Class SVM \cite{scholkopf2001} resuelve un problema de optimización cuadrática cuyo costo crece de manera cuadrática o cúbica con el número de ejemplos, lo que lo hace impracticable para datasets del orden de $10^5$ a $10^6$ peticiones. Por este motivo se adopta una variante escalable basada en descenso de gradiente estocástico (SGD), que mantiene el mismo principio conceptual —aprender la frontera del soporte normal— pero con un costo de entrenamiento lineal en el número de ejemplos \cite{sklearn_sgd_ocsvm}. Para conservar la capacidad de capturar relaciones no lineales entre características, se agrega una aproximación de kernel mediante el método de Nyström \cite{williams2001,sklearn_kernel_approx}, que transforma explícitamente el espacio de características antes de aplicar el detector lineal.

\subsubsection{Pipeline concreto adoptado}
La implementación seleccionada para todas las corridas one-class del estudio es el pipeline en tres etapas de la Ecuación~\eqref{eq:oneclass-pipeline}:
\begin{equation}
\label{eq:oneclass-pipeline}
\texttt{StandardScaler} \;\rightarrow\; \texttt{Nystroem} \;\rightarrow\; \texttt{SGDOneClassSVM}.
\end{equation}
Sea $x \in \mathbb{R}^d$ el vector de características de una petición HTTP. La primera etapa escala cada característica a media cero y varianza unitaria. La segunda aplica una transformación explícita $z(x) \in \mathbb{R}^m$ mediante la aproximación de Nyström, que aproxima el mapa de características asociado a un kernel RBF. La tercera etapa aprende un hiperplano en ese espacio transformado que separa el tráfico normal del origen según la función de decisión definida en la Ecuación~\eqref{eq:ocsvm-decision}. Las peticiones con $f(x) < 0$ se clasifican como anómalas. Esta formulación preserva el supuesto central del enfoque ---entrenar solo con tráfico normal--- y mantiene el costo de entrenamiento e inferencia en un rango compatible con el tamaño de los datasets utilizados \cite{sklearn_sgd_ocsvm,sklearn_kernel_approx}.

\subsection{Baseline supervisado: regresión logística lineal}
Se adopta regresión logística lineal como \textit{baseline} supervisado principal, preferentemente con \texttt{solver=lbfgs} o configuración equivalente eficiente, por su interpretabilidad mediante coeficientes, su costo computacional controlado y su solidez como referencia supervisada en la literatura \cite{islr_logreg,sklearn_logreg}.

\subsection{Baseline multietiqueta: One-vs-Rest}
Para la formulación multietiqueta se utiliza \texttt{OneVsRestClassifier} sobre regresión logística lineal. OvR es una descomposición estándar y metodológicamente defendible cuando interesa conservar trazabilidad por etiqueta y mantener interpretabilidad de los coeficientes por clase \cite{tsoumakas_multilabel_overview,rifkin_ova,sklearn_logreg}.

\section{Selección de hiperparámetros (Grid/Random Search) y validación}
La metodología incorpora un módulo de \textit{tuning} para alinear el pipeline con buenas prácticas experimentales, pero lo redefine para que no domine el costo total de ejecución. La política general es: \textbf{toda selección se hace solo dentro del bloque de entrenamiento} y con un \textbf{presupuesto computacional fijo} \cite{sklearn_randomsearch,sklearn_halving}.

\subsection{Tuning en supervisado (multiclase y multietiqueta)}
\begin{itemize}[leftmargin=*]
  \item \textbf{Bloque habilitado:} únicamente el \textbf{80\% de entrenamiento} del split externo.
  \item \textbf{Estrategia recomendada:} \textit{random search} acotado o \textit{successive halving} sobre un subconjunto interno del entrenamiento, evitando \textit{grid search} exhaustivo como protocolo principal.
  \item \textbf{Parámetros buscados:} principalmente \texttt{C}, penalización, \texttt{solver} y, cuando corresponda, \texttt{class\_weight}.
  \item \textbf{Métrica de selección:} $F_1$ weighted en multiclase y $F_1$ micro en multietiqueta.
  \item \textbf{Reentrenamiento final:} con la mejor configuración, el modelo se reentrena sobre todo el \textbf{80\% de entrenamiento} y se evalúa una única vez sobre el \textbf{20\% externo}.
  \item \textbf{Registro:} se guardan los scores por candidato a CSV para trazabilidad.
\end{itemize}

\subsection{Tuning en el detector one-class}
\begin{itemize}[leftmargin=*]
  \item \textbf{Bloque habilitado:} solo normales del \textbf{80\% de entrenamiento}; el test final nunca se toca durante la selección.
  \item \textbf{Estrategia recomendada:} búsqueda acotada sobre \texttt{n\_components}, \texttt{gamma}, regularización y número de épocas del backend escalable. 
  \item \textbf{Métrica de selección:} una métrica basada en normales del \textit{holdout} interno (por ejemplo, tasa de aceptación de normales o tasa de rechazo falsa). Si se usa una variante con anomalías internas, debe declararse como análisis secundario y no como protocolo principal.
  \item \textbf{Reentrenamiento final:} con la mejor configuración, el modelo se reentrena sobre todo el \textbf{80\% de normales} y se evalúa una única vez sobre el \textbf{20\% de normales en holdout} + \textbf{todas las anomalías/ataques disponibles}.
  \item \textbf{Registro:} se guardan los scores por candidato a CSV para trazabilidad.
\end{itemize}


\section{Configuración experimental efectivamente utilizada}
Las corridas incorporadas en los resultados finales siguen una parametrización homogénea por familia de modelos. La regularidad entre datasets no se utilizó solo como criterio de orden, sino como condición para sostener comparaciones transversales defendibles: cuando la tarea cambia, cambia la granularidad de la etiqueta; cuando la familia del modelo se mantiene, también se mantienen la lógica de particionado, el presupuesto de búsqueda y el bloque operativo de benchmark.

\subsection{Construcción y preparación de los datasets}
Los tres datasets se procesan sobre una representación común de características HTTP. En CSIC~2010 y CSIC/TorpEda~2012 se conserva el uso de información de cabeceras cuando está disponible, mientras que en SR-BH~2020 se mantiene el preprocesamiento compatible con separadores heterogéneos. El criterio adoptado es preservar un esquema de entrada comparable entre fuentes con distinto origen, sin introducir una definición distinta de características para cada dataset. De ese modo, las diferencias observadas en resultados quedan asociadas principalmente al realismo del tráfico, a la granularidad de las etiquetas y a la familia del modelo, y no a una representación de entrada completamente distinta.

\subsection{Correspondencia entre representación y variables objetivo}

\subsubsection{Vector de características}
En la implementación, cada petición HTTP se representa como un vector fijo de 25 características numéricas ($d = 25$), construido a partir de cuatro campos del protocolo: método HTTP, URI, cuerpo (\textit{body}) y, cuando está disponible, un campo derivado de encabezados (\texttt{req\_content\_length}). Las características cubren longitudes, estructura de la URI, composición de caracteres, grado de \textit{percent-encoding}, indicadores heurísticos basados en tokens sospechosos y una codificación \textit{one-hot} del método. El detalle de cada característica se presenta en el capítulo de representación de peticiones HTTP.

En el caso de SR-BH~2020, el dataset publicado originalmente incluye 24 características; en la versión adoptada para este estudio se recalcula el esquema completo y se agrega la característica \texttt{req\_content\_length} derivada de encabezados, unificando el conjunto en 25 para mantener comparabilidad con los demás datasets.

\subsubsection{Variables objetivo}
Las variables de salida utilizadas en los experimentos corresponden a tres formulaciones del problema, derivadas de la información de etiquetado disponible en cada dataset:

\begin{itemize}[leftmargin=*]
  \item \textbf{\texttt{label\_binary}:} variable binaria (0 = normal, 1 = ataque o anomalía). Disponible en los tres datasets.
  \item \textbf{\texttt{label\_multiclass}:} variable con una única clase por petición. En TorpEda~2012 corresponde a las clases nativas por tipo de ataque. En SR-BH~2020 se deriva desde la anotación multietiqueta CAPEC mediante una política determinista basada en severidad: se selecciona como clase primaria la etiqueta CAPEC con mayor nivel de severidad entre las activas en esa petición; si no hay etiquetas activas, la instancia queda como \texttt{NORMAL}. Para trazabilidad, el dataset procesado conserva columnas auxiliares \texttt{label\_multiclass\_first} y \texttt{label\_primary\_severity}. En CSIC~2010, donde no existe taxonomía por tipo de ataque, esta variable queda reducida de hecho a un esquema binario.
  \item \textbf{\texttt{label\_multilabel}:} vector binario de longitud $L$ (cantidad de etiquetas CAPEC posibles), donde cada componente indica si la petición activa esa etiqueta. Disponible plenamente en SR-BH~2020. En los demás datasets, la formulación multietiqueta se ejecuta en modo reducido según la granularidad efectivamente disponible.
\end{itemize}

\subsection{Detector one-class en CSIC~2010, CSIC/TorpEda~2012 y SR-BH~2020}
La familia \textit{one-class} se ejecuta en los tres datasets como problema binario \textit{normal vs. anómalo}, manteniendo una misma estructura: escalado, aproximación de kernel mediante Nyström y detector \texttt{SGDOneClassSVM}. El particionado conserva \textbf{80/20} sobre tráfico normal, con \textbf{semilla 42}, de modo que el ajuste se realiza exclusivamente con normales y la evaluación final utiliza el \textit{holdout} de normales junto con el total de anomalías o ataques disponibles. La búsqueda de hiperparámetros se limita a un presupuesto aleatorio pequeño sobre cuatro ejes: fracción esperada de outliers, escala del kernel aproximado, dimensionalidad de la aproximación y número máximo de iteraciones.

Esa configuración cumple dos propósitos. En primer lugar, preserva el supuesto metodológico central del enfoque \textit{one-class}: aprender la frontera del tráfico legítimo sin contaminar el entrenamiento con ataques. En segundo lugar, evita que el costo de explorar hiperparámetros vuelva impracticable una familia de modelos cuya principal justificación en el estudio es precisamente la escalabilidad. El bloque operativo también se mantiene uniforme, con un benchmark acotado y repetición mínima, suficiente para medir latencia, \textit{throughput}, memoria y estabilidad sin convertir la evaluación operativa en una prueba de carga ajena al alcance definido.

Además, en los tres datasets se ejecutan dos variantes: una base y otra sin variables de tokens sospechosos. La ablación mantiene intacto el resto del pipeline. Esa simetría permite interpretar cualquier diferencia de rendimiento como efecto del grupo de variables removido y no como efecto colateral de otra modificación simultánea.

\subsection{Regresión logística multiclase en CSIC~2010, CSIC/TorpEda~2012 y SR-BH~2020}
La familia multiclase se ejecuta con regresión logística lineal, \texttt{solver=lbfgs}, penalización \texttt{l2}, regularización inicial \texttt{C=1.0}, ponderación de clases \texttt{balanced}, máximo de 500 iteraciones y política explícita de clases raras basada en \texttt{keep} con mínimo de dos instancias. Sobre esa base fija, la búsqueda de hiperparámetros se restringe a una validación cruzada de tres particiones, presupuesto aleatorio 2 y una muestra interna acotada, explorando únicamente dos valores de regularización y dos alternativas de ponderación de clases.

La justificación de esta configuración no reside en exhaustividad, sino en estabilidad. \texttt{lbfgs} con \texttt{l2} conserva una línea base lineal robusta, con convergencia previsible y lectura posterior por coeficientes. La exploración se concentra en regularización y ponderación de clases porque, dentro de una familia lineal de baja complejidad, esos son los dos ejes con mayor impacto práctico bajo desbalance. Fijar el optimizador y la penalización reduce el espacio de búsqueda, disminuye ruido experimental y hace más clara la comparación entre datasets.

En SR-BH~2020 la tarea multiclase se apoya en la etiqueta CAPEC primaria derivada desde la anotación multietiqueta. En CSIC/TorpEda~2012 se conserva la taxonomía nativa por tipo de ataque. En CSIC~2010 la misma tubería se mantiene, pero la salida queda reducida de hecho a un problema binario, ya que el dataset no aporta una taxonomía multiclase rica. Esa configuración no contradice la matriz experimental: muestra, por el contrario, qué ocurre cuando una misma familia supervisada se aplica a datasets con distinta granularidad semántica.

También aquí se conservan dos variantes, base y sin tokens sospechosos, bajo el mismo particionado externo, el mismo presupuesto de \textit{tuning} y el mismo benchmark operativo. Esa regularidad refuerza la trazabilidad del contraste.

\subsection{OvR lineal para multietiqueta y modos reducidos por dataset}
La familia multietiqueta se ejecuta mediante una descomposición \textit{One-vs-Rest} sobre regresión logística lineal con \texttt{solver=lbfgs}, penalización \texttt{l2}, \texttt{class\_weight=balanced}, máximo de 300 iteraciones, tolerancia \texttt{1e-3} y paralelismo habilitado en el entrenamiento de las salidas. La búsqueda de hiperparámetros mantiene el mismo criterio acotado que en multiclase: validación cruzada interna, presupuesto 2, muestra interna limitada y exploración restringida a regularización y ponderación de clases.

La formulación plena multietiqueta se conserva en SR-BH~2020, que sí ofrece etiquetas CAPEC simultáneas por request. En CSIC~2010 se fuerza un modo binario reducido, porque el dataset no dispone de una salida multietiqueta real. En CSIC/TorpEda~2012 se utiliza un modo reducido automático, que adapta la salida a la granularidad efectivamente disponible. El esquema conserva una única familia de script y de evaluación para los tres datasets, pero sin atribuir una riqueza semántica que dos de ellos no poseen. De esa manera, la comparación conserva continuidad operacional sin introducir etiquetas artificiales.

La ablación de tokens sospechosos en esta familia se ejecuta dentro del mismo flujo del modelo y se registra como comparación específica entre variante base y variante ablacionada. El procedimiento evita duplicar lógica de entrenamiento y mantiene alineadas las métricas exportadas para ambas variantes. En el bloque operativo, el benchmark utiliza una ventana algo más amplia de calentamiento que en multiclase, porque la descomposición por etiqueta introduce un costo superior de inferencia y de carga en memoria.

\subsection{Lectura metodológica conjunta}
En conjunto, la forma de ejecución de cada familia responde a un mismo criterio: \textbf{máxima comparabilidad con complejidad controlada}. El detector \textit{one-class} conserva su supuesto de entrenamiento exclusivo con normales; la regresión logística multiclase mantiene una línea base lineal estable y fácilmente interpretable; y la formulación OvR multietiqueta preserva trazabilidad por etiqueta sin introducir arquitecturas más pesadas que desplacen la discusión central. Los presupuestos de \textit{tuning} son pequeños en los tres casos, no porque la optimización carezca de valor, sino porque la prioridad metodológica del estudio consiste en sostener una matriz experimental reproducible, coherente con el hardware disponible y suficiente para responder las preguntas de comparación planteadas.

\section{Métricas (efectividad y eficiencia)}
Dado que uno de los objetivos es comparar con trabajos relacionados, se selecciona un conjunto de métricas recurrente en la literatura sobre detección y clasificación de ataques web. El criterio de selección de métricas parte de que \textbf{accuracy por sí sola no alcanza} en presencia de desbalance, clases raras y tareas distintas (binaria, multiclase y multietiqueta). Por eso se combinan métricas de rendimiento, métricas sensibles al tipo de error y métricas operativas de costo temporal \cite{saito_pr_auc,chicco_mcc,srbh_paper,tsoumakas_multilabel_overview}.

\begin{itemize}[leftmargin=*]
  \item \textbf{Binario (normal vs ataque/anómalo):} matriz de confusión, Accuracy, Balanced Accuracy, Precision, Recall/TPR, F1, Specificity/TNR, FPR/FNR, MCC, ROC-AUC y PR-AUC. PR-AUC se incluye porque es más informativa que ROC-AUC cuando la clase positiva es minoritaria \cite{saito_pr_auc}.
  \item \textbf{Multiclase:} Accuracy, Balanced Accuracy, Precision/Recall/$F_1$ en agregación micro, macro y \textit{weighted}, además de MCC, matriz de confusión y métricas por clase. Se reportan varias agregaciones para no ocultar el comportamiento sobre clases raras detrás de un único promedio \cite{tsoumakas_multilabel_overview,zhou_2024}.
  \item \textbf{Multietiqueta (CAPEC):} Hamming Loss, Jaccard Similarity (micro, macro y por muestra cuando corresponde), Exact Match Ratio, Precision/Recall/$F_1$ micro y macro, y AUC micro/macro cuando se dispone de scores continuos. Estas métricas cubren error por etiqueta, solapamiento entre conjuntos y coincidencia exacta de la predicción multilabel. En modos binarios reducidos, Jaccard no se toma como métrica principal, sino como salida secundaria de consistencia \cite{srbh_paper,tsoumakas_multilabel_overview}.
  \item \textbf{Robustez al desbalance:} MCC como métrica complementaria a accuracy/F1, ya que resume de forma más equilibrada las cuatro celdas de la matriz de confusión \cite{chicco_mcc}.
  \item \textbf{Costo temporal del protocolo:} se reportan explícitamente el \textbf{tiempo de optimización de parámetros}, el \textbf{tiempo de construcción del modelo} y el \textbf{tiempo de aplicación del modelo}. El primero corresponde al tiempo total del bloque de búsqueda de hiperparámetros; el segundo, al tiempo medido alrededor de \texttt{fit} y exportado como \texttt{train\_time\_seconds}; y el tercero, al tiempo total del benchmark operativo, exportado como \texttt{wall\_time\_seconds} y complementado con latencias por request.
  \item \textbf{Eficiencia/tiempo real:} latencia media y percentiles P50/P95/P99 del pipeline (extracción de features + inferencia), throughput (req/s), tasa de fallos, CPU/RAM, tamaño del modelo y desagregación entre extracción e inferencia. Estas métricas se incluyen porque la evaluación no solo considera efectividad, sino también viabilidad operativa en un entorno tipo WAF \cite{ocswaf_zenodo}.
\end{itemize}

El detalle formal de cada métrica y de su cálculo exacto se encuentra en la sección de fundamentos matemáticos de la fundamentación teórica.

\section{Criterios de comparación con la literatura}
Comparar resultados entre papers puede ser engañoso si difieren \textit{splits}, preprocesamiento o definición de etiquetas. En este trabajo, la comparación se reduce al dominio discutido: detección y clasificación de ataques en tráfico HTTP, especialmente en escenarios cercanos a WAF o a datasets de ataques web. Aun dentro de ese dominio, los resultados de la literatura se usan como referencia contextual y no como ranking absoluto. Para maximizar comparabilidad:
\begin{itemize}[leftmargin=*]
  \item se reporta el \textbf{mismo conjunto de métricas} usado con frecuencia en la literatura;
  \item cuando un paper usa el mismo dataset y tarea, se replica la tarea y se reportan métricas homólogas;
  \item cuando un paper usa una configuración distinta, se compara solo a nivel de \textbf{métrica comparable} y se documenta explícitamente la amenaza a validez.
\end{itemize}

\begin{table}[H]
\centering
\caption{Métricas reportadas en trabajos relacionados y su correspondencia con el estudio.}
\label{tab:metricas_relacionadas}
\footnotesize
\setlength{\tabcolsep}{5pt}
\renewcommand{\arraystretch}{1.2}
\begin{tabularx}{\textwidth}{@{}L{3.6cm}Y@{}}
\toprule
\textbf{Trabajo} & \textbf{Métricas típicas reportadas} \\
\midrule
WAF basado en detección de anomalías & TPR/FPR, F1 y medición de overhead en tiempo de respuesta (ms) \cite{ocswaf_zenodo}. \\
SR-BH (paper del dataset) & Multietiqueta: Hamming Loss, Jaccard, EMR, F1 y AUC \cite{srbh_paper}. \\
Selección de features en web attacks (Riverol 2024) & Accuracy, TPR, FPR y AUC/ROC (curvas ROC) \cite{riverol_2024}. \\
Modelos recientes con énfasis en IDS (Fang 2024) & Accuracy, Precision, Recall, F1, FPR/FNR y MCC \cite{fang_2024,chicco_mcc}. \\
Ensamble/Deep Kernel Learning (Zhou 2024) & Accuracy, macro Precision/Recall/F1 y weighted Precision/F1 (desbalance) \cite{zhou_2024}. \\
\bottomrule
\end{tabularx}
\end{table}


\section{Importancia de características (FI)}
Esta sección se incluye porque la comparación no se limita a contrastar métricas finales entre modelos. El diseño experimental también requiere examinar el aporte relativo de los grupos de características utilizados y discutir hasta qué punto el rendimiento observado proviene de señales estructurales del request HTTP o de indicadores léxicos más cercanos a firmas. Por esa razón, además de medir desempeño, el protocolo incorpora un bloque específico para analizar \textbf{qué variables empujan la decisión del modelo} y \textbf{qué significado metodológico tiene ese aporte}.

En este contexto, \textbf{interpretabilidad} se entiende como la posibilidad de relacionar el comportamiento del modelo con variables concretas de entrada. En modelos lineales, esto significa poder revisar qué features reciben mayor peso y en qué dirección contribuyen a la predicción. \textbf{Explicabilidad}, en un sentido más amplio, refiere a la capacidad de justificar por qué un modelo rinde como rinde y qué tipo de señal está aprovechando. Aunque ambos términos suelen usarse juntos, aquí conviene distinguirlos: la interpretabilidad apunta a leer el modelo; la explicabilidad apunta a sostener una discusión metodológica sobre sus decisiones y sus límites.

La sigla \textbf{FI} (\textit{Feature Importance}) se usa como nombre corto para el análisis de importancia de variables. Aquí no se plantea como una etapa separada y costosa, sino como una lectura complementaria del mismo protocolo experimental. La idea no es construir un ranking absoluto y definitivo de variables, sino observar si ciertos grupos ---por ejemplo, longitudes, estructura de URI, método HTTP o tokens sospechosos--- sostienen una parte importante del rendimiento.

La estrategia adoptada tiene tres niveles:
\begin{itemize}[leftmargin=*]
  \item \textbf{Coeficientes estandarizados} (LR/OvR): como los modelos supervisados son lineales y se entrenan sobre variables escaladas, la magnitud de los coeficientes permite comparar de forma razonable el peso relativo de cada feature. En multiclase, la lectura se hace por clase; en multietiqueta, por etiqueta.
  \item \textbf{Ablaciones por grupos}: se comparan corridas del pipeline completo contra variantes en las que se elimina un grupo entero de variables. Esto permite discutir el aporte del grupo como bloque y no solo de features aisladas.
  \item \textbf{Ablación focal con y sin tokens sospechosos}: se reserva una comparación específica para ese grupo porque es el más delicado metodológicamente. Si gran parte del rendimiento dependiera de esas variables, la representación quedaría más cerca de una lógica de firma que de una descripción general del tráfico HTTP.
\end{itemize}

La sección está, entonces, para sostener una pregunta central del estudio: no solo cuál modelo rinde mejor, sino también \textbf{por qué} rinde mejor y \textbf{sobre qué tipo de señal} lo hace. Esa discusión resulta importante para interpretar resultados, comparar con la literatura y valorar la factibilidad de usar la representación adoptada en un contexto tipo WAF sin convertir el pipeline en un conjunto encubierto de reglas manuales.

\section{Amenazas a la validez}
\begin{itemize}[leftmargin=*]
  \item \textbf{Desbalance}: etiquetas CAPEC pueden ser raras; afecta F1 macro, estratificación y estabilidad del entrenamiento.
  \item \textbf{Reducción multilabel$\rightarrow$multiclase}: seleccionar una única etiqueta primaria simplifica la tarea, pero colapsa información multietiqueta \cite{capec_schema_docs,capec_schema_desc,owasp_risk_rating}.
  \item \textbf{Evaluación one-class}: si el conjunto mixto de prueba se reutiliza para tuning, las métricas se inflan; por eso el protocolo separa la validación interna basada en normales del conjunto final de prueba compuesto por holdout de normales y todas las anomalías/ataques disponibles.
  \item \textbf{Heurísticas de tokens}: pueden sesgar los modelos hacia firmas; por eso se exige la ablación correspondiente.
  \item \textbf{Domain shift}: datasets sintéticos/curados (CSIC/TorpEda) pueden no transferir a tráfico real (SR-BH).
  \item \textbf{Validez externa}: aun en SR-BH, el tráfico proviene de un honeypot específico y no de todas las configuraciones de producción posibles.
\end{itemize}

\chapter{Resultados y discusión}

\section{Criterio de presentación}
El capítulo de resultados adopta una organización estrictamente tabular por tarea, dataset y variante de modelo. La decisión responde a dos necesidades simultáneas. Por un lado, resulta necesario comparar enfoques \textit{one-class}, multiclase y multietiqueta sobre CSIC 2010, CSIC/TorpEda 2012 y SR-BH 2020, con métricas adecuadas a cada formulación y con una discusión explícita de factibilidad tipo WAF. Por otro lado, el propio protocolo metodológico del estudio establece que la comparación debe mantener separados entrenamiento, ajuste y prueba, y que además del rendimiento predictivo deben informarse tiempos, latencias, \textit{throughput}, consumo de recursos y tamaño del modelo.\

Para preservar comparabilidad entre artefactos heterogéneos, este capítulo distingue dos familias de tablas. Las primeras reúnen las métricas de efectividad reportadas por cada corrida. Las segundas reúnen las métricas operativas. En los bloques operativos, el tiempo de aplicación $t_{apply}^{\dagger}$ se expresa como el tiempo aproximado del lote principal de benchmark, calculado como \emph{solicitudes medidas / throughput global}; esta definición permite homogeneizar artefactos que no siempre exportan el mismo nombre de campo para el tiempo de pared del benchmark, pero sí informan de manera estable el volumen medido y el \textit{throughput}.\

\section{Lectura metodológica de las corridas}
Las corridas incorporadas en este libro siguen una lógica metodológica común: una representación compartida de peticiones HTTP, tres datasets con distinta granularidad de etiquetado y modelos compatibles con cada formulación del problema. CSIC 2010 se conserva como referencia controlada, útil para contrastar normalidad frente a tráfico anómalo y para observar qué ocurre cuando una tarea declarada como multiclase o multietiqueta se reduce en la práctica a un problema mucho más simple. CSIC/TorpEda 2012 aporta tipologías explícitas de ataque y un escenario genuinamente multiclase. SR-BH 2020, derivado de tráfico real de \textit{honeypot}, conserva el valor central de aportar mayor realismo y una semántica CAPEC más rica que la de los datasets puramente sintéticos o de laboratorio \cite{srbh_dataverse,srbh_paper,csic2010_impact,csic2010_doc,torpeda_2012}.\

La selección de modelos también sigue esa lógica de adecuación a la tarea y de costo computacional defendible. Para detección binaria de anomalías se utilizó un esquema \texttt{StandardScaler + Nystroem + SGDOneClassSVM}, porque permite aproximar una frontera no lineal de normalidad con costo mucho menor que el de un kernel exacto, manteniendo el supuesto central de entrenar solo con tráfico legítimo \cite{scholkopf2001,williams2001,sklearn_sgd_ocsvm,sklearn_kernel_approx}. Para multiclase y multietiqueta se adoptó regresión logística lineal, y en el caso multietiqueta se la desplegó en esquema One-vs-Rest. La elección es coherente con los objetivos del estudio: interpretabilidad, trazabilidad, bajo costo de inferencia y una línea base suficientemente fuerte para comparar formulaciones sin que la complejidad del modelo eclipse la discusión experimental \cite{islr_logreg,sklearn_logreg,rifkin_ova,tsoumakas_multilabel_overview}.\

El ajuste de hiperparámetros se mantuvo acotado. En los modelos supervisados se usó búsqueda aleatoria con presupuestos muy pequeños y validación cruzada interna, mientras que en el detector one-class se restringió la exploración a parámetros que controlan la forma de la frontera, la aproximación del kernel y la convergencia. No se presenta una optimización exhaustiva, sino una sintonización ligera y reproducible, suficiente para evitar configuraciones arbitrarias sin convertir el tuning en el costo dominante del protocolo \cite{bergstra2012,sklearn_randomsearch}.\

La ablación de variables basadas en \textit{suspicious tokens} se mantuvo en los tres bloques porque cumple una función metodológica importante: separar el aporte de señales estructurales y de codificación de aquellas señales que se aproximan más a firmas léxicas explícitas. Cuando la ablación apenas modifica los resultados, puede sostenerse que el modelo depende poco de esos atajos. Cuando la ablación degrada fuertemente el rendimiento, la interpretación correcta no es que el modelo sea ``malo, sino que parte de su capacidad discriminativa descansa efectivamente en conocimiento léxico de dominio, algo que conviene declarar para no sobredimensionar la supuesta neutralidad de la representación.\

\section{Detección one-class}
\subsection{Métricas de efectividad}
\begin{table}[H]
\centering
\caption{Resultados de detección one-class por dataset y variante.}
\label{tab:res-one-eff}
\scriptsize
\renewcommand{\arraystretch}{1.12}
\begin{adjustbox}{max width=\textwidth}
\begin{tabular}{llccccccccccc}
\toprule
\textbf{Dataset} & \textbf{Variante} & \textbf{Acc.} & \textbf{Bal. Acc.} & \textbf{Prec.} & \textbf{Recall} & \textbf{F1} & \textbf{MCC} & \textbf{TNR} & \textbf{FPR} & \textbf{FNR} & \textbf{ROC AUC} & \textbf{PR AUC} \\
\midrule
CSIC 2010 & Base & 0.6951 & 0.7512 & 0.9581 & 0.5437 & 0.6938 & 0.5037 & 0.9586 & 0.0414 & 0.4563 & 0.7688 & 0.8801 \\CSIC 2010 & Sin tokens & 0.7013 & 0.7631 & 0.9916 & 0.5341 & 0.6943 & 0.5340 & 0.9922 & 0.0078 & 0.4659 & 0.7630 & 0.8762 \\SR-BH 2020 & Base & 0.3175 & 0.5617 & 0.9813 & 0.1327 & 0.2338 & 0.1648 & 0.9908 & 0.0092 & 0.8673 & 0.5158 & 0.8547 \\SR-BH 2020 & Sin tokens & 0.3124 & 0.5585 & 0.9804 & 0.1261 & 0.2235 & 0.1596 & 0.9908 & 0.0092 & 0.8739 & 0.4736 & 0.8276 \\CSIC/TorpEda 2012 & Base & 0.7440 & 0.8615 & 0.9995 & 0.7379 & 0.8490 & 0.2504 & 0.9851 & 0.0149 & 0.2621 & 0.8358 & 0.9954 \\CSIC/TorpEda 2012 & Sin tokens & 0.7369 & 0.7661 & 0.9930 & 0.7354 & 0.8450 & 0.1848 & 0.7968 & 0.2032 & 0.2646 & 0.8215 & 0.9949 \\\bottomrule
\end{tabular}
\end{adjustbox}
\end{table}

\subsection{Métricas operativas}
\begin{table}[H]
\centering
\caption{Indicadores operativos de las corridas one-class.}
\label{tab:res-one-op}
\scriptsize
\renewcommand{\arraystretch}{1.12}
\begin{adjustbox}{max width=\textwidth}
\begin{tabular}{llcccccccc}
\toprule
\textbf{Dataset} & \textbf{Variante} & $t_{opt}$ (s) & $t_{build}$ (s) & $t_{apply}^{\dagger}$ (s) & \textbf{Thr.} & \textbf{P95} (ms) & \textbf{CPU} (\%) & \textbf{Peak RSS} (MB) & \textbf{Modelo} (KiB) \\
\midrule
CSIC 2010 & Base & 4.120 & 1.071 & 2.076 & 192.708 & 9.916 & 130.56 & 873.45 & 569.5 \\CSIC 2010 & Sin tokens & 4.264 & 1.161 & 2.429 & 164.639 & 10.697 & 144.47 & 861.43 & 561.2 \\SR-BH 2020 & Base & 3.859 & 5.270 & 1.629 & 245.480 & 6.032 & 123.97 & 5448.73 & 569.5 \\SR-BH 2020 & Sin tokens & 5.142 & 6.797 & 1.648 & 242.650 & 5.458 & 124.36 & 5429.17 & 561.2 \\CSIC/TorpEda 2012 & Base & 1.330 & 0.380 & 1.611 & 248.260 & 5.135 & 146.47 & 942.27 & 569.5 \\CSIC/TorpEda 2012 & Sin tokens & 1.238 & 0.249 & 1.469 & 272.250 & 4.518 & 159.26 & 902.93 & 561.2 \\\bottomrule
\end{tabular}
\end{adjustbox}
\end{table}

En conjunto, las corridas one-class muestran un patrón estable. La variante base conserva la mejor lectura global en Harvard y en TorpEda, especialmente por su mejor \textit{balanced accuracy}, MCC y control del error sobre la clase normal. En CSIC 2010, en cambio, la ablación mejora levemente Accuracy, Balanced Accuracy y MCC, aunque lo hace a costa de una caída pequeña en ROC AUC y PR AUC. Estos resultados indican que el bloque one-class se beneficia de una representación no lineal suficientemente rica como para no depender de manera uniforme de los tokens sospechosos.\

Operativamente, el costo de entrenamiento final del detector es bajo en los tres datasets, pero el costo de memoria no lo es. Harvard es, con mucha diferencia, el caso más exigente: supera los 5.4 GB de pico de RSS en ambas variantes y, aunque mantiene \textit{throughput} alto, lo hace con una huella de memoria muy superior al resto. TorpEda ofrece la mejor combinación entre discriminación y costo. CSIC queda en una posición intermedia, con resultados predictivos razonables y tiempos compatibles con una discusión de viabilidad, aunque sin la solidez de TorpEda en separación binaria.

\section{Clasificación multiclase}
\subsection{Métricas de efectividad}
\begin{table}[H]
\centering
\caption{Resultados multiclase por dataset y variante. En CSIC 2010 las AUC corresponden a la formulación binaria efectiva observada en prueba; en Harvard y TorpEda corresponden a la lectura OvR ponderada reportada por los artefactos.}
\label{tab:res-mc-eff}
\scriptsize
\renewcommand{\arraystretch}{1.12}
\begin{adjustbox}{max width=\textwidth}
\begin{tabular}{llccccccccc}
\toprule
\textbf{Dataset} & \textbf{Variante} & \textbf{Acc.} & \textbf{Bal. Acc.} & \textbf{F1 macro} & \textbf{F1 weighted} & \textbf{Prec. macro} & \textbf{Recall macro} & \textbf{MCC} & \textbf{ROC AUC} & \textbf{PR AUC} \\
\midrule
CSIC 2010 & Base & 0.8286 & 0.6989 & 0.7291 & 0.8085 & 0.8248 & 0.6989 & 0.5084 & 0.8430 & 0.9229 \\CSIC 2010 & Sin tokens & 0.8294 & 0.7011 & 0.7313 & 0.8098 & 0.8249 & 0.7011 & 0.5112 & 0.8439 & 0.9226 \\SR-BH 2020 & Base & 0.8014 & 0.5039 & 0.5252 & 0.7854 & 0.6240 & 0.5039 & 0.6450 & 0.9363 & 0.8631 \\SR-BH 2020 & Sin tokens & 0.7173 & 0.3426 & 0.3837 & 0.6882 & 0.5388 & 0.3426 & 0.4875 & 0.8958 & 0.8003 \\CSIC/TorpEda 2012 & Base & 0.8408 & 0.2916 & 0.3011 & 0.8092 & 0.5303 & 0.2916 & 0.7384 & 0.9675 & 0.8499 \\CSIC/TorpEda 2012 & Sin tokens & 0.8136 & 0.2556 & 0.2531 & 0.7747 & 0.3251 & 0.2556 & 0.6900 & 0.9569 & 0.8114 \\\bottomrule
\end{tabular}
\end{adjustbox}
\end{table}

\subsection{Métricas operativas}
\begin{table}[H]
\centering
\caption{Indicadores operativos de las corridas multiclase.}
\label{tab:res-mc-op}
\scriptsize
\renewcommand{\arraystretch}{1.12}
\begin{adjustbox}{max width=\textwidth}
\begin{tabular}{llcccccccc}
\toprule
\textbf{Dataset} & \textbf{Variante} & $t_{opt}$ (s) & $t_{build}$ (s) & $t_{apply}^{\dagger}$ (s) & \textbf{Thr.} & \textbf{P95} (ms) & \textbf{CPU} (\%) & \textbf{Peak RSS} (MB) & \textbf{Modelo} (KiB) \\
\midrule
CSIC 2010 & Base & 2.4525 & 1.0872 & 1.5929 & 376.6683 & 3.2893 & 114.26 & 479.48 & 3.42 \\CSIC 2010 & Sin tokens & 2.3335 & 1.0055 & 1.4508 & 413.5759 & 2.6209 & 118.56 & 467.23 & 3.08 \\SR-BH 2020 & Base & 10.0885 & 137.3066 & 1.5420 & 388.8543 & 2.9597 & 119.90 & 2380.93 & 6.87 \\SR-BH 2020 & Sin tokens & 9.5461 & 144.5713 & 1.5916 & 376.9847 & 3.1997 & 118.12 & 2269.09 & 6.12 \\CSIC/TorpEda 2012 & Base & 3.8222 & 14.6148 & 1.4965 & 400.9321 & 2.8933 & 118.28 & 478.72 & 5.51 \\CSIC/TorpEda 2012 & Sin tokens & 3.7528 & 12.0153 & 1.5008 & 399.7902 & 2.8655 & 116.61 & 472.93 & 4.88 \\\bottomrule
\end{tabular}
\end{adjustbox}
\end{table}

Las tres corridas multiclase dejan una señal metodológica nítida: el efecto de la ablación depende fuertemente del dataset. En CSIC 2010, donde la prueba termina comportándose como una discriminación de dos etiquetas efectivas, la diferencia entre variantes es mínima y la versión sin tokens obtiene una mejora marginal en Accuracy, Balanced Accuracy, F1 macro y MCC. En Harvard la situación es la opuesta: la ablación degrada todas las métricas globales de manera marcada, lo que sugiere que las señales léxicas aportan información importante para separar varias tipologías de ataque sobre un espacio de clases amplio y heterogéneo. En TorpEda el resultado también favorece al modelo base, sobre todo en macro F1, MCC y AUC.\

La lectura por clases, analizada en los artefactos específicos de cada corrida, muestra además que la degradación no se distribuye de forma pareja. En Harvard, por ejemplo, el impacto negativo de quitar tokens sospechosos es especialmente visible en clases como CAPEC-242, CAPEC-34 y CAPEC-194, mientras que otras categorías prácticamente no cambian. Eso refuerza una conclusión importante para el libro: la utilidad de las señales léxicas no es uniforme ni universal, pero tampoco puede tratarse como un simple adorno.\

En el plano operativo, multiclase es el bloque más liviano en inferencia para CSIC y TorpEda, con latencias P95 por debajo de 3.3 ms y \textit{throughput} cercano o superior a 400 req/s. Harvard mantiene una inferencia rápida, pero el costo de construcción del modelo es muy superior al de los otros datasets y su pico de memoria durante benchmark también crece con fuerza. Por lo tanto, la discusión de viabilidad debe separar claramente costo de entrenamiento y costo de aplicación: Harvard multiclase es operativo una vez entrenado, pero no es el escenario más cómodo para reentrenamientos frecuentes.

\section{Clasificación multietiqueta}
\subsection{Métricas globales}
\begin{table}[H]
\centering
\caption{Métricas globales multietiqueta por dataset y variante.}
\label{tab:res-ml-eff}
\scriptsize
\renewcommand{\arraystretch}{1.12}
\begin{adjustbox}{max width=\textwidth}
\begin{tabular}{llcccccccc}
\toprule
\textbf{Dataset} & \textbf{Variante} & \textbf{F1 micro} & \textbf{F1 macro} & \textbf{Hamming} & \textbf{Jaccard micro} & \textbf{EMR} & $t_{build}$ (s) & $t_{apply}^{\dagger}$ (s) & \textbf{Thr.} \\
\midrule
CSIC 2010 & Base & 0.8336 & 0.7361 & 0.1664 & 0.7147 & 0.8336 & 2.4590 & 2.0293 & 295.67 \\CSIC 2010 & Sin tokens & 0.7771 & 0.7338 & 0.2229 & 0.6355 & 0.7771 & 1.9575 & 1.9460 & 308.33 \\SR-BH 2020 & Base & 0.6582 & 0.1079 & 0.0039 & 0.4906 & 0.7820 & 19.96 & 6.57 & 91.32 \\SR-BH 2020 & Sin tokens & 0.2052 & 0.1013 & 0.0508 & 0.1143 & 0.0723 & 37.69 & 6.36 & 94.29 \\CSIC/TorpEda 2012 & Base & 0.8498 & 0.2601 & 0.0288 & 0.7388 & 0.7681 & 2.6578 & 2.2321 & 268.84 \\CSIC/TorpEda 2012 & Sin tokens & 0.4725 & 0.2393 & 0.2085 & 0.3094 & 0.3902 & 1.5397 & 2.3763 & 252.49 \\\bottomrule
\end{tabular}
\end{adjustbox}
\end{table}

\begin{table}[H]
\centering
\caption{Métricas extendidas disponibles para las corridas multietiqueta de Harvard y TorpEda.}
\label{tab:res-ml-extra}
\scriptsize
\renewcommand{\arraystretch}{1.12}
\begin{adjustbox}{max width=\textwidth}
\begin{tabular}{llcccccc}
\toprule
\textbf{Dataset} & \textbf{Variante} & \textbf{Prec. micro} & \textbf{Recall micro} & \textbf{PR AUC micro} & \textbf{PR AUC macro} & \textbf{ROC AUC micro} & \textbf{ROC AUC macro} \\
\midrule
SR-BH 2020 & Base & 0.9376 & 0.5071 & 0.7297 & 0.1383 & 0.9903 & 0.7075 \\SR-BH 2020 & Sin tokens & 0.1160 & 0.8886 & 0.1271 & 0.1269 & 0.9718 & 0.9571 \\CSIC/TorpEda 2012 & Base & 0.8746 & 0.8263 & 0.9244 & --- & 0.9864 & --- \\CSIC/TorpEda 2012 & Sin tokens & 0.3148 & 0.9473 & 0.7721 & --- & 0.9555 & --- \\\bottomrule
\end{tabular}
\end{adjustbox}
\end{table}

El Cuadro~\ref{tab:res-ml-extra} requiere una aclaración explícita. En SR-BH~2020, la variante base conserva la mejor combinación entre precisión micro, \textit{PR AUC} micro y \textit{ROC AUC} micro, mientras que la ablación sin tokens sospechosos empuja el \textit{recall} micro a valores muy altos a costa de una caída extrema de precisión. Esa combinación es típica de una sobrepredicción agresiva: el modelo recupera más etiquetas positivas, pero también dispara muchos más falsos positivos. El hecho de que \textit{ROC AUC} macro resulte mayor en la ablación no contradice esa lectura, porque AUC resume capacidad de ordenamiento sobre distintos umbrales, mientras que precisión, \textit{exact match}, Jaccard y $F_1$ dependen de la decisión binaria finalmente aplicada en inferencia.

En CSIC/TorpEda~2012, los símbolos --- de \textit{PR AUC} macro y \textit{ROC AUC} macro no representan rendimiento nulo ni error de cálculo. Responden a una decisión metodológica ligada a la forma en que esas corridas fueron formuladas. El script multietiqueta incorpora un modo de compatibilidad \texttt{reduced-binary-mode} para datasets cuya columna multietiqueta queda efectivamente reducida a una única dimensión positiva; en ese escenario el problema deja de comportarse como un multietiqueta rico y pasa a interpretarse como un caso binario reducido, con bloque específico de métricas binarias y con menor utilidad real de los promedios entre etiquetas. En los comandos utilizados, esa lógica se activa de manera forzada para CSIC y se habilita en modo automático para TorpEda. Bajo esa formulación, los promedios macro de \textit{PR AUC} y \textit{ROC AUC} dejan de aportar una lectura sustantivamente distinta: con una sola dimensión efectiva, el promedio macro colapsa hacia la misma señal de ordenamiento que ya queda representada por la evaluación binaria del objetivo activo y por las métricas micro conservadas en el cuadro. Por esa razón se prefirió marcar esas celdas con --- en lugar de exhibirlas como si fueran evidencia adicional de un comportamiento multietiqueta pleno, algo que habría sobreinterpretado la estructura real de esas corridas.

\begin{table}[H]
\centering
\caption{Métricas binarias derivadas para CSIC 2010 en su formulación multietiqueta reducida a binario.}
\label{tab:res-ml-csic-bin}
\scriptsize
\renewcommand{\arraystretch}{1.12}
\begin{adjustbox}{max width=\textwidth}
\begin{tabular}{lccccccccc}
\toprule
\textbf{Variante} & \textbf{Acc.} & \textbf{Bal. Acc.} & \textbf{Prec.+} & \textbf{Recall+} & \textbf{F1+} & \textbf{MCC} & \textbf{TNR} & \textbf{FPR} & \textbf{FNR} \\
\midrule
Base & 0.8336 & 0.7046 & 0.8353 & 0.4392 & 0.5757 & 0.5231 & 0.9700 & 0.0300 & 0.5608 \\Sin tokens & 0.7771 & 0.7608 & 0.5502 & 0.7272 & 0.6264 & 0.4812 & 0.7944 & 0.2056 & 0.2728 \\\bottomrule
\end{tabular}
\end{adjustbox}
\end{table}

\begin{table}[H]
\centering
\caption{Indicadores operativos consolidados de las corridas multietiqueta.}
\label{tab:res-ml-op}
\scriptsize
\renewcommand{\arraystretch}{1.12}
\begin{adjustbox}{max width=\textwidth}
\begin{tabular}{llcccccc}
\toprule
\textbf{Dataset} & \textbf{Variante} & \textbf{P95} (ms) & \textbf{CPU} (\%) & \textbf{Peak RSS} (MB) & \textbf{Modelo} (KiB) & \textbf{Tasa de solicitudes fallidas} & \textbf{Lectura} \\
\midrule
CSIC 2010 & Base & 4.2743 & 132.07 & 445.46 & 4.17 & 0.0000 & mejor efectividad global \\CSIC 2010 & Sin tokens & 4.2951 & 122.30 & 445.56 & 3.85 & 0.0000 & más rápido bajo carga leve \\SR-BH 2020 & Base & 13.60 & 122.52 & 2862.48 & 50.76 & 0.0000 & mejor equilibrio semántico \\SR-BH 2020 & Sin tokens & 12.77 & 121.79 & 2920.11 & 49.59 & 0.0000 & alto recall, muy baja precisión \\CSIC/TorpEda 2012 & Base & 3.9211 & 124.56 & 436.51 & 9.95 & 0.0000 & baseline claramente superior \\CSIC/TorpEda 2012 & Sin tokens & 4.8011 & 122.88 & 437.49 & 9.53 & 0.0000 & predicción excesivamente expansiva \\\bottomrule
\end{tabular}
\end{adjustbox}
\end{table}

La lectura del bloque multietiqueta requiere más cautela que la de los otros dos. En CSIC 2010, el propio artefacto declara una tarea efectiva \texttt{binary\_reduced}; en otras palabras, el pipeline multietiqueta se conserva por consistencia metodológica, pero el problema observado en prueba ya no es genuinamente multietiqueta. En ese contexto, el modelo base gana con claridad en F1 micro, Jaccard micro y exact match, aunque la ablación obtenga mayor recall positivo y mejor balanced accuracy en la lectura binaria derivada. La diferencia entre ambas lecturas no es una contradicción: una cosa es maximizar coincidencia global bajo predicción más conservadora y otra maximizar recuperación de la clase positiva aceptando muchos más falsos positivos.\

En SR-BH 2020 aparece el caso más exigente del estudio. Allí la variante base domina ampliamente en F1 micro, exact match, Jaccard y precisión micro, mientras que la ablación solo conserva como aparente ventaja un recall micro muy alto acompañado por una caída drástica de precisión. El patrón es claro: sin tokens sospechosos, el modelo pasa a sobrepredecir etiquetas con una agresividad que destruye la coincidencia exacta y el solapamiento útil entre conjuntos reales y predichos.\

TorpEda 2012 muestra un comportamiento similar. Aunque la tarea se declara como multietiqueta, el mismo artefacto advierte que la cardinalidad real observada en prueba equivale en la práctica a una sola etiqueta o instancia vacía. Aun así, la comparación sigue siendo informativa: el baseline mantiene la mejor combinación entre F1 micro, Jaccard, exact match y AUC, mientras que la ablación incrementa fuertemente la cardinalidad predicha y la proporción de instancias clasificadas como multietiqueta, produciendo una salida excesivamente expansiva.\

\section{Síntesis transversal}
La comparación completa confirma que ninguna formulación domina de manera universal a las otras. El bloque one-class cumple mejor cuando la pregunta principal es si una solicitud se aparta o no del soporte de normalidad. El bloque multiclase es el más eficiente para tipificar cuando las clases están suficientemente definidas y la inferencia debe permanecer liviana. El bloque multietiqueta conserva su mayor valor analítico en SR-BH 2020, donde la riqueza semántica del problema justifica el costo adicional y donde la simplificación a una sola etiqueta empobrecería la lectura del fenómeno.\

También se observa que la utilidad de las variables de tokens sospechosos depende del escenario. En Harvard multiclase y multietiqueta, y en TorpEda multiclase y multietiqueta, esas variables aportan una fracción visible del rendimiento. En cambio, en CSIC multiclase y en parte del bloque one-class la ablación genera cambios menores o incluso mejoras marginales. Por eso no corresponde defender una postura absoluta; lo defendible es mostrar la evidencia de ambas variantes y discutir cuándo esas señales funcionan como atajos demasiado específicos y cuándo constituyen información realmente útil.\

Por último, el costo operativo vuelve a separar tareas. Multiclase es, en general, la formulación más liviana durante la aplicación. One-class conserva tiempos de entrenamiento muy bajos, pero con una huella de memoria mucho más alta. Multietiqueta es la formulación más costosa en SR-BH 2020, aunque también la más fiel a la estructura semántica del dataset. El resultado global del capítulo, por tanto, no es la proclamación de un único mejor modelo, sino la delimitación bastante precisa de qué enfoque conviene según el tipo de etiqueta disponible, el costo operativo tolerable y el grado de riqueza semántica que se desea preservar.\

\section{Limitaciones de las corridas}
La primera limitación relevante es de soporte de clases. En Harvard multiclase, la clase CAPEC-248 Command Injection no queda representada en el conjunto de prueba y, además, el dataset completo solo aporta un registro de esa categoría. En términos prácticos, esto impide evaluar esa clase con estabilidad estadística y obliga a leer con cautela cualquier promedio macro que dependa de etiquetas declaradas pero ausentes en test. El problema no invalida la corrida, pero sí limita la interpretación de la cobertura multiclase sobre el espacio completo de CAPEC.\

La segunda limitación es de granularidad efectiva. CSIC 2010 en configuración multiclase termina comportándose como un problema de dos etiquetas efectivas en prueba; CSIC 2010 y TorpEda 2012 en configuración multietiqueta se reducen de hecho a escenarios casi monolabel. Estas corridas siguen siendo útiles para sostener la matriz experimental completa y para contrastar el comportamiento del pipeline bajo tareas formalmente distintas, pero no deben presentarse como evidencia de multietiqueta genuina en el mismo sentido que SR-BH 2020.\

La tercera limitación es el alcance del tuning. Todas las corridas utilizaron presupuestos pequeños. Eso es coherente con el carácter reproducible y acotado del estudio, pero no autoriza a afirmar que se halló el óptimo global de cada familia de modelos. Los resultados son comparables dentro del protocolo definido; no constituyen, sin embargo, una búsqueda exhaustiva del máximo desempeño posible.\

La cuarta limitación es de entorno operativo. Los benchmarks informan latencia, throughput, CPU, memoria, tamaño del modelo y tasa de fallos sobre un pipeline local de extracción de características más inferencia. Esa evidencia es útil para estimar factibilidad, pero no sustituye un despliegue real sobre un proxy reverso, con concurrencia sostenida, tráfico de red, reglas de decisión y costos de integración propios de un WAF completo.\


\section{Protocolo de evaluación de viabilidad en tiempo real en un entorno tipo WAF}

\subsection{Requisitos operativos}
El alcance no incluye el despliegue de un WAF completo, pero sí exige justificar si los modelos estudiados podrían integrarse, al menos teóricamente, dentro de un entorno de inspección HTTP con restricciones de latencia y de recursos. Desde esa perspectiva, los requisitos mínimos son cuatro: (i) mantener una latencia baja y estable en el pipeline de extracción e inferencia; (ii) sostener un \textit{throughput} razonable bajo carga liviana; (iii) evitar picos de memoria incompatibles con un despliegue práctico; y (iv) conservar artefactos de tamaño acotado, de manera que el costo de persistencia y carga del modelo no domine el sistema.\

\subsection{Consolidado operativo por dataset y modelo}
\begin{table}[H]
\centering
\caption{Resumen consolidado de viabilidad operativa.}
\label{tab:waf-consolidado}
\scriptsize
\renewcommand{\arraystretch}{1.10}
\begin{adjustbox}{max width=\textwidth}
\begin{tabular}{lllcccccccc}
\toprule
\textbf{Tarea} & \textbf{Dataset} & \textbf{Variante} & $t_{opt}$ & $t_{build}$ & $t_{apply}^{\dagger}$ & \textbf{Thr.} & \textbf{P95} & \textbf{CPU} & \textbf{Peak RSS} & \textbf{Modelo} \\
 &  &  & \textbf{(s)} & \textbf{(s)} & \textbf{(s)} & \textbf{(req/s)} & \textbf{(ms)} & \textbf{(\%)} & \textbf{(MB)} & \textbf{(KiB)} \\
\midrule
One-class & CSIC 2010 & Base & 4.120 & 1.071 & 2.076 & 192.708 & 9.916 & 130.56 & 873.45 & 569.5 \\One-class & CSIC 2010 & Sin tokens & 4.264 & 1.161 & 2.429 & 164.639 & 10.697 & 144.47 & 861.43 & 561.2 \\One-class & SR-BH 2020 & Base & 3.859 & 5.270 & 1.629 & 245.480 & 6.032 & 123.97 & 5448.73 & 569.5 \\One-class & SR-BH 2020 & Sin tokens & 5.142 & 6.797 & 1.648 & 242.650 & 5.458 & 124.36 & 5429.17 & 561.2 \\One-class & CSIC/TorpEda 2012 & Base & 1.330 & 0.380 & 1.611 & 248.260 & 5.135 & 146.47 & 942.27 & 569.5 \\One-class & CSIC/TorpEda 2012 & Sin tokens & 1.238 & 0.249 & 1.469 & 272.250 & 4.518 & 159.26 & 902.93 & 561.2 \\Multiclase & CSIC 2010 & Base & 2.4525 & 1.0872 & 1.5929 & 376.6683 & 3.2893 & 114.26 & 479.48 & 3.42 \\Multiclase & CSIC 2010 & Sin tokens & 2.3335 & 1.0055 & 1.4508 & 413.5759 & 2.6209 & 118.56 & 467.23 & 3.08 \\Multiclase & SR-BH 2020 & Base & 10.0885 & 137.3066 & 1.5420 & 388.8543 & 2.9597 & 119.90 & 2380.93 & 6.87 \\Multiclase & SR-BH 2020 & Sin tokens & 9.5461 & 144.5713 & 1.5916 & 376.9847 & 3.1997 & 118.12 & 2269.09 & 6.12 \\Multiclase & CSIC/TorpEda 2012 & Base & 3.8222 & 14.6148 & 1.4965 & 400.9321 & 2.8933 & 118.28 & 478.72 & 5.51 \\Multiclase & CSIC/TorpEda 2012 & Sin tokens & 3.7528 & 12.0153 & 1.5008 & 399.7902 & 2.8655 & 116.61 & 472.93 & 4.88 \\Multietiqueta & CSIC 2010 & Base & 0.3287 & 2.4590 & 2.0293 & 295.67 & 4.2743 & 132.07 & 445.46 & 4.17 \\Multietiqueta & CSIC 2010 & Sin tokens & 0.3287 & 1.9575 & 1.9460 & 308.33 & 4.2951 & 122.30 & 445.56 & 3.85 \\Multietiqueta & SR-BH 2020 & Base & 2.98 & 19.96 & 6.57 & 91.32 & 13.60 & 122.52 & 2862.48 & 50.76 \\Multietiqueta & SR-BH 2020 & Sin tokens & N/A & 37.69 & 6.36 & 94.29 & 12.77 & 121.79 & 2920.11 & 49.59 \\Multietiqueta & CSIC/TorpEda 2012 & Base & 0.7330 & 2.6578 & 2.2321 & 268.84 & 3.9211 & 124.56 & 436.51 & 9.95 \\Multietiqueta & CSIC/TorpEda 2012 & Sin tokens & 0.7330 & 1.5397 & 2.3763 & 252.49 & 4.8011 & 122.88 & 437.49 & 9.53 \\\bottomrule
\end{tabular}
\end{adjustbox}
\end{table}

\noindent\textbf{Aclaración de lectura.} En la fila \textit{Multietiqueta / SR-BH 2020 / Sin tokens}, la celda \texttt{N/A} en $t_{opt}$ indica que la corrida fue ejecutada sin búsqueda de hiperparámetros. No se trata de un tiempo de optimización igual a cero, sino de una etapa omitida por diseño. La implementación de la ablación multietiqueta ejecuta la variante sin tokens sospechosos con el ajuste de hiperparámetros desactivado salvo que se solicite explícitamente una reoptimización. En esa configuración, las diferencias observadas entre la variante base y la variante sin tokens quedan asociadas a la remoción de variables léxicas sospechosas y no a una segunda búsqueda que pudiera modificar el resultado por otra vía. Por esa razón, el valor correcto es \textit{no aplica}. Mantener esta distinción impide leer como eficiencia un bloque de \textit{tuning} que no fue ejecutado.

\subsection{Lectura de factibilidad}
El consolidado operativo permite distinguir con bastante claridad tres perfiles. El primero es el de las corridas multiclase sobre CSIC y TorpEda: alta velocidad, latencias muy bajas, artefactos pequeños y memoria moderada. Son las variantes más cómodas para una integración tipo WAF si la tarea de interés es tipificar una única clase por solicitud. El segundo perfil es el de las corridas one-class: su entrenamiento final es rápido y su discusión metodológica es atractiva cuando se quiere entrenar solo con tráfico normal, pero la memoria de proceso es considerablemente mayor, especialmente en Harvard. El tercer perfil es el de las corridas multietiqueta sobre SR-BH 2020: son las más costosas, pero también las que preservan mejor la estructura semántica del dataset más realista del estudio.\

Desde una lectura estrictamente operacional, el cuello de botella más severo del conjunto no está en la inferencia multiclase, sino en la memoria de one-class y en la latencia/throughput de multietiqueta sobre Harvard. Aun así, ninguna corrida muestra tasa de fallos distinta de cero en los benchmarks reportados, lo que permite sostener que, al menos dentro del alcance local del protocolo evaluado, no aparecen señales de inestabilidad catastrófica del pipeline.\

La consecuencia práctica es que la factibilidad teórica sí puede defenderse, pero de forma condicionada. Si el objetivo fuera un componente liviano y de baja latencia para clasificación de ataques con etiqueta única, las variantes multiclase resultan las candidatas más naturales. Si el objetivo fuera una alarma de anomalía entrenada solo con normales, el esquema one-class sigue siendo defendible, aunque con un costo de memoria que debe declararse. Si lo que se busca es conservar la semántica CAPEC multietiqueta de SR-BH 2020, entonces la formulación multietiqueta es la más fiel al problema, pero no la más barata de desplegar.\


\chapter{Conclusiones}

\section{Conclusiones}

Este Trabajo Final de Grado diseñó, implementó y evaluó un flujo reproducible de aprendizaje automático para la detección de anomalías y la clasificación de ataques en tráfico HTTP. El estudio comparó tres formulaciones complementarias del problema: detección \textit{one-class}, clasificación multiclase y clasificación multietiqueta. La comparación se realizó sobre tres fuentes de datos con distinta granularidad y realismo: CSIC~2010, CSIC/TorpEda~2012 y SR-BH~2020.

La primera conclusión es que la formulación del problema condiciona fuertemente la utilidad del modelo. El enfoque \textit{one-class} resulta metodológicamente adecuado cuando solo se dispone de tráfico normal para entrenamiento y se desea detectar desviaciones respecto de ese comportamiento. No obstante, su salida es limitada: permite marcar una petición como anómala, pero no explica por sí sola el tipo de ataque. La clasificación multiclase, en cambio, ofrece una tipificación directa cuando existe una única clase objetivo por instancia. La clasificación multietiqueta conserva mejor la riqueza semántica de SR-BH~2020, ya que una misma petición puede asociarse con más de un patrón CAPEC.

La segunda conclusión es que los modelos lineales evaluados constituyen líneas base defendibles para el alcance del trabajo. La regresión logística y su extensión One-vs-Rest proporcionan bajo costo de inferencia, artefactos pequeños y una interpretación más directa mediante coeficientes. Esto resulta relevante para un contexto tipo WAF, donde no solo importa el rendimiento predictivo, sino también la posibilidad de justificar qué señales del tráfico HTTP influyen en la predicción.

La tercera conclusión es que las variables basadas en tokens sospechosos aportan información útil en ciertos escenarios, pero deben interpretarse con cautela. En SR-BH~2020 y CSIC/TorpEda~2012, especialmente en tareas multiclase y multietiqueta, la eliminación de esas variables produce una degradación visible de varias métricas. Esto indica que los modelos aprovechan señales léxicas asociadas a ataques conocidos. Sin embargo, también muestra que parte del rendimiento puede acercarse a una lógica de firma, por lo que la comparación con y sin tokens resulta necesaria para evitar una interpretación exagerada de la generalización del modelo.

La cuarta conclusión es que la viabilidad operativa depende de la familia de modelo y del dataset. Las variantes multiclase presentan el perfil más liviano durante la aplicación, con baja latencia y alto rendimiento. Las variantes \textit{one-class} mantienen un entrenamiento final rápido, pero presentan mayor consumo de memoria, especialmente sobre SR-BH~2020. La formulación multietiqueta conserva la semántica más rica, aunque con mayor costo operativo en el dataset real de honeypot.

Finalmente, el trabajo confirma que no existe un único modelo dominante para todos los escenarios. La mejor elección depende del objetivo operativo: detección temprana de anomalías con pocos datos rotulados, tipificación de una clase principal de ataque o preservación de múltiples etiquetas semánticas. La contribución principal del estudio consiste en comparar esas alternativas bajo una representación común de peticiones HTTP, con separación explícita entre entrenamiento, ajuste y prueba, e incorporando tanto métricas predictivas como métricas de costo operativo.

\chapter{Trabajos futuros}

A partir de los resultados obtenidos, se identifican las siguientes líneas de continuidad:

\begin{itemize}[leftmargin=*]
  \item \textbf{Evaluación en un WAF o proxy real:} integrar el flujo de extracción de características e inferencia en un proxy reverso o componente WAF experimental, para medir latencia, concurrencia, estabilidad y costo de integración bajo tráfico HTTP real.
  \item \textbf{Ampliación de datasets:} incorporar nuevas fuentes de tráfico web, incluyendo tráfico de producción anonimizado o nuevos honeypots, para evaluar con mayor fuerza el cambio de dominio entre datasets sintéticos, benchmarks y tráfico real.
  \item \textbf{Estrategias avanzadas de umbral:} estudiar calibración de probabilidades y selección de umbrales por clase o etiqueta, especialmente en la formulación multietiqueta, donde un umbral global puede producir sobrepredicción o subpredicción de etiquetas.
  \item \textbf{Modelos no lineales y profundos:} comparar las líneas base utilizadas con modelos más complejos, como árboles de ensamble, redes neuronales o modelos híbridos, manteniendo siempre la medición de costo operativo para evitar mejoras predictivas impracticables.
  \item \textbf{Mejor análisis de explicabilidad:} complementar los coeficientes estandarizados con técnicas como importancia por permutación o explicaciones locales, con el objetivo de identificar qué características influyen en cada familia de ataque.
  \item \textbf{Tratamiento jerárquico de CAPEC:} aprovechar la estructura jerárquica de CAPEC para evaluar clasificación por niveles de abstracción, por ejemplo agrupando patrones específicos en familias superiores de ataque.
  \item \textbf{Gestión de clases raras:} aplicar técnicas de re-muestreo, agrupamiento, aprendizaje sensible al costo o evaluación específica para etiquetas con muy bajo soporte, especialmente en SR-BH~2020.
  \item \textbf{Ablaciones adicionales:} extender el análisis con ablaciones por grupos de características, separando longitud, estructura, codificación, método HTTP y tokens sospechosos, para cuantificar con mayor precisión el aporte de cada bloque.
\end{itemize}

% ============================================================================
% APÉNDICE
% ============================================================================
\appendix

\chapter{Comandos reproducibles}

Este apéndice resume una plantilla de comandos para reproducir el flujo experimental. Los nombres exactos de rutas pueden ajustarse según la organización final del repositorio, pero la secuencia refleja las etapas utilizadas: preparación del entorno, construcción de características, entrenamiento, evaluación y consolidación de resultados.

\section*{Preparación del entorno}

\begin{verbatim}
python -m venv .venv
source .venv/bin/activate
pip install -r requirements.txt
\end{verbatim}

\section*{Construcción de datasets procesados}

\begin{verbatim}
python build_features_csic.py \\
  --input data/raw/csic2010 \\
  --output data/processed/csic2010_features.parquet

python build_features_torpeda.py \\
  --input data/raw/torpeda2012 \\
  --output data/processed/torpeda2012_features.parquet

python build_features_srbh.py \\
  --input data/raw/srbh2020 \\
  --output data/processed/srbh2020_features.parquet
\end{verbatim}

\section*{Entrenamiento y evaluación one-class}

\begin{verbatim}
python train_oneclass.py \\
  --dataset data/processed/csic2010_features.parquet \\
  --target label_binary \\
  --output results/oneclass/csic2010_base

python train_oneclass.py \\
  --dataset data/processed/csic2010_features.parquet \\
  --target label_binary \\
  --drop-suspicious-tokens \\
  --output results/oneclass/csic2010_sin_tokens
\end{verbatim}

\section*{Entrenamiento y evaluación multiclase}

\begin{verbatim}
python train_supervised.py \\
  --task multiclass \\
  --dataset data/processed/torpeda2012_features.parquet \\
  --target label_multiclass \\
  --output results/multiclass/torpeda2012_base

python train_supervised.py \\
  --task multiclass \\
  --dataset data/processed/srbh2020_features.parquet \\
  --target label_multiclass \\
  --drop-suspicious-tokens \\
  --output results/multiclass/srbh2020_sin_tokens
\end{verbatim}

\section*{Entrenamiento y evaluación multietiqueta}

\begin{verbatim}
python train_supervised.py \\
  --task multilabel \\
  --dataset data/processed/srbh2020_features.parquet \\
  --target label_multilabel \\
  --output results/multilabel/srbh2020_base

python train_supervised.py \\
  --task multilabel \\
  --dataset data/processed/srbh2020_features.parquet \\
  --target label_multilabel \\
  --drop-suspicious-tokens \\
  --output results/multilabel/srbh2020_sin_tokens
\end{verbatim}

\section*{Consolidación de resultados}

\begin{verbatim}
python scripts/collect_metrics.py \\
  --input results \\
  --output tables/consolidated_metrics.csv

python scripts/collect_operational_metrics.py \\
  --input results \\
  --output tables/operational_metrics.csv
\end{verbatim}

La reproducción completa requiere conservar las mismas semillas, particiones externas y configuración de hiperparámetros indicadas en los capítulos de metodología y diseño experimental. Cuando se ejecuten variantes ablacionadas, debe verificarse que la única diferencia sea la eliminación del grupo de variables indicado.

\begin{thebibliography}{99}

\bibitem{rfc9110}
Fielding, R., Nottingham, M., Reschke, J. \emph{RFC 9110: HTTP Semantics}. RFC Editor, 2022. \url{https://www.rfc-editor.org/rfc/rfc9110.html}.

\bibitem{rfc9112}
Fielding, R., Nottingham, M., Reschke, J. \emph{RFC 9112: HTTP/1.1}. RFC Editor, 2022. \url{https://www.rfc-editor.org/rfc/rfc9112.html}.



\bibitem{owasp_waf}
OWASP Community. \emph{Web Application Firewall}. \url{https://owasp.org/www-community/Web_Application_Firewall}.

\bibitem{modsecurity_site}
ModSecurity Project. \emph{Open Source Web Application Firewall}. \url{https://modsecurity.org}.

\bibitem{capec_home}
MITRE. \emph{CAPEC (Common Attack Pattern Enumeration and Classification)}. \url{https://capec.mitre.org}.

\bibitem{capec_about}
MITRE. \emph{About CAPEC}. \url{https://capec.mitre.org/about/index.html}.

\bibitem{capec_schema_docs}
MITRE. \emph{CAPEC Schema (XML) documentation}. \url{https://capec.mitre.org/documents/schema/index.html}.

\bibitem{capec_schema_desc}
MITRE. \emph{CAPEC Schema Description (v2.5)}. \url{https://capec.mitre.org/documents/documentation/CAPEC_Schema_Description_v2.5.pdf}.

\bibitem{owasp_risk_rating}
OWASP Community. \emph{OWASP Risk Rating Methodology}. \url{https://owasp.org/www-community/OWASP_Risk_Rating_Methodology}.

\bibitem{srbh_dataverse}
Sureda Riera, T. et al. \emph{SR-BH 2020 multi-label dataset}. Harvard Dataverse.
DOI: \url{https://doi.org/10.7910/DVN/OGOIXX}.

\bibitem{srbh_paper}
Sureda Riera, T. et al. \emph{A new multi-label dataset for Web attacks CAPEC classification using machine learning techniques}.
Expert Systems with Applications, 2022.

\bibitem{kruegel_vigna_2003}
Kruegel, C., Vigna, G. \emph{Anomaly Detection of Web-based Attacks}. ACM CCS, 2003.
\url{https://sites.cs.ucsb.edu/~vigna/publications/2003_kruegel_vigna_ccs03.pdf}.

\bibitem{ingham_inoue_2007}
Ingham, K. L., Inoue, H.
\emph{Comparing Anomaly Detection Techniques for HTTP}.
Recent Advances in Intrusion Detection (RAID 2007), Lecture Notes in Computer Science 4637, 42--62, 2007. DOI: \url{https://doi.org/10.1007/978-3-540-74320-0_3}.

\bibitem{wang_stolfo_2004}
Wang, K., Stolfo, S. J.
\emph{Anomalous Payload-Based Network Intrusion Detection}.
Recent Advances in Intrusion Detection (RAID 2004), Lecture Notes in Computer Science 3224, 203--222, 2004. DOI: \url{https://doi.org/10.1007/978-3-540-30143-1_11}.

\bibitem{rifkin_ova}
Rifkin, R., Klautau, A. \emph{In Defense of One-Vs-All Classification}. JMLR, 2004.
\url{https://www.jmlr.org/papers/volume5/rifkin04a/rifkin04a.pdf}.

\bibitem{tsoumakas_multilabel_overview}
Tsoumakas, G., Katakis, I. \emph{Multi-Label Classification: An Overview}. International Journal of Data Warehousing and Mining, 3(3), 1--13, 2007.

\bibitem{read2011}
Read, J., Pfahringer, B., Holmes, G., Frank, E.
\emph{Classifier Chains for Multi-label Classification}.
Machine Learning, 85, 333--359, 2011. DOI: \url{https://doi.org/10.1007/s10994-011-5256-5}.

\bibitem{sokolova_lapalme_2009}
Sokolova, M., Lapalme, G.
\emph{A Systematic Analysis of Performance Measures for Classification Tasks}.
Information Processing \& Management, 45(4), 427--437, 2009. DOI: \url{https://doi.org/10.1016/j.ipm.2009.03.002}.

\bibitem{powers2011}
Powers, D. M. W.
\emph{Evaluation: From Precision, Recall and F-Measure to ROC, Informedness, Markedness \& Correlation}.
Journal of Machine Learning Technologies, 2(1), 37--63, 2011.

\bibitem{islr_logreg}
James, G., Witten, D., Hastie, T., Tibshirani, R. \emph{An Introduction to Statistical Learning}. Cap. 4 (Logistic Regression).

\bibitem{sklearn_logreg}
scikit-learn developers. \emph{LogisticRegression documentation}. \url{https://scikit-learn.org/stable/modules/generated/sklearn.linear_model.LogisticRegression.html}.

\bibitem{sklearn_gridsearch}
scikit-learn developers. \emph{GridSearchCV documentation}. \url{https://scikit-learn.org/stable/modules/generated/sklearn.model_selection.GridSearchCV.html}.

\bibitem{sklearn_randomsearch}
scikit-learn developers. \emph{RandomizedSearchCV documentation}. \url{https://scikit-learn.org/stable/modules/generated/sklearn.model_selection.RandomizedSearchCV.html}.

\bibitem{sklearn_perm_importance}
scikit-learn developers. \emph{Permutation importance documentation}. \url{https://scikit-learn.org/stable/modules/permutation_importance.html}.

\bibitem{sklearn_sgd_ocsvm}
scikit-learn developers. \emph{SGDOneClassSVM documentation}. \url{https://scikit-learn.org/stable/modules/generated/sklearn.linear_model.SGDOneClassSVM.html}.

\bibitem{sklearn_kernel_approx}
scikit-learn developers. \emph{Kernel Approximation documentation (Nystroem / RBFSampler)}. \url{https://scikit-learn.org/stable/modules/kernel_approximation.html}.

\bibitem{sklearn_halving}
scikit-learn developers. \emph{Successive Halving Iterations documentation}. \url{https://scikit-learn.org/stable/modules/grid_search.html#searching-for-optimal-parameters-with-successive-halving}.

\bibitem{csic2010_impact}
IMPACT Cyber Trust. \emph{HTTP DATASET CSIC 2010}. \url{https://www.impactcybertrust.org/dataset_view?idDataset=940}.

\bibitem{torpeda_2012}
Torrano-Gim{\'e}nez, C., P{\'e}rez, A., {\'A}lvarez, G.
\emph{TORPEDA: Una especificación abierta de conjuntos de datos de ataques a aplicaciones web}. RECSI 2012.
\url{https://recsi2012.mondragon.edu/es/programa/recsi2012_submission_73.pdf}.

\bibitem{csic2010_doc}
Information Security Institute (CSIC).
\emph{HTTP DATASET CSIC 2010 (dataset description)}.
\url{https://petescully.co.uk/wp-content/uploads/2018/04/http_dataset_csic_2010.pdf}.

\bibitem{ocswaf_zenodo}
Epp, N., Funk, R., Cappo, C.
\emph{Anomaly-based Web Application Firewall using HTTP-specific features and One-Class SVM} (artefacto/código, Zenodo).
\url{https://zenodo.org/records/1336812}, 2018.

\bibitem{riverol_2024}
Riverol, A., Betarte, G., Martínez, R., Pardo, Á.
\emph{Capturing the security expert knowledge in feature selection for web application attack detection}.
LADC 2024. arXiv:2407.18445.

\bibitem{fang_2024}
Fang, M., Wang, Y., Yang, L., Wu, H., Yin, Z., Liu, X., Xie, Z., Kong, Z.
\emph{Reinventing web security: An enhanced cycle-consistent generative adversarial network approach to intrusion detection}.
Electronics, 13(9):1711, 2024.

\bibitem{zhou_2024}
Zhou, Y., Zhu, H., Chai, Y., Jiang, Y., Liu, Y.
\emph{Towards Trustworthy Web Attack Detection: An Uncertainty-Aware Ensemble Deep Kernel Learning Model}.
arXiv:2410.07725, 2024.

\bibitem{fawcett2006}
Fawcett, T.
\emph{An Introduction to ROC Analysis}.
Pattern Recognition Letters, 27(8), 861--874, 2006. DOI: \url{https://doi.org/10.1016/j.patrec.2005.10.010}.

\bibitem{matthews1975}
Matthews, B. W.
\emph{Comparison of the Predicted and Observed Secondary Structure of T4 Phage Lysozyme}.
Biochimica et Biophysica Acta, 405(2), 442--451, 1975. DOI: \url{https://doi.org/10.1016/0005-2795(75)90109-9}.

\bibitem{saito_pr_auc}
Saito, T., Rehmsmeier, M.
\emph{The Precision-Recall Plot Is More Informative than the ROC Plot When Evaluating Binary Classifiers on Imbalanced Datasets}.
PLOS ONE 10(3):e0118432, 2015. DOI: \url{https://doi.org/10.1371/journal.pone.0118432}.

\bibitem{chicco_mcc}
Chicco, D., Jurman, G.
\emph{The advantages of the Matthews correlation coefficient (MCC) over F1 score and accuracy in binary classification evaluation}.
BMC Genomics 21, 6 (2020). DOI: \url{https://doi.org/10.1186/s12864-019-6413-7}.


\bibitem{scholkopf2001}
Schölkopf, B., Platt, J. C., Shawe-Taylor, J., Smola, A. J., Williamson, R. C.
\emph{Estimating the Support of a High-Dimensional Distribution}.
Neural Computation, 13(7), 2001.

\bibitem{williams2001}
Williams, C. K. I., Seeger, M.
\emph{Using the Nyström Method to Speed Up Kernel Machines}.
Advances in Neural Information Processing Systems 13, 2001.

\bibitem{bergstra2012}
Bergstra, J., Bengio, Y.
\emph{Random Search for Hyper-Parameter Optimization}.
Journal of Machine Learning Research, 13, 2012.

\bibitem{brodersen2010}
Brodersen, K. H., Ong, C. S., Stephan, K. E., Buhmann, J. M.
\emph{The Balanced Accuracy and Its Posterior Distribution}.
Proceedings of the 20th International Conference on Pattern Recognition, 2010.

\bibitem{gorodkin2004}
J. Gorodkin. ``Comparing two K-category assignments by a K-category correlation coefficient''. \textit{Computational Biology and Chemistry}, 28(5--6):367--374, 2004.

\end{thebibliography}

\end{document}
